% Leçon 19 — Vocabulaire et compréhension de texte : vie culturelle au Cameroun et en Chine — Chinois 1ère A — Cours signé OnBuch+
% Programme officiel MINESEC. Compiler avec : tectonic ZHA19-vocabulaire-vie-culturelle-au-cameroun-et-en-chine.tex
\newcommand{\DOCMATIERE}{Chinois}
\newcommand{\DOCNIVEAU}{1ère A}
\newcommand{\DOCMODULE}{Module 3 : Citoyenneté et ouverture au monde}
\newcommand{\DOCLECON}{ZHA19}
\newcommand{\DOCTITRE}{Leçon 19 — Vocabulaire et compréhension de texte : vie culturelle au Cameroun et en Chine}
% OnBuch+ — préambule « cours signés » ENRICHI (compilé avec tectonic / XeLaTeX)
% Système visuel « Pop » d'OnBuch+ : fond crème (jamais blanc), bordures encre
% épaisses, ombres « dures » décalées, titres Archivo Black en capitales.
% Version MATHÉMATIQUES : graphes (pgfplots/tikz), boîtes pédagogiques
% (définition, propriété, méthode, attention, le savais-tu, exercice résolu,
% exercice, TP, bilan), page de garde et signature OnBuch+ ancrée en bas.
%
% Macros attendues dans main.tex AVANT \input{preamble} :
%   \DOCMATIERE  \DOCNIVEAU  \DOCMODULE  \DOCLECON (numéro)  \DOCTITRE
\documentclass[11pt]{article}

\usepackage{fontspec}
\usepackage[french]{babel} % français = langue principale ; le chinois via \zh{..} / textebox (xeCJK)
\usepackage{xeCJK}
\usepackage{pifont} % \ding{51} \ding{55} utilisés par les modèles
\usepackage[a4paper,margin=2cm,top=3.4cm,bottom=2.8cm,headheight=1.4cm,footskip=1.3cm]{geometry}
\usepackage{amsmath,amssymb,amsfonts}
\usepackage{mathtools} % \xmapsto, \coloneqq... souvent utilisés par les modèles
\usepackage{cancel} % simplifications barrées (bilans de forces)
% \abs{z} (module, valeur absolue) et \norm{u} (norme), utilisés par les modèles
\providecommand{\abs}{}\let\abs\relax\DeclarePairedDelimiter{\abs}{\lvert}{\rvert}
\providecommand{\norm}{}\let\norm\relax\DeclarePairedDelimiter{\norm}{\lVert}{\rVert}
\usepackage[table]{xcolor} % [table] : fournit aussi \rowcolors (alternance de lignes)
\usepackage{pagecolor}
\usepackage{tikz}
\usetikzlibrary{arrows.meta,positioning,calc,shapes.geometric,shapes.misc,shapes.arrows,decorations.pathmorphing,decorations.pathreplacing,decorations.markings,patterns,shadows,mindmap,trees,fit,backgrounds,matrix,intersections,angles,quotes}
\usepackage{pgfplots}
\usepackage[version=4]{mhchem} % équations nucléaires \ce{^{235}_{92}U}
\usepackage[european,siunitx]{circuitikz} % schémas électriques
\pgfplotsset{compat=1.18}
\usepgfplotslibrary{fillbetween,groupplots} % aires entre courbes (calcul intégral)
\usepackage{enumitem}
\usepackage{fancyhdr}
% [most] charge toutes les bibliothèques tcolorbox (listings, minted, raster,
% documentation...) et gonfle inutilement la mémoire TeX sur un document long
% (~300 boîtes) : on ne charge que ce qui est réellement utilisé.
\usepackage[skins,breakable]{tcolorbox}
\usepackage{graphicx}
\usepackage{colortbl}
\usepackage{booktabs}
\usepackage{tabularx}
\usepackage{diagbox} % case d'en-tête diagonale des tableaux à double entrée
% Colonnes extensibles centrée / gauche / droite (C, L, R), souvent utilisées par les modèles
\newcolumntype{C}{>{\centering\arraybackslash}X}
\newcolumntype{L}{>{\raggedright\arraybackslash}X}
\newcolumntype{R}{>{\raggedleft\arraybackslash}X}
\usepackage{array}
\usepackage{multicol}
\usepackage{multirow}
\usepackage{siunitx}
% parse-numbers=false : accepte aussi \SI{2,5 \times 10^{-3}}{...}
\sisetup{locale=FR,output-decimal-marker={,},group-minimum-digits=4,parse-numbers=false}
\DeclareSIUnit\ohm{\ensuremath{\symup{\Omega}}} % U+2126 absent de la police
\DeclareSIUnit\division{div} % réglages d'oscilloscope (ms/div, V/div)
\DeclareSIUnit\tr{tr} % tours (vitesse de rotation, ex. \SI{1500}{\tr\per\minute})
\DeclareSIUnit\molar{M} % concentration molaire (ex. \SI{3}{\molar}, \SI{1}{\micro\molar})
\DeclareSIUnit\an{an} % année (le modèle confond parfois avec \year, un compteur TeX) — ex. \SI{5}{\kilo\meter\per\an}
\DeclareSIUnit\FCFA{FCFA} % franc CFA (ex. \SI{500}{\FCFA\per\kilo\gram})
\DeclareSIUnit\degreeBrix{\textdegree Brix} % teneur en sucre d'un sirop/jus (réfractométrie)
\DeclareSIUnit\month{mois} % le modèle utilise parfois \month, un compteur TeX, comme unité de durée
\usepackage{qrcode}
% \degree : commande fréquemment utilisée par les modèles pour un angle ou une
% température en dehors de \SI{}{} (non fournie par les paquets chargés).
\providecommand{\degree}{^{\circ}}
% \qed (fin de démonstration), utilisé par les modèles sans amsthm
\providecommand{\qed}{\hfill$\square$}
% \barre : double barre des valeurs interdites dans un tableau de variations (tkz-tab)
\providecommand{\barre}{\ensuremath{\|}}
% environnement proof (amsthm non chargé) utilisé par les modèles
\ifdefined\proof\else\newenvironment{proof}[1][Démonstration]{\par\noindent\textit{#1.}\ }{\qed\par}\fi
\usepackage{lastpage}
\usepackage{hyperref}
\hypersetup{hidelinks}

% ============================================================
% Palette « Pop » OnBuch+ — mêmes hex que la classe P (plus_ui.dart)
% ============================================================
\definecolor{poporange}{HTML}{F59321}
\definecolor{popdark}{HTML}{E07A0C}
\definecolor{popcream}{HTML}{FFF6E8}
\definecolor{popink}{HTML}{1C1714}
\definecolor{poppurple}{HTML}{7A5AE0}
\definecolor{popgreen}{HTML}{2E9E6A}
\definecolor{poppink}{HTML}{E0457B}
\definecolor{popblue}{HTML}{2F7DE1}
\definecolor{popgold}{HTML}{C98A12}
\definecolor{popmuted}{HTML}{8A7F76}
\definecolor{popline}{HTML}{EADFD2}
\definecolor{popgray}{HTML}{8A7F76}
\colorlet{popgrey}{popgray}
% Alias tolérants (le modèle nomme parfois les couleurs différemment)
\colorlet{popwhite}{white}
\colorlet{popblack}{popink}
\colorlet{popred}{poppink}
\colorlet{popyellow}{popgold}
% Teintes claires (fonds de tableaux/encadrés) — le modèle nomme parfois
% les couleurs avec un suffixe L (ex. popblueL) sans les avoir définies.
\colorlet{poporangeL}{poporange!20}
\colorlet{popdarkL}{popdark!20}
\colorlet{poppurpleL}{poppurple!20}
\colorlet{popgreenL}{popgreen!20}
\colorlet{poppinkL}{poppink!20}
\colorlet{popblueL}{popblue!20}
\colorlet{popgoldL}{popgold!20}
\colorlet{popredL}{popred!20}
\colorlet{popgrayL}{popgray!20}
% Teintes claires pour fonds de tableaux / zones
\colorlet{popblueL}{popblue!10!white}
\colorlet{poporangeL}{poporange!14!white}
\colorlet{popgreenL}{popgreen!12!white}
\colorlet{poppurpleL}{poppurple!12!white}
\colorlet{poppinkL}{poppink!12!white}
\colorlet{popgoldL}{popgold!14!white}

\pagecolor{popcream}

% ============================================================
% Typographie
% ============================================================
\defaultfontfeatures{Ligatures=TeX}
\setmainfont{PlusJakartaSans-Medium.ttf}[Path=../../../fonts/,BoldFont=PlusJakartaSans-Bold.ttf]
% Symboles, API et emojis absents de la police principale : repli sur DejaVu Sans (caractères actifs)
\newfontface\symfont{DejaVuSans.ttf}[Path=../../../fonts/]
\makeatletter
\newcommand\symactive[1]{\catcode#1=13 \begingroup\lccode`\~=#1 \lowercase{\endgroup\def~}{{\symfont\char#1}}}
\newcommand\symblank[1]{\catcode#1=13 \begingroup\lccode`\~=#1 \lowercase{\endgroup\def~}{}}
\newcommand\symhyph[1]{\catcode#1=13 \begingroup\lccode`\~=#1 \lowercase{\endgroup\def~}{\mbox{-}}}
\newcount\symc
\newcommand\symrange[2]{\symc=#1 \@whilenum\symc<#2 \do{\expandafter\symactive\expandafter{\the\symc}\advance\symc by 1 }}
\newcommand\symblankrange[2]{\symc=#1 \@whilenum\symc<#2 \do{\expandafter\symblank\expandafter{\the\symc}\advance\symc by 1 }}
\makeatother
\symrange{"0250}{"02FF}\symactive{"2713}\symactive{"2714}\symactive{"2717}\symactive{"2718}\symactive{"2610}\symactive{"2611}\symactive{"2612}
\symrange{"2600}{"27BF}
\catcode"2705=13 \begingroup\lccode`\~="2705 \lowercase{\endgroup\def~}{{\symfont\char"2713}}\symblank{"26BD}\symblank{"26A1}
\symrange{"2460}{"257F}\symactive{"223C}\symactive{"2248}\symactive{"2260}
\symactive{"01F8}\symactive{"01F9}\symactive{"1E3E}\symactive{"1E3F}
\symhyph{"2011}\symhyph{"2E1A}\symblankrange{"1F300}{"1FAFF}\symblank{"FE0F}
% Caractères chinois : WenQuanYi Zen Hei (sous-ensemble GB2312 + pinyin), gras/italique simulés
\setCJKmainfont{WenQuanYiZenHei-subset.ttf}[Path=../../../fonts/,AutoFakeBold=3,AutoFakeSlant=0.2]
\xeCJKsetup{CJKspace=false}
% Pinyin : voyelles à caron (ǎ ǐ ǒ ǔ ǖ ǘ ǚ ǜ) absentes de la police principale -> caractères actifs
\newfontface\pyfont{WenQuanYiZenHei-subset.ttf}[Path=../../../fonts/]
\newcommand\pyactive[1]{\catcode#1=13 \begingroup\lccode`\~=#1 \lowercase{\endgroup\def~}{{\pyfont\char#1}}}
\pyactive{"01CD}\pyactive{"01CE}\pyactive{"01CF}\pyactive{"01D0}\pyactive{"01D1}\pyactive{"01D2}\pyactive{"01D3}\pyactive{"01D4}\pyactive{"01D5}\pyactive{"01D6}\pyactive{"01D7}\pyactive{"01D8}\pyactive{"01D9}\pyactive{"01DA}\pyactive{"01DB}\pyactive{"01DC}
% Maths en sans-serif (Fira Math) avec chiffres et lettres droites de Plus
% Jakarta, pour que les formules se fondent dans le texte.
\usepackage[math-style=ISO,bold-style=ISO]{unicode-math}
\setmathfont{FiraMath-Regular.otf}[Path=../../../fonts/]
\setmathfont{PlusJakartaSans-Medium.ttf}[Path=../../../fonts/,range={up/num,up/{Latin,latin},\mathup}]
\usepackage{icomma} % virgule décimale sans espace en mode math
% Caractères mathématiques Unicode absents des polices de texte (Plus Jakarta,
% Archivo Black) : rendus par leur équivalent mathématique, en texte comme en
% formule — sinon ils s'affichent en carré vide (ex. « Arithmétique dans ℤ »).
\usepackage{newunicodechar}
\newunicodechar{ℝ}{\ensuremath{\mathbb{R}}}
\newunicodechar{ℤ}{\ensuremath{\mathbb{Z}}}
\newunicodechar{ℂ}{\ensuremath{\mathbb{C}}}
\newunicodechar{ℕ}{\ensuremath{\mathbb{N}}}
\newunicodechar{ℚ}{\ensuremath{\mathbb{Q}}}
\newunicodechar{𝔻}{\ensuremath{\mathbb{D}}}
\newunicodechar{α}{\ensuremath{\alpha}}
\newunicodechar{β}{\ensuremath{\beta}}
\newunicodechar{γ}{\ensuremath{\gamma}}
\newunicodechar{δ}{\ensuremath{\delta}}
\newunicodechar{ε}{\ensuremath{\varepsilon}}
\newunicodechar{θ}{\ensuremath{\theta}}
\newunicodechar{λ}{\ensuremath{\lambda}}
\newunicodechar{μ}{\ensuremath{\mu}}
\newunicodechar{σ}{\ensuremath{\sigma}}
\newunicodechar{τ}{\ensuremath{\tau}}
\newunicodechar{φ}{\ensuremath{\varphi}}
\newunicodechar{ψ}{\ensuremath{\psi}}
\newunicodechar{ω}{\ensuremath{\omega}}
\newunicodechar{Σ}{\ensuremath{\Sigma}}
% majuscules produites par \MakeUppercase dans les titres (α-aminés -> Α)
\newunicodechar{Α}{\ensuremath{\alpha}}
\newunicodechar{Β}{\ensuremath{\beta}}
\newunicodechar{Γ}{\ensuremath{\Gamma}}
\newunicodechar{Δ}{\ensuremath{\Delta}}
\newunicodechar{Ε}{\ensuremath{\varepsilon}}
\newunicodechar{Θ}{\ensuremath{\Theta}}
\newunicodechar{Λ}{\ensuremath{\Lambda}}
\newunicodechar{Μ}{\ensuremath{\mu}}
\newunicodechar{Τ}{\ensuremath{\tau}}
\newunicodechar{Φ}{\ensuremath{\Phi}}
\newunicodechar{Ψ}{\ensuremath{\Psi}}
\newunicodechar{Ω}{\ensuremath{\Omega}}
\newunicodechar{↦}{\ensuremath{\mapsto}}
\newunicodechar{∈}{\ensuremath{\in}}
\newunicodechar{∉}{\ensuremath{\notin}}
\newunicodechar{∧}{\ensuremath{\wedge}}
\newunicodechar{⇔}{\ensuremath{\Leftrightarrow}}
\newunicodechar{⟺}{\ensuremath{\Longleftrightarrow}}
\newunicodechar{⇒}{\ensuremath{\Rightarrow}}
\newunicodechar{⟹}{\ensuremath{\Longrightarrow}}
\newunicodechar{∘}{\ensuremath{\circ}}
\newunicodechar{∪}{\ensuremath{\cup}}
\newunicodechar{∩}{\ensuremath{\cap}}
\newunicodechar{≡}{\ensuremath{\equiv}}
\newunicodechar{⊂}{\ensuremath{\subset}}
\newunicodechar{⊥}{\ensuremath{\perp}}
\newunicodechar{∃}{\ensuremath{\exists}}
\newunicodechar{∀}{\ensuremath{\forall}}
\newunicodechar{∛}{\ensuremath{\sqrt[3]{\,}}}
\newunicodechar{∜}{\ensuremath{\sqrt[4]{\,}}}
\newunicodechar{⃗}{\ensuremath{{}^{\to}}}
\newunicodechar{ⁿ}{\ifmmode^{n}\else\textsuperscript{\ensuremath{n}}\fi}
\newunicodechar{ˣ}{\ifmmode^{x}\else\textsuperscript{\ensuremath{x}}\fi}
\newunicodechar{ᵇ}{\ifmmode^{b}\else\textsuperscript{\ensuremath{b}}\fi}
\newunicodechar{ʳ}{\ifmmode^{r}\else\textsuperscript{\ensuremath{r}}\fi}
\newunicodechar{ᵘ}{\ifmmode^{u}\else\textsuperscript{\ensuremath{u}}\fi}
\newunicodechar{ʸ}{\ifmmode^{y}\else\textsuperscript{\ensuremath{y}}\fi}
\newunicodechar{ᵃ}{\ifmmode^{a}\else\textsuperscript{\ensuremath{a}}\fi}
\newunicodechar{ᵉ}{\ifmmode^{e}\else\textsuperscript{\ensuremath{e}}\fi}
\newunicodechar{ᵏ}{\ifmmode^{k}\else\textsuperscript{\ensuremath{k}}\fi}
\newunicodechar{ᵖ}{\ifmmode^{p}\else\textsuperscript{\ensuremath{p}}\fi}
\newunicodechar{ᴺ}{\ifmmode^{N}\else\textsuperscript{\ensuremath{N}}\fi}
\newunicodechar{ᶜ}{\ifmmode^{c}\else\textsuperscript{\ensuremath{c}}\fi}
\newunicodechar{ʰ}{\ifmmode^{h}\else\textsuperscript{\ensuremath{h}}\fi}
\newunicodechar{ᵠ}{\ifmmode^{q}\else\textsuperscript{\ensuremath{q}}\fi}
\newunicodechar{ₙ}{\ifmmode_{n}\else\textsubscript{\ensuremath{n}}\fi}
\newunicodechar{ₐ}{\ifmmode_{a}\else\textsubscript{\ensuremath{a}}\fi}
\newunicodechar{ᵢ}{\ifmmode_{i}\else\textsubscript{\ensuremath{i}}\fi}
\newunicodechar{ᵣ}{\ifmmode_{r}\else\textsubscript{\ensuremath{r}}\fi}
\newunicodechar{ₖ}{\ifmmode_{k}\else\textsubscript{\ensuremath{k}}\fi}
\newunicodechar{ᵦ}{\ifmmode_{\beta}\else\textsubscript{\ensuremath{\beta}}\fi}
\newunicodechar{ₓ}{\ifmmode_{x}\else\textsubscript{\ensuremath{x}}\fi}
\newfontfamily\popdisplay{ArchivoBlack-Regular.ttf}[Path=../../../fonts/]
\newfontfamily\poplogo{SpaceGrotesk-Bold.ttf}[Path=../../../fonts/]
\newfontfamily\popsemi{PlusJakartaSans-SemiBold.ttf}[Path=../../../fonts/]
\newfontfamily\popxbold{PlusJakartaSans-ExtraBold.ttf}[Path=../../../fonts/]
\newfontfamily\popxboldls{PlusJakartaSans-ExtraBold.ttf}[Path=../../../fonts/,LetterSpace=3.0]

\setlist[enumerate]{leftmargin=1.7em,itemsep=5pt,topsep=4pt}
\setlist[itemize]{leftmargin=1.7em,itemsep=4pt,topsep=4pt,label={\color{poporange}\textbullet}}
\setlength{\parindent}{0pt}
\setlength{\parskip}{5pt}
\renewcommand{\arraystretch}{1.25}

% pgfplots : style Pop par défaut
\pgfplotsset{
  every axis/.append style={axis line style={popink,line width=1pt},tick style={popink},
    grid style={popline,line width=0.5pt},label style={font=\small},tick label style={font=\footnotesize}},
  popcurve/.style={poporange,line width=1.8pt,no markers,smooth},
}
% Même style disponible dans un \draw TikZ simple (hors pgfplots)
\tikzset{popcurve/.style={poporange,line width=1.8pt}}
\tikzset{popgrid/.style={popline,very thin}}
\colorlet{popcurve}{poporange} % le nom du style est parfois utilisé comme couleur

% ============================================================
% Pill / squiggle
% ============================================================
\newcommand{\poppill}[3]{%
  \tikz[baseline=-0.55ex]\node[draw=popink,line width=1.2pt,rounded corners=2.6pt,%
    fill=#2,inner xsep=6.5pt,inner ysep=3pt]{%
    \popxboldls\fontsize{7}{7}\selectfont\color{#3}\MakeUppercase{#1}};%
}
\newcommand{\popsquiggle}[1]{%
  \tikz[baseline=-0.4ex]{
    \draw[#1,line width=2.6pt,line cap=round]
      plot [smooth] coordinates {(0,0.09) (0.34,0.01) (0.68,0.21) (1.02,0.01) (1.36,0.10)};
  }%
}
% Mot-clé surligné (marqueur orange sous le texte)
\newcommand{\cle}[1]{{\popxbold\color{popdark}#1}}

% ============================================================
% En-tête / pied
% ============================================================
\pagestyle{fancy}
\fancyhf{}
\renewcommand{\headrulewidth}{0pt}
\renewcommand{\footrulewidth}{0pt}
\lhead{\raisebox{-0.15ex}{\poplogo\fontsize{13}{13}\selectfont\color{popink}OnBuch\color{poporange}+}}
\rhead{\poppill{\DOCMATIERE\ \textbullet\ \DOCNIVEAU\ \textbullet\ Leçon \DOCLECON}{white}{popink}}
\lfoot{\popxbold\fontsize{7.6}{7.6}\selectfont\color{popmuted}\MakeUppercase{Cours signé OnBuch+}\ \textbullet\ {\popsemi\DOCTITRE}}
\rfoot{\tikz[baseline=-0.6ex]\node[rounded corners=8.5pt,fill=popink,minimum height=17pt,minimum width=17pt,inner xsep=5pt,inner sep=0]{\popxbold\fontsize{8}{8}\selectfont\color{popcream}\thepage\,/\,\pageref*{LastPage}};}

% ============================================================
% Titres
% ============================================================
\newcommand{\chaptitre}[2]{%
  \begin{tcolorbox}[enhanced,colback=poporange,colframe=popink,boxrule=2.6pt,arc=11pt,
      drop shadow=popink,left=15pt,right=15pt,top=12pt,bottom=11pt,before skip=3pt,after skip=18pt]
    \poppill{#2}{popink}{popcream}\par\vspace{9pt}
    {\raggedright\hyphenpenalty=10000\exhyphenpenalty=10000\popdisplay\fontsize{22}{24}\selectfont\color{popcream}\MakeUppercase{#1}\par}
  \end{tcolorbox}
}
% Section numérotée : pastille orange avec le numéro + titre Archivo
\newcounter{popsec}
\newcommand{\coursec}[1]{%
  \stepcounter{popsec}\setcounter{popsub}{0}%
  \par\vspace{16pt}\noindent
  \begin{minipage}{\linewidth}
  \tikz[baseline=-0.7ex]\node[draw=popink,line width=1.6pt,fill=poporange,rounded corners=4pt,minimum size=22pt,inner sep=0]{\popdisplay\fontsize{12}{12}\selectfont\color{popcream}\thepopsec};%
  \hspace{8pt}{\popdisplay\fontsize{15}{17}\selectfont\color{popink}\MakeUppercase{#1}}\par
  \vspace{4pt}\noindent{\color{poporange}\rule{36pt}{3.6pt}}
  \end{minipage}\par\nopagebreak\vspace{8pt}
}
\newcounter{popsub}
\newcommand{\courssub}[1]{%
  \stepcounter{popsub}%
  \par\vspace{9pt}\noindent{\popxbold\fontsize{11.5}{13}\selectfont\color{popink}{\color{poporange}\thepopsec.\thepopsub}\hspace{6pt}#1}\par\nopagebreak\vspace{4pt}
}

% ============================================================
% Boîtes pédagogiques — même recette : fond blanc, bordure épaisse colorée,
% ombre dure encre, badge plein. Titre optionnel [#1].
% ============================================================
\tcbset{popbase/.style={breakable,enhanced,colback=white,boxrule=2.3pt,arc=9pt,
  shadow={3pt}{-3pt}{0pt}{popink},left=14pt,right=14pt,top=11pt,bottom=11pt,before skip=11pt,after skip=15pt,
  fontupper=\popsemi\fontsize{10.6}{14.5}\selectfont\color{popink},
  fontlower=\popsemi\fontsize{10.6}{14.5}\selectfont\color{popink}}}
% Coche dessinée (le glyphe ✓ n'existe pas dans les polices du document)
\newcommand{\popcheck}[1]{\tikz[baseline=-0.5ex]\draw[#1,line width=1.6pt,line cap=round,line join=round](-0.13,0.0)--(-0.04,-0.09)--(0.13,0.1);}
\providecommand{\coche}{\popcheck{popgreen}}
\providecommand{\resultat}[1]{\par\smallskip\noindent\fcolorbox{popgreen}{popgreenL}{\parbox{\dimexpr\linewidth-2\fboxsep-2\fboxrule\relax}{\textbf{Résultat.} #1}}\par\smallskip}
\newcommand{\popboxhead}[3]{\poppill{#1}{#2}{white}\ifx\relax#3\relax\else\hspace{6pt}{\popxbold\fontsize{10.6}{12}\selectfont\color{popink}#3}\fi\par\vspace{7pt}}

\newtcolorbox{aretenir}[1][]{popbase,colframe=popgreen,before upper={\popboxhead{À retenir}{popgreen}{#1}}}
% Texte en chinois (document de lecture, dialogue, extrait) : cadre
% d'étude. Un titre optionnel (source, auteur) s'affiche dans l'en-tête.
\newtcolorbox{textebox}[1][]{popbase,colframe=popblue,before upper={\popboxhead{文本 · Texte}{popblue}{#1}},after upper={}}
% Marques ✓ / ✗ (forme correcte / fautive) souvent utilisées par les modèles
\providecommand{\cross}{\textcolor{poppink}{\ensuremath{\times}}}
\providecommand{\xmark}{\cross}\providecommand{\crossmark}{\cross}\providecommand{\cmark}{\ensuremath{\checkmark}}
% Ligne de réponse / justification à compléter (inventée par les modèles)
\providecommand{\justification}[1]{\par\noindent\textit{Justification~:} \dotfill\par}
% \textipa (package tipa non chargé) : la police ne couvre pas l'API -> simple passage
\providecommand{\textipa}[1]{#1}
% Soulignement (ulem non chargé) : \uline, \ul
\providecommand{\uline}[1]{\underline{#1}}\providecommand{\ul}[1]{\underline{#1}}
% \glossaire{\item ...} inventée par les modèles : liste de glossaire
\providecommand{\glossaire}[1]{\begin{itemize}#1\end{itemize}}
% \alert (beamer) inventée par les modèles : texte mis en évidence
\providecommand{\alert}[1]{\textbf{\textcolor{poppink}{#1}}}
% variantes ASCII d'ordinaux écrites en macro
\providecommand{\iere}{\textsuperscript{re}}\providecommand{\ieme}{\textsuperscript{e}}\providecommand{\ere}{\textsuperscript{re}}\providecommand{\eme}{\textsuperscript{e}}\providecommand{\er}{\textsuperscript{er}}\providecommand{\ie}{\textsuperscript{e}}
% Ordinaux français écrits en macro par les modèles : 1\ière, 3\ième
\newcommand{\ière}{\textsuperscript{re}}\newcommand{\ième}{\textsuperscript{e}}
% \exemple (inventée par les modèles) : étiquette « Exemple : »
\providecommand{\exemple}{\textbf{Exemple~:}\ }
% \square utilisable aussi hors mode math (cases à cocher)
\AtBeginDocument{\let\squareMath\square\renewcommand{\square}{\ensuremath{\squareMath}}}
% Mot ou phrase en chinois à l'intérieur d'un paragraphe français
\newcommand{\zh}[1]{\ifmmode\text{#1}\else{#1}\fi}
\let\eng\zh \let\esp\zh \let\all\zh % alias tolérés
% flèches d'intonation inventées par les modèles
\providecommand{\textsearrow}{}\renewcommand{\textsearrow}{\ensuremath{\searrow}}\providecommand{\textnearrow}{}\renewcommand{\textnearrow}{\ensuremath{\nearrow}}
% \fr{..} : mot français (inventé par les modèles)
\providecommand{\fr}[1]{\textit{#1}}
% zhbox : encadré de texte chinois inventé par les modèles
\newenvironment{zhbox}{\begin{quote}}{\end{quote}}
% \courssubsub : titre de niveau 3 inventé par les modèles
\providecommand{\courssubsub}[1]{\par\medskip\noindent\textbf{#1}\par\smallskip}
% \vs : « versus » inventé par les modèles
\providecommand{\vs}{\textit{vs}\ }
% \blank : case à compléter inventée par les modèles
\providecommand{\blank}[1][]{\underline{\hspace{1.6em}}}
\newtcolorbox{exemplebox}[1][]{popbase,colframe=poporange,before upper={\popboxhead{Exemple}{poporange}{#1}}}
\newtcolorbox{definition}[1][]{popbase,colframe=popblue,before upper={\popboxhead{Définition}{popblue}{#1}}}
\newtcolorbox{propriete}[1][]{popbase,colframe=poppurple,before upper={\popboxhead{Propriété}{poppurple}{#1}}}
\newtcolorbox{methode}[1][]{popbase,colframe=popgold,before upper={\popboxhead{Méthode}{popgold}{#1}}}
\newtcolorbox{attention}[1][]{popbase,colframe=poppink,before upper={\popboxhead{Attention}{poppink}{#1}}}
\newtcolorbox{experience}[1][]{popbase,colframe=popblue,colback=popblueL,before upper={\popboxhead{Expérience}{popblue}{#1}}}
% « Le savais-tu ? » : carte violette pleine (contexte, culture, Cameroun)
\newtcolorbox{savaistu}[1][]{popbase,colback=poppurple,colframe=popink,
  fontupper=\popsemi\fontsize{10.4}{14.2}\selectfont\color{white},
  before upper={\poppill{Le savais-tu\,?}{white}{poppurple}\ifx\relax#1\relax\else\hspace{6pt}{\popxbold\color{white}#1}\fi\par\vspace{7pt}}}
% Exercice résolu : énoncé en haut, \tcblower puis la solution sur fond crème
\newcounter{popexo}\newcounter{popexr}
\newtcolorbox{exoresolu}[1][]{popbase,colframe=poporange,colbacklower=popcream,
  segmentation style={popink,line width=1.2pt,dashed},
  before upper={\stepcounter{popexr}\popboxhead{Exercice résolu \thepopexr}{poporange}{#1}},
  before lower={\poppill{Solution}{popink}{popcream}\par\vspace{6pt}}}
% Exercice (non résolu en place) — la correction est dans la partie Corrigés
\newtcolorbox{exercice}[1][]{popbase,colframe=popink,
  before upper={\stepcounter{popexo}\popboxhead{Exercice \thepopexo}{popink}{#1}}}
% Corrigé
\newtcolorbox{corrige}[1][]{popbase,colframe=popgreen,colback=popgreenL,
  before upper={\popboxhead{Corrigé}{popgreen}{#1}}}
% Objectifs / prérequis / bilan
\newtcolorbox{objectifs}{popbase,colframe=popink,colback=poporangeL,
  before upper={\popboxhead{Objectifs de la leçon}{popink}{}}}
\newtcolorbox{prerequis}{popbase,colframe=popink,colback=popblueL,
  before upper={\popboxhead{Prérequis}{popblue}{}}}
\newtcolorbox{bilan}{popbase,colframe=popink,colback=white,boxrule=2.8pt,
  before upper={\popboxhead{Fiche bilan}{poporange}{L'essentiel à connaître pour l'examen}}}
% Figure encadrée avec légende
\newtcolorbox{popfigure}[1][]{enhanced,breakable=false,colback=white,colframe=popline,boxrule=1.4pt,arc=8pt,
  left=10pt,right=10pt,top=10pt,bottom=8pt,before skip=10pt,after skip=12pt,halign=center,
  lower separated=false,halign lower=center,
  fontlower=\popsemi\fontsize{9}{11.5}\selectfont\color{popmuted},
  #1}
\newcommand{\legende}[1]{\par\vspace{6pt}{\popsemi\fontsize{9}{11.5}\selectfont\color{popmuted}#1}}

% ============================================================
% Signature OnBuch+
% ============================================================
\newcommand{\poplogoinline}[1]{{\poplogo\fontsize{#1}{#1}\selectfont\color{popink}OnBuch\color{poporange}+}}
\newcommand{\signatureonbuch}{%
  \begin{tcolorbox}[enhanced,colback=popink,colframe=popink,boxrule=2.2pt,arc=10pt,
      drop shadow=poporange,left=14pt,right=14pt,top=10pt,bottom=10pt,before skip=0pt,after skip=0pt]
    \tikz[baseline=-0.7ex]\node[circle,fill=popgreen,draw=popcream,line width=1.3pt,minimum size=22pt,inner sep=0]{\popcheck{white}};%
    \hspace{9pt}\begin{minipage}[c]{0.62\linewidth}
      {\popxbold\fontsize{12}{13}\selectfont\color{popcream}Signé\ \textbullet\ {\poplogo OnBuch\color{poporange}+}}\par\vspace{2pt}
      {\raggedright\popsemi\fontsize{8}{10.5}\selectfont\color{popline}\MakeUppercase{Équipe pédagogique OnBuch+}\\\MakeUppercase{Conforme au programme officiel MINESEC}\par}
    \end{minipage}\hfill
    {\poplogo\fontsize{22}{22}\selectfont\color{popcream}OnBuch\color{poporange}+}
  \end{tcolorbox}%
}

% ============================================================
% Page de garde
% #1 titre · #2 matière · #3 niveau · #4 module · #5 n° leçon · #6 durée ·
% #7 description / objectifs (texte ou itemize)
% ============================================================
\newcommand{\pagegarde}[7]{%
  \thispagestyle{empty}
  \newgeometry{a4paper,margin=1.7cm,top=1.6cm,bottom=1.4cm}%
  \begin{tikzpicture}[remember picture,overlay]
    \fill[poporange!18] ([xshift=-3.2cm,yshift=-3.6cm]current page.north east) circle (4.6cm);
    \fill[poppurple!14] ([xshift=2.4cm,yshift=7.6cm]current page.south west) circle (3.8cm);
  \end{tikzpicture}%
  {\poplogo\fontsize{30}{30}\selectfont\color{popink}OnBuch\color{poporange}+}\hfill
  \raisebox{0.6ex}{\poppill{Cours signé \textbullet\ Premium}{popink}{popcream}}\par
  \vspace{1pt}\popsquiggle{poppurple}\par
  \vspace{8pt}
  \begin{tcolorbox}[enhanced,colback=poporange,colframe=popink,boxrule=2.8pt,arc=13pt,
      drop shadow=popink,left=18pt,right=18pt,top=12pt,bottom=13pt,after skip=10pt]
    \poppill{#2 \textbullet\ #3}{popink}{popcream}\hspace{5pt}\poppill{#4}{white}{popink}\par\vspace{8pt}
    {\popdisplay\fontsize{14}{14}\selectfont\color{popink}LEÇON #5}\par\vspace{4pt}
    {\raggedright\hyphenpenalty=10000\exhyphenpenalty=10000\popdisplay\fontsize{25}{27}\selectfont\color{popcream}\MakeUppercase{#1}\par}
    \vspace{8pt}
    {\popxbold\fontsize{9.5}{9.5}\selectfont\color{popink}\MakeUppercase{Durée conseillée : #6}}
  \end{tcolorbox}
  \begin{tcolorbox}[enhanced,colback=white,colframe=popink,boxrule=2.2pt,arc=10pt,
      drop shadow=popink,left=16pt,right=16pt,top=10pt,bottom=10pt,after skip=8pt]
    \poppill{Dans cette leçon}{poppurple}{white}\par\vspace{5pt}
    \popsemi\fontsize{9.3}{12.6}\selectfont\color{popink}#7
  \end{tcolorbox}
  \noindent\begin{minipage}[t]{0.487\linewidth}
    \begin{tcolorbox}[enhanced,colback=popgreen,colframe=popink,boxrule=2.2pt,arc=10pt,
        drop shadow=popink,left=11pt,right=11pt,top=10pt,bottom=10pt,height=3.1cm,valign=center]
      \tcbox[colback=white,colframe=popink,boxrule=1.3pt,arc=5pt,boxsep=0pt,nobeforeafter,
        left=3pt,right=3pt,top=3pt,bottom=3pt,valign=center]{\qrcode[height=1.9cm]{https://play.google.com/store/apps/details?id=com.luvvix.onbuch&hl=fr}}%
      \hspace{8pt}\begin{minipage}[c]{0.5\linewidth}
        {\popdisplay\fontsize{11}{12.5}\selectfont\color{white}\MakeUppercase{Télécharge OnBuch}}\par\vspace{4pt}
        {\raggedright\popsemi\fontsize{8.4}{11}\selectfont\color{white}Gratuit \textbullet\ Google Play\\Tuteur IA, annales, concours, résultats\par}
      \end{minipage}
    \end{tcolorbox}
  \end{minipage}\hfill
  \begin{minipage}[t]{0.487\linewidth}
    \begin{tcolorbox}[enhanced,colback=poppurple,colframe=popink,boxrule=2.2pt,arc=10pt,
        drop shadow=popink,left=11pt,right=11pt,top=10pt,bottom=10pt,height=3.1cm,valign=center]
      \tikz[baseline=-0.7ex]\node[circle,fill=white,draw=popink,line width=1.3pt,minimum size=1.6cm,inner sep=0]{\popdisplay\fontsize{24}{24}\selectfont\color{poppurple}+};%
      \hspace{8pt}\begin{minipage}[c]{0.6\linewidth}
        {\popdisplay\fontsize{11}{12.5}\selectfont\color{white}\MakeUppercase{Passe à OnBuch+}}\par\vspace{4pt}
        {\raggedright\popsemi\fontsize{8.4}{11}\selectfont\color{white}Tous les cours signés, Tuteur IA illimité, fiches PDF — onglet OnBuch+ de l'app\par}
      \end{minipage}
    \end{tcolorbox}
  \end{minipage}
  \vspace{8pt}
  \popcontenu
  \vfill
  \signatureonbuch
  \restoregeometry
  \newpage
}

% Rangée « Ce cours contient » (page de garde)
\newcommand{\poptile}[3]{% #1 couleur #2 gros texte #3 légende
  \begin{tcolorbox}[enhanced,colback=#1,colframe=popink,boxrule=2pt,arc=9pt,shadow={2.5pt}{-2.5pt}{0pt}{popink},
      left=6pt,right=6pt,top=9pt,bottom=9pt,halign=center,valign=center,height=1.75cm,width=\linewidth,nobeforeafter]
    {\popdisplay\fontsize{12.5}{13}\selectfont\color{white}\MakeUppercase{#2}}\par\vspace{3pt}
    {\centering\popsemi\fontsize{7.8}{9.5}\selectfont\color{white}#3\par}
  \end{tcolorbox}}
\newcommand{\popcontenu}{%
  \par\noindent{\popxboldls\fontsize{7.5}{7.5}\selectfont\color{popmuted}CE COURS SIGNÉ CONTIENT}\par\vspace{7pt}
  \noindent\begin{minipage}[t]{0.235\linewidth}\poptile{popblue}{Cours}{complet, illustré, pas à pas}\end{minipage}\hfill
  \begin{minipage}[t]{0.235\linewidth}\poptile{poporange}{Méthodes}{et exercices résolus}\end{minipage}\hfill
  \begin{minipage}[t]{0.235\linewidth}\poptile{poppink}{Exercices}{type examen + corrigés}\end{minipage}\hfill
  \begin{minipage}[t]{0.235\linewidth}\poptile{popgreen}{Fiche bilan}{l'essentiel à retenir}\end{minipage}\par}

% Sommaire
\newtcolorbox{sommaire}{enhanced,colback=white,colframe=popink,boxrule=2.2pt,arc=10pt,
  drop shadow=popink,left=18pt,right=18pt,top=13pt,bottom=13pt,before skip=6pt,after skip=20pt,
  before upper={\poppill{Sommaire}{poppurple}{white}\par\vspace{9pt}\popsemi\fontsize{10.6}{14.5}\selectfont\color{popink}}}

% Signature de fin de cours — poussée tout en bas de la dernière page
\newcommand{\findecours}{%
  \par\vfill\nopagebreak
  \begin{center}\popsquiggle{poporange}\end{center}\vspace{4pt}
  \signatureonbuch
}
\AtBeginDocument{\def\llbracket{\mathopen{[\mkern-3.5mu[}}\def\rrbracket{\mathclose{]\mkern-3.5mu]}}} % intervalles d'entiers (glyphe ⟦ absent de la police)
% Glyphes absents de Fira Math : redéfinis à partir de glyphes disponibles
\AtBeginDocument{%
  \def\setminus{\mathbin{\backslash}}\let\smallsetminus\setminus
  \def\vdots{\mathord{\vbox{\baselineskip3.2pt\lineskiplimit0pt\kern4pt\hbox{.}\hbox{.}\hbox{.}}}}%
  \def\ddots{\mathinner{\mkern1mu\raise6.5pt\hbox{.}\mkern2mu\raise3.6pt\hbox{.}\mkern2mu\raise0.7pt\hbox{.}\mkern1mu}}%
  \def\triangle{\mathord{\scalebox{0.9}{$\symup{\Delta}$}}}%
  \def\nabla{\mathord{\raisebox{\depth}{\scalebox{1}[-1]{$\symup{\Delta}$}}}}%
  \def\checkmark{\popcheck{popdark}}%
  \def\star{\mathbin{\ast}}\def\bigstar{\mathord{\ast}}%
  \def\diamond{\mathbin{\scalebox{0.6}{$\Diamond$}}}%
  \def\bigcup{\mathop{\mathchoice{\raisebox{-0.3em}{\scalebox{1.7}{$\cup$}}}{\scalebox{1.25}{$\cup$}}{\cup}{\cup}}\displaylimits}%
  \def\bigcap{\mathop{\mathchoice{\raisebox{-0.3em}{\scalebox{1.7}{$\cap$}}}{\scalebox{1.25}{$\cap$}}{\cap}{\cap}}\displaylimits}%
  \def\lesseqqgtr{\mathrel{\lessgtr}}\def\bigoplus{\mathop{\oplus}\displaylimits}%
}
\providecommand{\ssi}{si et seulement si} % abréviation fréquente
\AtBeginDocument{\providecommand{\wideparen}{\overparen}} % arc de cercle

%% FIGFIX-BEGIN v3 — correctifs globaux des figures (docs/onbuch-generation/tools/figfix)
\usepackage{adjustbox}\usepackage{etoolbox}
% toute figure TikZ/pgfplots plus large que la ligne est réduite à la largeur disponible
\BeforeBeginEnvironment{tikzpicture}{\begin{adjustbox}{max width=\linewidth}}
\AfterEndEnvironment{tikzpicture}{\end{adjustbox}}
% légendes pgfplots lisibles par-dessus les courbes
\pgfplotsset{every axis/.append style={legend style={fill=white,fill opacity=0.92,text opacity=1,draw=black!15,inner sep=3pt}}}
% étiquettes d'arêtes (syntaxe edge["..."]) avec fond blanc
\tikzset{every edge quotes/.append style={fill=white,inner sep=1.2pt}}
% bulles de cartes mentales : texte à la ligne dans la bulle, bulle un peu plus grande
\tikzset{every concept/.append style={text width=2.0cm,align=center,minimum size=2.3cm,inner sep=2pt}}
%% FIGFIX-END
\begin{document}

\pagegarde{Leçon 19 — Vocabulaire et compréhension de texte : vie culturelle au Cameroun et en Chine}{Chinois}{1ère A}{Module 3 : Citoyenneté et ouverture au monde}{ZHA19}{2 h}{Cette leçon permet de découvrir et d'employer le vocabulaire chinois de la vie culturelle (fêtes, arts, traditions) à travers un texte support original comparant le Cameroun et la Chine. L'élève apprend à lire, comprendre et commenter un court texte culturel en chinois simplifié.
\vspace{4pt}
\begin{multicols}{2}\small
\begin{itemize}
\item À la fin de cette leçon, je sais lire à voix haute un court texte chinois sur la vie culturelle en respectant les tons.
\item Je sais reconnaître, prononcer et écrire au moins 15 mots de vocabulaire liés aux fêtes, arts et traditions.
\item Je sais identifier les clés (radicaux) et l'ordre des traits de quelques caractères du thème.
\item Je sais répondre à des questions de compréhension portant sur le texte support.
\item Je sais employer les collocations usuelles du thème (过节日, 表演节目, 庆祝...).
\item Je sais comparer oralement et par écrit une fête camerounaise et une fête chinoise.
\end{itemize}\end{multicols}}

\begin{sommaire}
\begin{multicols}{2}
\begin{enumerate}
\item Mise en route et découverte du thème
\item Texte support : lecture et traduction
\item Glossaire du thème
\item Clés (radicaux) et ordre des traits
\item Compréhension du texte et collocations
\item Production guidée et ouverture interculturelle
\item Activité d'intégration
\item Méthodes et exercices résolus
\item Exercices
\item Corrigés des exercices
\item Fiche bilan
\end{enumerate}
\end{multicols}
\end{sommaire}

\begin{objectifs}
\begin{itemize}
\item À la fin de cette leçon, je sais lire à voix haute un court texte chinois sur la vie culturelle en respectant les tons.
\item Je sais reconnaître, prononcer et écrire au moins 15 mots de vocabulaire liés aux fêtes, arts et traditions.
\item Je sais identifier les clés (radicaux) et l'ordre des traits de quelques caractères du thème.
\item Je sais répondre à des questions de compréhension portant sur le texte support.
\item Je sais employer les collocations usuelles du thème (过节日, 表演节目, 庆祝...).
\item Je sais comparer oralement et par écrit une fête camerounaise et une fête chinoise.
\item Je sais produire un court paragraphe guidé présentant une fête culturelle de mon pays.
\end{itemize}
\end{objectifs}

\begin{prerequis}
\begin{itemize}
\item Vocabulaire des pays, nationalités et langues (Seconde : 中国, 喀麦隆, 人, 语)
\item Structures de base : phrase sujet-verbe-objet, 是, 有, 很
\item Interrogation simple : 什么, 吗, 呢
\item Notions de temps : 年, 月, 日, 时候
\item Radicaux de base déjà rencontrés (人, 口, 女, 木)
\end{itemize}
\end{prerequis}

\newpage

\coursec{Mise en route et découverte du thème}

\courssub{Situation d'accroche : la fête de la jeunesse au lycée de Yaoundé}

C'est le 11 février, jour de la fête de la jeunesse. Dans la cour du lycée de Yaoundé, tout est en effervescence : un groupe d'élèves exécute une danse traditionnelle bamiléké au rythme des tambours, une classe de Terminale joue un sketch comique, et des stands proposent du ndolé, des brochettes et des beignets. Les drapeaux camerounais flottent partout.

Parmi les invités se trouve \zh{李明} (Lǐ Míng), un correspondant chinois qui étudie le français à Yaoundé. Il observe, fasciné, puis se tourne vers son amie Aïcha et lui demande :

\zh{你们国家有什么传统节日？}

\emph{Nǐmen guójiā yǒu shénme chuántǒng jiérì ?}

« Quelles fêtes traditionnelles y a-t-il dans votre pays ? »

Aïcha voudrait lui parler de la fête du Ngondo chez les Sawa de Douala, des danses bamiléké, de la fête nationale du 20 mai avec ses défilés. Mais elle voudrait aussi découvrir la grande fête chinoise : la fête du Printemps, avec ses lanternes rouges et ses danses du dragon. Problème : comment dire tout cela en chinois ?

\textbf{Problématique de la leçon :} pour échanger sur la vie culturelle du Cameroun et de la Chine, quels mots et quelles expressions faut-il connaître ? C'est précisément l'objectif de cette leçon : construire le vocabulaire des fêtes, des arts et des traditions, puis lire et comprendre un texte chinois sur ce thème.

\courssub{Brainstorming : les fêtes et activités culturelles que vous connaissez}

Avant d'ouvrir le dictionnaire chinois, faisons le point en français. L'enseignant écrit au tableau la question : \emph{Quelles fêtes et activités culturelles connaissez-vous au Cameroun ?} Chaque élève propose une idée. Voici les réponses attendues, déjà traduites en chinois pour préparer la suite :

\begin{center}
\begin{tabularx}{\linewidth}{|X|X|X|X|}
\hline
\rowcolor{popblueL} Caractères & Pinyin & Nature & Traduction \\
\hline
文化 & wénhuà & nom & la culture \\
\hline
节日 & jiérì & nom & la fête \\
\hline
传统节日 & chuántǒng jiérì & locution nominale & la fête traditionnelle \\
\hline
国庆节 & guóqìngjié & nom & la fête nationale \\
\hline
青年节 & qīngniánjié & nom & la fête de la jeunesse \\
\hline
舞蹈 & wǔdǎo & nom & la danse \\
\hline
音乐 & yīnyuè & nom & la musique \\
\hline
艺术 & yìshù & nom & l'art \\
\hline
春节 & Chūnjié & nom propre & la fête du Printemps (Nouvel An chinois) \\
\hline
\end{tabularx}
\end{center}

Observez bien la première colonne : beaucoup de mots chinois des fêtes contiennent le caractère \zh{节} (jié). C'est un indice précieux que nous exploiterons dans la section sur les radicaux : en chinois, les mots d'une même famille partagent souvent un même caractère, comme en français \emph{fête}, \emph{festif}, \emph{festival} partagent la même racine. Ainsi \zh{节日} (la fête), \zh{国庆节} (la fête nationale), \zh{青年节} (la fête de la jeunesse) et \zh{春节} (la fête du Printemps) forment une petite famille de mots facile à mémoriser.

\begin{definition}[Le mot-clé de la leçon : 文化]
\zh{文化} (wénhuà) signifie « la culture ». C'est le mot central de cette leçon : les fêtes (\zh{节日}), la musique (\zh{音乐}), la danse (\zh{舞蹈}) et l'art (\zh{艺术}) sont tous des éléments de la \cle{文化} d'un pays. On dit par exemple : \zh{中国文化} (Zhōngguó wénhuà, « la culture chinoise ») et \zh{喀麦隆文化} (Kāmàilóng wénhuà, « la culture camerounaise »).
\end{definition}

\begin{textebox}[Les deux mots à retenir dès maintenant]
\zh{文化} — wénhuà — la culture.\\
\zh{节日} — jiérì — la fête.\\
Phrase modèle : \zh{中国和喀麦隆都有很多传统节日。}\\
\emph{Zhōngguó hé Kāmàilóng dōu yǒu hěn duō chuántǒng jiérì.}\\
« La Chine et le Cameroun ont toutes les deux beaucoup de fêtes traditionnelles. »
\end{textebox}

Remarquez la structure de la phrase modèle : \textbf{Sujet 1 + 和 + Sujet 2 + 都 + verbe}. Le mot \zh{和} (hé) signifie « et » entre deux noms, et l'adverbe \zh{都} (dōu, « tous les deux, tous ») se place toujours \emph{avant} le verbe, jamais après. C'est une structure de base du niveau HSK 1 que nous réutiliserons dans tout le texte support.

\courssub{Observation d'images : deux mondes de fêtes}

L'enseignant projette (ou décrit) deux séries d'images.

\textbf{Première série : la Chine.} Des lanternes rouges suspendues dans les rues, une danse du dragon avec un long dragon de tissu porté par des dizaines de personnes, des familles réunies autour d'un grand repas, des enveloppes rouges offertes aux enfants. C'est \zh{春节} (Chūnjié), la fête du Printemps, c'est-à-dire le Nouvel An chinois.

\textbf{Deuxième série : le Cameroun.} Le défilé du 20 mai sur le boulevard du 20 mai à Yaoundé, les pirogues décorées de la fête du Ngondo à Douala, les danseurs bamiléké avec leurs masques et leurs costumes colorés, les élèves en uniforme chantant lors de la fête de la jeunesse.

L'enseignant demande alors : \emph{Qu'est-ce que ces images ont en commun ?} Réponse attendue : dans les deux pays, les fêtes rassemblent les gens, on danse, on mange, on porte de beaux vêtements, on transmet des traditions. En chinois, on peut déjà dire :

\zh{中国人喜欢春节，喀麦隆人喜欢国庆节。}

\emph{Zhōngguórén xǐhuan Chūnjié, Kāmàilóngrén xǐhuan guóqìngjié.}

« Les Chinois aiment la fête du Printemps, les Camerounais aiment la fête nationale. »

Notez au passage comment on forme le nom des habitants : \textbf{pays + 人} (rén, « personne ») : \zh{中国人} (Zhōngguórén, « les Chinois »), \zh{喀麦隆人} (Kāmàilóngrén, « les Camerounais »). Le verbe \zh{喜欢} (xǐhuan, « aimer ») se construit directement avec un nom : \zh{喜欢 + nom}.

\begin{savaistu}[La fête du Printemps, un géant des fêtes]
La fête du Printemps (\zh{春节}, Chūnjié) est la fête la plus importante de Chine. Elle est célébrée par plus d'un milliard de personnes, en Chine et dans les communautés chinoises du monde entier. Les familles se réunissent, décorent leurs maisons en rouge (couleur porte-bonheur), mangent des raviolis (\zh{饺子}, jiǎozi) et regardent ensemble la grande soirée télévisée du Nouvel An. On peut la comparer au 20 mai au Cameroun : une fête qui rassemble tout le pays, des grandes villes aux plus petits villages, autour des mêmes symboles et de la même fierté nationale.
\end{savaistu}

\courssub{Carte mentale : le vocabulaire de la culture}

Pour organiser le vocabulaire du thème, construisons ensemble une carte mentale autour du mot-clé \zh{文化}. Cette carte sera complétée au fur et à mesure de la leçon.

\begin{popfigure}
\begin{tikzpicture}[mindmap, grow cyclic, every node/.style={concept, align=center, font=\small}, root concept/.append style={concept color=popblue, font=\large\bfseries}, level 1/.append style={level distance=3.4cm, sibling angle=90}]
\node[root concept, align=center]{文化\\ wénhuà\\ la culture}
  child[concept color=poporange] { node[align=center]{节日\\ jiérì\\ les fêtes} }
  child[concept color=popgreen] { node[align=center]{音乐\\ yīnyuè\\ la musique} }
  child[concept color=poppurple] { node[align=center]{舞蹈\\ wǔdǎo\\ la danse} }
  child[concept color=popgold] { node[align=center]{艺术\\ yìshù\\ l'art} };
\end{tikzpicture}
\legende{Carte mentale du thème : la culture et ses quatre grandes branches.}
\end{popfigure}

\begin{methode}[Aborder un texte chinois sans paniquer]
Face à un texte en caractères chinois, beaucoup d'élèves francophones veulent traduire mot à mot dès la première ligne. Mauvaise idée ! Voici la bonne démarche, en trois étapes :
\begin{enumerate}
\item \textbf{Écouter} d'abord le texte lu par l'enseignant, sans support écrit, pour saisir le rythme, les tons et le sens global.
\item \textbf{Repérer} ensuite, à la lecture, les mots déjà connus (noms de pays, mots du thème comme \zh{节日} ou \zh{文化}) : ce sont des « îlots de compréhension ».
\item \textbf{Deviner} le sens global à partir de ces îlots et du contexte, \emph{avant} de traduire phrase par phrase. Le sens global guide la traduction, jamais l'inverse.
\end{enumerate}
\end{methode}

\courssub{Première écoute du texte support (sans support écrit)}

Fermez vos cahiers. L'enseignant lit le texte support deux fois, à vitesse normale puis lente. Votre tâche : écouter et noter en français les idées que vous croyez comprendre. Voici le texte que l'enseignant lit (vous ne le regarderez qu'à la section suivante) :

\begin{textebox}[Texte lu par l'enseignant — écoute seule]
\zh{中国和喀麦隆都有很丰富的文化。在中国，最重要的节日是春节。春节的时候，家家户户挂红灯笼，吃饺子，看舞龙舞狮，非常热闹。在喀麦隆，五月二十日是国庆节，人们看阅兵，唱歌跳舞。杜阿拉的萨瓦人还有恩贡多节，他们在水上庆祝。两个国家的节日不一样，但是人们都一样快乐。文化让我们互相了解，成为朋友。}\\
\emph{Zhōngguó hé Kāmàilóng dōu yǒu hěn fēngfù de wénhuà. Zài Zhōngguó, zuì zhòngyào de jiérì shì Chūnjié. Chūnjié de shíhou, jiājiā-hùhù guà hóng dēnglong, chī jiǎozi, kàn wǔlóng wǔshī, fēicháng rènao. Zài Kāmàilóng, wǔ yuè èrshí rì shì guóqìngjié, rénmen kàn yuèbīng, chànggē tiàowǔ. Dù'ālā de Sàwǎrén hái yǒu Ēngòngduōjié, tāmen zài shuǐ shàng qìngzhù. Liǎng ge guójiā de jiérì bù yīyàng, dànshì rénmen dōu yīyàng kuàilè. Wénhuà ràng wǒmen hùxiāng liǎojiě, chéngwéi péngyou.}
\end{textebox}

Traduction (à ne montrer qu'après l'écoute) : « La Chine et le Cameroun ont toutes les deux une culture très riche. En Chine, la fête la plus importante est la fête du Printemps. Pendant la fête du Printemps, chaque famille accroche des lanternes rouges, mange des raviolis et regarde les danses du dragon et du lion : c'est très animé. Au Cameroun, le 20 mai est la fête nationale ; les gens regardent le défilé militaire, chantent et dansent. Les Sawa de Douala ont aussi la fête du Ngondo, qu'ils célèbrent sur l'eau. Les fêtes des deux pays sont différentes, mais les gens sont tout aussi heureux. La culture nous permet de nous comprendre mutuellement et de devenir amis. »

\courssub{Hypothèses des élèves : qu'avez-vous compris ?}

Après la double écoute, l'enseignant interroge la classe : \emph{De quoi parle ce texte ? Quels mots avez-vous reconnus ?} Les hypothèses attendues :

\begin{itemize}
\item Le texte parle de la \cle{Chine} et du \cle{Cameroun} (on a entendu \zh{中国} et \zh{喀麦隆}).
\item Il parle de \cle{fêtes} (on a reconnu \zh{节日} et peut-être \zh{春节}).
\item On a peut-être entendu une date : \zh{五月二十日} (wǔ yuè èrshí rì), le 20 mai.
\item Le texte semble \cle{comparer} les deux pays.
\end{itemize}

Toutes les hypothèses, même incomplètes, sont acceptées : elles seront vérifiées à la lecture, dans la section suivante. C'est exactement la démarche d'un bon lecteur : émettre des hypothèses, puis les confirmer ou les corriger.

\begin{exoresolu}[Reconnaître les mots du thème à l'écoute]
Énoncé : l'enseignant prononce cinq mots : \emph{wénhuà, jiérì, yīnyuè, wǔdǎo, yìshù}. Associez chaque mot à sa traduction : a) la danse, b) la culture, c) l'art, d) la fête, e) la musique.
\tcblower
Solution détaillée : \emph{wénhuà} (\zh{文化}) = b) la culture ; \emph{jiérì} (\zh{节日}) = d) la fête ; \emph{yīnyuè} (\zh{音乐}) = e) la musique ; \emph{wǔdǎo} (\zh{舞蹈}) = a) la danse ; \emph{yìshù} (\zh{艺术}) = c) l'art. Astuce : \zh{乐} (yuè) évoque la musique et la joie, et \zh{舞} (wǔ) se retrouve dans tous les mots liés à la danse : \zh{跳舞} (tiàowǔ, « danser »), \zh{舞蹈} (wǔdǎo, « la danse »), \zh{舞龙} (wǔlóng, « la danse du dragon »).
\end{exoresolu}

\begin{attention}[Pièges de prononciation pour les francophones]
\begin{itemize}
\item \zh{节日} (jiérì) : le \emph{j} chinois ne se prononce pas comme le « j » français de « jour » ; c'est un son proche de « tsie », avec la langue contre le palais. Et \emph{rì} ne se dit pas « ri » à la française : la langue est recourbée vers l'arrière.
\item \zh{音乐} (yīnyuè) : attention aux tons ! Premier ton sur \emph{yīn} (haut et plat), quatrième ton sur \emph{yuè} (descendant). Un mauvais ton peut rendre le mot incompréhensible.
\item \zh{文化} (wénhuà) : \emph{wén} au deuxième ton (montant, comme une question), \emph{huà} au quatrième ton (descendant, sec).
\item Ne confondez pas \zh{节日} (jiérì, la fête) et \zh{节目} (jiémù, l'émission, le numéro de spectacle) : même premier caractère, sens différents !
\item \zh{热闹} (rènao, « animé, bruyant de joie ») : le second caractère se prononce avec un ton léger, presque sans ton. Cet adjectif est parfait pour décrire une fête camerounaise ou chinoise.
\end{itemize}
\end{attention}

\begin{experience}[Jeu de rôle d'ouverture : Aïcha et Li Ming]
Par groupes de deux, rejouez la scène d'accroche. L'élève A (Li Ming) pose la question : \zh{你们国家有什么传统节日？} (Nǐmen guójiā yǒu shénme chuántǒng jiérì ?). L'élève B (Aïcha) répond avec les mots de la carte mentale, par exemple : \zh{我们国家有国庆节和青年节。} (Wǒmen guójiā yǒu guóqìngjié hé qīngniánjié. — « Dans notre pays, il y a la fête nationale et la fête de la jeunesse. »). Puis échangez les rôles : Aïcha interroge Li Ming sur la Chine, et Li Ming répond : \zh{中国有春节。} (Zhōngguó yǒu Chūnjié. — « La Chine a la fête du Printemps. »). Objectif : prononcer correctement les tons et mémoriser les cinq mots-clés avant la lecture du texte.
\end{experience}

\begin{aretenir}[Objectifs de la leçon]
À l'issue de cette leçon, vous serez capables de :
\begin{enumerate}
\item reconnaître, prononcer (pinyin et tons) et employer le vocabulaire de la vie culturelle : \zh{文化}, \zh{节日}, \zh{音乐}, \zh{舞蹈}, \zh{艺术}, \zh{春节}, \zh{国庆节} ;
\item lire et comprendre un court texte chinois comparant la vie culturelle du Cameroun et de la Chine ;
\item identifier les clés (radicaux) et l'ordre des traits de quelques caractères du thème ;
\item répondre à des questions de compréhension et parler d'une fête de votre pays en chinois.
\end{enumerate}
Mot-clé à retenir dès maintenant : \cle{文化} (wénhuà), la culture — le fil conducteur de toute la leçon.
\end{aretenir}

\coursec{Texte support : lecture et traduction}

\courssub{De la mise en route à la lecture}

Dans la section précédente, nous avons découvert le thème de la vie culturelle au Cameroun et en Chine à travers des images et quelques mots-clés : \zh{春节} (Chūnjié, la fête du Printemps), \zh{舞龙舞狮} (wǔlóng wǔshī, la danse du dragon et du lion), \zh{传统节日} (chuántǒng jiérì, les fêtes traditionnelles). Nous allons maintenant lire un court texte qui met en parallèle les deux pays. L'objectif est double : comprendre le sens global du texte, puis observer comment la langue chinoise exprime des idées que le français formule autrement.

\courssub{Première lecture : écoute et sens global}

\begin{methode}[Comment aborder un texte chinois ?]
\begin{enumerate}
\item \textbf{Lire d'abord en silence} les caractères, sans chercher à tout traduire : repérer les mots déjà connus.
\item \textbf{Écouter le professeur} (ou l'enregistrement) en suivant le pinyin du doigt.
\item \textbf{Lire à voix haute} phrase par phrase, en soignant les tons.
\item \textbf{Traduire phrase par phrase}, puis reformuler le texte entier en français avec ses propres mots.
\end{enumerate}
\end{methode}

\begin{textebox}[中国和喀麦隆的文化生活 — La vie culturelle en Chine et au Cameroun]
中国和喀麦隆都有丰富的文化生活。\\
Zhōngguó hé Kāmàilóng dōu yǒu fēngfù de wénhuà shēnghuó.\\
《La Chine et le Cameroun ont tous deux une vie culturelle riche.》\\[4pt]
中国人过春节，全家一起吃饭、看舞龙舞狮。\\
Zhōngguórén guò Chūnjié, quánjiā yìqǐ chīfàn, kàn wǔlóng wǔshī.\\
《Les Chinois célèbrent la fête du Printemps ; toute la famille mange ensemble et regarde la danse du dragon et du lion.》\\[4pt]
喀麦隆人有很多传统节日，人们唱歌、跳舞、穿漂亮的传统服装。\\
Kāmàilóngrén yǒu hěn duō chuántǒng jiérì, rénmen chànggē, tiàowǔ, chuān piàoliang de chuántǒng fúzhuāng.\\
《Les Camerounais ont beaucoup de fêtes traditionnelles ; les gens chantent, dansent et portent de beaux vêtements traditionnels.》\\[4pt]
两个国家的人民都很喜欢音乐和舞蹈。\\
Liǎng ge guójiā de rénmín dōu hěn xǐhuan yīnyuè hé wǔdǎo.\\
《Les peuples des deux pays aiment tous beaucoup la musique et la danse.》\\[4pt]
文化让我们互相了解。\\
Wénhuà ràng wǒmen hùxiāng liǎojiě.\\
《La culture nous permet de nous comprendre mutuellement.》
\end{textebox}

\courssub{Traduction commentée, phrase par phrase}

\textbf{Phrase 1 :} \zh{中国和喀麦隆都有丰富的文化生活。}

\emph{Zhōngguó hé Kāmàilóng dōu yǒu fēngfù de wénhuà shēnghuó.}

Structure : Sujet (\zh{中国和喀麦隆}) + \zh{都} + verbe \zh{有} + complément (\zh{丰富的文化生活}). Le mot \zh{和} (hé) signifie « et » et relie deux noms, jamais deux phrases. L'adjectif \zh{丰富} (fēngfù, « riche, abondant ») qualifie le nom grâce à la particule \zh{的} (de).

\textbf{Phrase 2 :} \zh{中国人过春节，全家一起吃饭、看舞龙舞狮。}

\emph{Zhōngguórén guò Chūnjié, quánjiā yìqǐ chīfàn, kàn wǔlóng wǔshī.}

Le verbe \zh{过} (guò) suivi d'une fête signifie « célébrer, passer (une fête) ». \zh{全家} (quánjiā) signifie « toute la famille » ; \zh{一起} (yìqǐ) signifie « ensemble » et se place avant le verbe. La virgule chinoise \zh{、} (dùnhào) sépare des verbes ou des noms coordonnés dans une énumération : \zh{吃饭、看舞龙舞狮}.

\textbf{Phrase 3 :} \zh{喀麦隆人有很多传统节日，人们唱歌、跳舞、穿漂亮的传统服装。}

\emph{Kāmàilóngrén yǒu hěn duō chuántǒng jiérì, rénmen chànggē, tiàowǔ, chuān piàoliang de chuántǒng fúzhuāng.}

\zh{很多} (hěn duō) = « beaucoup de » ; \zh{人们} (rénmen) = « les gens », avec le suffixe de pluriel \zh{们} prononcé au ton neutre. Le verbe \zh{穿} (chuān) signifie « porter, mettre (un vêtement) ».

\textbf{Phrase 4 :} \zh{两个国家的人民都很喜欢音乐和舞蹈。}

\emph{Liǎng ge guójiā de rénmín dōu hěn xǐhuan yīnyuè hé wǔdǎo.}

Devant un classificateur, « deux » se dit \zh{两} (liǎng) et non \zh{二} (èr) : \zh{两个国家} « les deux pays ». Le mot \zh{人民} (rénmín) signifie « le peuple ».

\textbf{Phrase 5 :} \zh{文化让我们互相了解。}

\emph{Wénhuà ràng wǒmen hùxiāng liǎojiě.}

Le verbe \zh{让} (ràng) signifie « faire, laisser, permettre » : Sujet + \zh{让} + personne + verbe. L'adverbe \zh{互相} (hùxiāng) signifie « mutuellement, l'un l'autre » et se place avant le verbe \zh{了解} (liǎojiě, « comprendre, connaître »).

\courssub{Grammaire observée dans le texte}

\begin{propriete}[L'adverbe 都 (dōu) : « tous, tous les deux »]
\zh{都} se place \textbf{avant le verbe} (ou l'adjectif) et indique que l'action ou la qualité concerne la totalité du sujet : « tous », « tous les deux ».\\
Structure : Sujet (pluriel) + \zh{都} + Verbe.\\
Exemples :\\
\zh{我们都喜欢跳舞。} \emph{Wǒmen dōu xǐhuan tiàowǔ.} « Nous aimons tous danser. »\\
\zh{两个国家的人民都很喜欢音乐。} \emph{Liǎng ge guójiā de rénmín dōu hěn xǐhuan yīnyuè.} « Les peuples des deux pays aiment tous la musique. »\\
Attention : \zh{都} ne se traduit pas toujours mot à mot ; il souligne l'idée de totalité.
\end{propriete}

\begin{attention}[过 (guò) + fête : ne pas traduire par « traverser »]
Dans \zh{过春节} (guò Chūnjié), le verbe \zh{过} signifie « célébrer, passer (une fête, un anniversaire) », et non « traverser ».\\
Exemples : \zh{过生日} (guò shēngrì) « fêter son anniversaire » ; \zh{过新年} (guò xīnnián) « célébrer le Nouvel An ».\\
Erreur fréquente des francophones : traduire \zh{中国人过春节} par « Les Chinois traversent le Printemps ». Retenez : \cle{过 + fête = célébrer}.
\end{attention}

\begin{propriete}[La structure 让 (ràng) : causatif / permissif]
\zh{让} exprime la causation ou la permission. Structure : \textbf{Sujet 1 + 让 + Sujet 2 (personne) + Verbe}.\\
\zh{文化让我们互相了解。} La culture (S1) fait/permet que nous (S2) nous comprenions mutuellement (V).\\
Autres exemples : \zh{老师让我们写字。} (Lǎoshī ràng wǒmen xiě zì.) « Le professeur fait écrire les élèves / demande aux élèves d'écrire. » ; \zh{请让我去。} (Qǐng ràng wǒ qù.) « Permettez-moi d'y aller. »
\end{propriete}

\begin{exemplebox}[Phrases modèles à réemployer]
\zh{人们唱歌、跳舞。} \emph{Rénmen chànggē, tiàowǔ.} « Les gens chantent et dansent. »\\
\zh{文化让我们互相了解。} \emph{Wénhuà ràng wǒmen hùxiāng liǎojiě.} « La culture nous permet de nous comprendre mutuellement. »\\
\zh{全家一起吃饭。} \emph{Quánjiā yìqǐ chīfàn.} « Toute la famille mange ensemble. »
\end{exemplebox}

\courssub{Lecture chorale puis individuelle : le travail des tons}

Le professeur lit d'abord chaque phrase, la classe répète en chœur, puis des élèves lisent individuellement. Trois points de prononciation exigent une attention particulière :

\begin{itemize}
\item \textbf{Le 3e ton de \zh{有} (yǒu) et de \zh{很} (hěn)} : le 3e ton descend puis remonte. Devant un autre 3e ton, le premier devient un 2e ton : \zh{很好} se prononce \emph{hén hǎo}. Dans \zh{有丰富} (yǒu fēngfù) et \zh{很多} (hěn duō), le 3e ton est souvent réalisé « à moitié » (bas, sans remontée) devant un ton non-3e.
\item \textbf{Le ton neutre de \zh{们} dans \zh{人们} (rénmen)} : le suffixe \zh{们} se prononce court et léger, sans ton propre. Même chose dans \zh{我们} (wǒmen).
\item \textbf{La modification de \zh{一} dans \zh{一起} (yìqǐ)} : devant un 4e ton, \zh{一} (yī) passe au 2e ton : on prononce \emph{yíqǐ}.
\end{itemize}

\begin{experience}[Lecture en écho]
\begin{enumerate}
\item Lecture chorale : la classe répète chaque phrase après le professeur, en frappant doucement le rythme des tons sur la table.
\item Lecture individuelle : chaque élève lit une phrase ; les autres vérifient les tons de \zh{有}, \zh{很}, \zh{人们}.
\item Jeu du « détecteur de tons » : un élève lit, la classe lève la main dès qu'un ton est faux, puis corrige.
\end{enumerate}
\end{experience}

\courssub{Repérage : mots connus et mots nouveaux}

\begin{center}
\begin{tabularx}{\linewidth}{|X|X|X|}
\hline
\rowcolor{popblueL} Catégorie & Mots du texte & Sens \\
\hline
Mots déjà connus & \zh{中国} \zh{人} \zh{有} \zh{吃} \zh{看} \zh{很} \zh{多} \zh{喜欢} \zh{和} & Chine, personne, avoir, manger, regarder, très, beaucoup, aimer, et \\
\hline
Mots nouveaux (1) & \zh{丰富} \zh{文化} \zh{生活} \zh{过} \zh{春节} \zh{全家} \zh{一起} & riche, culture, vie, célébrer, fête du Printemps, toute la famille, ensemble \\
\hline
Mots nouveaux (2) & \zh{传统} \zh{节日} \zh{人们} \zh{唱歌} \zh{跳舞} \zh{穿} \zh{漂亮} \zh{服装} & traditionnel, fête, les gens, chanter, danser, porter, beau, vêtement \\
\hline
Mots nouveaux (3) & \zh{国家} \zh{人民} \zh{音乐} \zh{舞蹈} \zh{让} \zh{互相} \zh{了解} & pays, peuple, musique, danse, permettre, mutuellement, comprendre \\
\hline
\end{tabularx}
\end{center}

\begin{popfigure}
\begin{tikzpicture}
\node[draw=popblue, fill=popblueL, rounded corners, align=center, minimum width=4.2cm, minimum height=2.6cm] (chine) at (-3.6,0) {\textbf{中国 Chine}\\[3pt] \zh{春节} Chūnjié\\ \zh{舞龙舞狮} wǔlóng wǔshī};
\node[draw=poporange, fill=poporangeL, rounded corners, align=center, minimum width=4.2cm, minimum height=2.6cm] (cam) at (3.6,0) {\textbf{喀麦隆 Cameroun}\\[3pt] \zh{传统节日} chuántǒng jiérì\\ \zh{唱歌跳舞} chànggē tiàowǔ};
\node[draw=popgreen, fill=popgreenL, rounded corners, align=center, minimum width=3.4cm] (commun) at (0,-3) {\textbf{Élément commun}\\ \zh{音乐和舞蹈}\\ yīnyuè hé wǔdǎo};
\draw[-{Stealth}, thick, popgreen] (chine.south) -- (commun.north west);
\draw[-{Stealth}, thick, popgreen] (cam.south) -- (commun.north east);
\end{tikzpicture}
\legende{Carte mentale du texte : deux pays, deux traditions, une passion commune.}
\end{popfigure}

\courssub{Exercice résolu et reformulation}

\begin{exoresolu}[Traduire et retrouver la phrase du texte]
Énoncé : traduisez en chinois la phrase « Les Camerounais ont beaucoup de fêtes traditionnelles », puis retrouvez dans le texte la phrase qui exprime « La culture nous permet de nous comprendre mutuellement ».
\tcblower
Solution détaillée :\\
1. « Les Camerounais » = \zh{喀麦隆人} ; « avoir » = \zh{有} ; « beaucoup de » = \zh{很多} ; « fêtes traditionnelles » = \zh{传统节日}. On obtient : \zh{喀麦隆人有很多传统节日。} \emph{Kāmàilóngrén yǒu hěn duō chuántǒng jiérì.}\\
2. La phrase cherchée est la dernière du texte : \zh{文化让我们互相了解。} \emph{Wénhuà ràng wǒmen hùxiāng liǎojiě.} On reconnaît la structure Sujet + \zh{让} + personne + verbe, avec l'adverbe \zh{互相} placé avant \zh{了解}.
\end{exoresolu}

\begin{methode}[Reformuler le texte en français]
\begin{enumerate}
\item Relire le texte chinois une dernière fois, sans le pinyin si possible.
\item Fermer le livre et raconter le contenu en français, à l'oral, en trois ou quatre phrases : qui ? quoi ? quelle fête ? quel point commun ?
\item Vérifier ensuite que les idées essentielles sont présentes : la fête du Printemps en Chine, les fêtes traditionnelles au Cameroun, l'amour commun de la musique et de la danse, le rôle de la culture.
\end{enumerate}
Exemple de reformulation attendue : « En Chine, on célèbre la fête du Printemps en famille, avec des repas et des danses du dragon. Au Cameroun, les fêtes traditionnelles rassemblent les gens qui chantent, dansent et portent de beaux vêtements. Dans les deux pays, la musique et la danse rapprochent les peuples. »
\end{methode}

\begin{attention}[Pièges de traduction relevés dans ce texte]
\begin{itemize}
\item Ne pas traduire \zh{都} par « tous » à la fin de la phrase : en chinois, il se place avant le verbe, jamais après.
\item \zh{两个国家} : devant un classificateur, « deux » se dit \zh{两}, jamais \zh{二}.
\item \zh{全家一起吃饭} : \zh{一起} se place avant le verbe \zh{吃}, contrairement au français « manger ensemble » où « ensemble » suit le verbe.
\item \zh{人们} est déjà pluriel grâce à \zh{们} ; inutile d'ajouter un autre marqueur.
\end{itemize}
\end{attention}

\begin{aretenir}[L'essentiel de cette section]
\begin{itemize}
\item Le texte compare la vie culturelle de la Chine (\zh{春节}, \zh{舞龙舞狮}) et du Cameroun (\zh{传统节日}, \zh{唱歌跳舞}), unies par l'amour de \zh{音乐和舞蹈}.
\item \zh{都} (dōu) avant le verbe = « tous » ; \zh{过} + fête = « célébrer » ; \zh{让} = « permettre, faire » ; \zh{互相} = « mutuellement ».
\item Soigner le 3e ton de \zh{有} et \zh{很}, le ton neutre de \zh{们}, et la modification de \zh{一} dans \zh{一起} (yìqǐ).
\end{itemize}
\end{aretenir}

% ============================================================
% SECTION 3 : GLOSSAIRE DU THÈME (Vocabulaire complet)
% ============================================================
\coursec{Glossaire du thème : vie culturelle}

\courssub{Vocabulaire de base (HSK 1--3) -- Tableau de référence}

\begin{center}
\begin{tabularx}{\linewidth}{|X|X|X|X|}
\hline
\rowcolor{popblueL} \textbf{Caractères} & \textbf{Pinyin} & \textbf{Nature} & \textbf{Traduction} \\
\hline
中国 & Zhōngguó & n. & Chine \\
\hline
喀麦隆 & Kāmàilóng & n. & Cameroun \\
\hline
和 & hé & conj. & et (relie deux noms) \\
\hline
都 & dōu & adv. & tous, tous les deux (place avant le verbe) \\
\hline
有 & yǒu & v. & avoir, il y a \\
\hline
丰富 & fēngfù & adj. & riche, abondant, varié \\
\hline
文化 & wénhuà & n. & culture \\
\hline
生活 & shēnghuó & n./v. & vie, vivre \\
\hline
人 & rén & n. & personne, gens \\
\hline
过 & guò & v. & passer (le temps), célébrer (une fête) \\
\hline
春节 & Chūnjié & n. & Fête du Printemps (Nouvel An chinois) \\
\hline
全家 & quánjiā & n. & toute la famille \\
\hline
一起 & yìqǐ & adv. & ensemble (place avant le verbe) \\
\hline
吃饭 & chīfàn & v.o. & manger (un repas) \\
\hline
看 & kàn & v. & regarder, voir \\
\hline
舞龙舞狮 & wǔlóng wǔshī & v.o. & danser le dragon et le lion (danse traditionnelle) \\
\hline
很多 & hěn duō & adj./adv. & beaucoup (de) \\
\hline
传统 & chuántǒng & adj./n. & tradition, traditionnel \\
\hline
节日 & jiérì & n. & fête, jour de fête \\
\hline
人们 & rénmen & n. & les gens, les personnes (pluriel de 人) \\
\hline
唱歌 & chànggē & v.o. & chanter \\
\hline
跳舞 & tiàowǔ & v.o. & danser \\
\hline
穿 & chuān & v. & porter (vêtements, chaussures) \\
\hline
漂亮 & piàoliang & adj. & beau, joli \\
\hline
服装 & fúzhuāng & n. & vêtements, tenue, costume \\
\hline
两 & liǎng & num. & deux (devant classificateur) \\
\hline
个 & gè & m. & classificateur général \\
\hline
国家 & guójiā & n. & pays, État, nation \\
\hline
人民 & rénmín & n. & le peuple, la population \\
\hline
音乐 & yīnyuè & n. & musique \\
\hline
舞蹈 & wǔdǎo & n. & danse (art, spectacle) \\
\hline
让 & ràng & v. & laisser, faire, permettre (causatif) \\
\hline
互相 & hùxiāng & adv. & mutuellement, l'un l'autre \\
\hline
了解 & liǎojiě & v. & comprendre, connaître (une personne, une situation) \\
\hline
\end{tabularx}
\end{center}

\courssub{Collocations utiles (associations de mots fréquentes)}

\begin{center}
\begin{tabularx}{\linewidth}{|X|X|X|}
\hline
\rowcolor{poporangeL} \textbf{Collocation (chinois + pinyin)} & \textbf{Structure} & \textbf{Traduction} \\
\hline
\zh{过春节} (guò Chūnjié) & V + N (fête) & Célébrer le Nouvel An chinois \\
\hline
\zh{过生日} (guò shēngrì) & V + N (fête) & Fêter son anniversaire \\
\hline
\zh{全家一起} (quánjiā yìqǐ) & Sujet + Adv & Toute la famille ensemble \\
\hline
\zh{传统节日} (chuántǒng jiérì) & Adj + N & Fête traditionnelle \\
\hline
\zh{传统服装} (chuántǒng fúzhuāng) & Adj + N & Vêtement traditionnel / costume traditionnel \\
\hline
\zh{唱歌跳舞} (chànggē tiàowǔ) & V + V (série) & Chanter et danser \\
\hline
\zh{穿服装} (chuān fúzhuāng) & V + N & Porter un costume / s'habiller \\
\hline
\zh{音乐和舞蹈} (yīnyuè hé wǔdǎo) & N + 和 + N & La musique et la danse \\
\hline
\zh{让我们了解} (ràng wǒmen liǎojiě) & 让 + Sujet + V & Nous permettre de comprendre \\
\hline
\zh{互相了解} (hùxiāng liǎojiě) & Adv + V & Se comprendre mutuellement \\
\hline
\end{tabularx}
\end{center}

% ============================================================
% SECTION 4 : CLÉS (RADICAUX) ET ORDRE DES TRAITS
% ============================================================
\coursec{Clés (radicaux) et ordre des traits}

\courssub{Rappel méthodologique}

\begin{methode}[Comment apprendre un caractère ?]
\begin{enumerate}
\item \textbf{Identifier la clé (radical)} : elle donne souvent un indice sémantique (le sens) ou phonétique.
\item \textbf{Compter les traits restants} (hors clé) pour chercher dans le dictionnaire.
\item \textbf{Respecter l'ordre des traits} : horizontal avant vertical, gauche avant droite, haut avant bas, extérieur avant intérieur, trait central avant les ailes, point final.
\item \textbf{Écrire en répétant} à voix haute le nom des traits ou le pinyin.
\end{enumerate}
\end{methode}

\courssub{Huit caractères clés de la leçon}

\begin{center}
\begin{tabularx}{\linewidth}{|c|X|X|X|X|}
\hline
\rowcolor{poppurpleL} \textbf{Car.} & \textbf{Pinyin / Sens} & \textbf{Clé (Radical)} & \textbf{Nb traits total} & \textbf{Ordre des traits (description)} \\
\hline
文 & wén (culture, écriture) & \zh{文} (wén, radical n°67) & 4 & 1. Point (丶) 2. Horizontal (一) 3. Crochet vertical (丿) 4. Point (丶) \\
\hline
化 & huà (changer, -isation) & \zh{亻} (rén, homme debout) & 4 & 1. Homme (亻: traits 1-2) 2. Haut de 匕 (horizontal) 3. Crochet vertical de 匕 4. Trait final descendant droit \\
\hline
春 & chūn (printemps) & \zh{日} (rì, soleil) & 9 & 1. 三 (trois horizontales) 2. 人 (homme: traits 4-5) 3. 日 (soleil: traits 6-9 : vertical, pli, horizontales) \\
\hline
节 & jié (fête, nœud, articuler) & \zh{竹} (zhú, bambou) & 13 & 1-6. Clé 竹 (deux fois : haut en haut, bas en bas) 7. Horizontal 8. Vertical 9. Horizontal 10. Pli horizontal-crochet 11. Horizontal 12. Vertical 13. Horizontal \\
\hline
传 & chuán (transmettre, tradition) & \zh{亻} (rén, homme) & 6 & 1-2. Homme (亻) 3. Point 4. Horizontal 5. Crochet vertical 6. Point (clé phonétique 云 simplifiée) \\
\hline
统 & tǒng (unifier, tout) & \zh{纟} (sī, fil de soie) & 7 & 1-2. Fil (纟: diagonale, point) 3. Horizontal 4. Vertical 5. Horizontal 6. Crochet vertical 7. Horizontal (partie droite 充) \\
\hline
乐 & yuè (musique) / lè (joyeux) & \zh{丿} (piě, trait oblique) & 5 & 1. 丿 2. 丨 (vertical) 3. 一 (horizontal) 4. 丿 (oblique) 5. 乀 (trait final descendant droit) \\
\hline
舞 & wǔ (danser) & \zh{夕} (xī, soir) & 14 & 1-3. 夕 (haut: horizontales, vertical) 4-13. Partie basse 无 (sans: horizontales, vertical, crochets, point) 14. Trait final descendant (si forme traditionnelle 無, ici simplifié 无 + 夕) \textbf{Note :} \zh{舞} s'écrit 夕 au-dessus de 无. Ordre: écrire 夕 (3 traits), puis 无 (11 traits). \\
\hline
\end{tabularx}
\end{center}

\begin{attention}[Pièges d'écriture pour cette leçon]
\begin{itemize}
\item \zh{化} : ne pas confondre avec \zh{北} (běi, nord). \zh{化} a un trait final descendant droit (乀), \zh{北} se termine par une horizontale.
\item \zh{节} : la clé \zh{竹} s'écrit en deux fois (le haut puis le bas), ne pas faire les six traits d'un seul tenant sans lever le stylo.
\item \zh{舞} : attention à l'ordre : le haut \zh{夕} (3 traits) \textbf{avant} le bas \zh{无} (11 traits). Ne pas écrire le trait final du 无 avant d'avoir fini le haut.
\item \zh{传} vs \zh{专} (zhuān) : \zh{传} a la clé \zh{亻} (homme) à gauche ; \zh{专} n'en a pas.
\end{itemize}
\end{attention}

\begin{experience}[Atelier écriture : « Course aux radicaux »]
\begin{enumerate}
\item Le professeur affiche un radical (ex: \zh{亻}, \zh{纟}, \zh{日}, \zh{竹}).
\item Par groupes de 3, les élèves listent en 2 minutes le maximum de caractères de la leçon (ou connus) contenant ce radical.
\item Mise en commun au tableau : vérification de l'ordre des traits pour les caractères de la leçon.
\end{enumerate}
\end{experience}

% ============================================================
% SECTION 5 : COMPRÉHENSION DU TEXTE ET COLLOCATIONS
% ============================================================
\coursec{Compréhension du texte et collocations}

\courssub{Questions de compréhension globale et détaillée}

\begin{exercice}[Questions sur le texte]
Répondez en français (ou en chinois pour les questions 5 et 6) d'après le texte \zh{中国和喀麦隆的文化生活}.
\begin{enumerate}
\item Quels sont les deux pays comparés dans le texte ?
\item Citez trois activités que font les Chinois pendant la fête du Printemps (phrase 2).
\item Comment les Camerounais s'habillent-ils pendant les fêtes traditionnelles ? (phrase 3)
\item Quel est le point commun entre les peuples chinois et camerounais ? (phrase 4)
\item Traduisez en chinois : « La culture nous permet de nous comprendre mutuellement. »
\item Complétez en caractères : \zh{中国人 \_\_\_ 春节，全家 \_\_\_ 吃饭。} (Utilisez \zh{过} et \zh{一起}).
\end{enumerate}
\end{exercice}

\begin{corrige}[Exercice : Questions sur le texte]
\begin{enumerate}
\item La Chine (\zh{中国}) et le Cameroun (\zh{喀麦隆}).
\item Manger ensemble (\zh{吃饭}), regarder la danse du dragon et du lion (\zh{看舞龙舞狮}), célébrer la fête (\zh{过春节}).
\item Ils portent de beaux vêtements traditionnels (\zh{穿漂亮的传统服装}).
\item Ils aiment tous beaucoup la musique et la danse (\zh{都很喜欢音乐和舞蹈}).
\item \zh{文化让我们互相了解。} (Wénhuà ràng wǒmen hùxiāng liǎojiě.)
\item \zh{中国人 \textbf{过} 春节，全家 \textbf{一起} 吃饭。}
\end{enumerate}
\end{corrige}

\courssub{Exercices sur les collocations et la structure}

\begin{exercice}[Mots-outils et order des mots]
Remettez les mots dans l'ordre pour former une phrase correcte. Écrivez en caractères et en pinyin.
\begin{enumerate}
\item \zh{都 / 中国 / 喀麦隆 / 有 / 丰富 / 的 / 文化生活 / 和}
\item \zh{人们 / 传统节日 / 唱歌 / 跳舞 / 在 / 穿 / 服装 / 漂亮 / 的} (Ajoutez \zh{在} pour « pendant » si nécessaire, ou structure simple).
\item \zh{让 / 文化 / 我们 / 互相 / 了解}
\item \zh{两个 / 国家 / 人民 / 音乐 / 和 / 舞蹈 / 都 / 喜欢 / 很}
\end{enumerate}
\end{exercice}

\begin{corrige}[Exercice : Mots-outils et ordre des mots]
\begin{enumerate}
\item \zh{中国和喀麦隆都有丰富的文化生活。} (Zhōngguó hé Kāmàilóng dōu yǒu fēngfù de wénhuà shēnghuó.)
\item \zh{人们在传统节日唱歌、跳舞、穿漂亮的传统服装。} (Rénmen zài chuántǒng jiérì chànggē, tiàowǔ, chuān piàoliang de chuántǒng fúzhuāng.) \textit{Note : \zh{在}... indique le moment. Sans \zh{在} : \zh{人们传统节日唱歌...} (sujet + temps + verbe) est aussi correct.}
\item \zh{文化让我们互相了解。} (Wénhuà ràng wǒmen hùxiāng liǎojiě.)
\item \zh{两个国家的人民都很喜欢音乐和舞蹈。} (Liǎng ge guójiā de rénmín dōu hěn xǐhuan yīnyuè hé wǔdǎo.)
\end{enumerate}
\end{corrige}

\begin{exercice}[Remplacement lexical (Substitution)]
Remplacez l'élément souligné dans la phrase modèle par les mots proposés. Réécrivez la phrase complète.
\emph{Phrase modèle :} \zh{中国人过春节，全家一起吃饭。}
\begin{enumerate}
\item Remplacer \zh{中国人} par \zh{喀麦隆人} ; \zh{春节} par \zh{传统节日} ; \zh{吃饭} par \zh{唱歌、跳舞}。
\item Remplacer \zh{全家} par \zh{我们} ; \zh{一起} par \zh{都} (attention à la place !).
\end{enumerate}
\end{exercice}

\begin{corrige}[Exercice : Remplacement lexical]
\begin{enumerate}
\item \zh{喀麦隆人过传统节日，全家一起唱歌、跳舞。} (Kāmàilóngrén guò chuántǒng jiérì, quánjiā yìqǐ chànggē, tiàowǔ.)
\item \zh{我们都过春节，一起吃饭。} (Wǒmen dōu guò Chūnjié, yìqǐ chīfàn.) \textit{Note : \zh{都} se place avant \zh{过}. \zh{一起} reste avant \zh{吃}.}
\end{enumerate}
\end{corrige}

% ============================================================
% SECTION 6 : PRODUCTION GUIDÉE ET OUVERTURE INTERCULTURELLE
% ============================================================
\coursec{Production guidée et ouverture interculturelle}

\courssub{Expression écrite : Ma fête préférée}

\begin{methode}[Rédiger un court paragraphe descriptif (niveau HSK 2-3)]
\begin{enumerate}
\item \textbf{Présenter la fête} : Nom + temps (\zh{是}... \zh{的} / \zh{在}...).
\item \textbf{Décrire les activités} : Sujet + \zh{一起} / \zh{都} + Verbes (\zh{吃}, \zh{看}, \zh{唱}, \zh{跳}, \zh{穿}...).
\item \textbf{Exprimer le sentiment} : \zh{很有意思} (hěn yǒu yìsi), \zh{很开心} (hěn kāixīn), \zh{让我们互相了解}。
\item \textbf{Relire} : vérifier l'ordre \textbf{Sujet - Temps - Lieu - Adverbe - Verbe - Complément}, les particules \zh{的/得/地}, les classificateurs.
\end{enumerate}
\end{methode}

\begin{exercice}[Production écrite guidée]
Rédigez un paragraphe de 5 à 6 phrases en caractères (avec pinyin si besoin) sur le thème : \textbf{« Une fête traditionnelle au Cameroun »} (ex: Ngondo, Fête du Maïs, Medumba, Fête nationale du 20 mai, Tabaski, Noël).
Utilisez au moins 8 mots du glossaire de cette leçon. Suivez le modèle ci-dessous.

\emph{Modèle (Fête du Printemps en Chine) :}
\zh{春节是中国的传统节日。人们过春节，全家一起吃饺子、看舞龙舞狮。大家都穿新衣服，很开心。文化让我们互相了解。}
\end{exercice}

\begin{corrige}[Exercice : Production écrite guidée (Exemple : Ngondo)]
\shortstack[l]{\zh{恩贡多是喀麦隆的传统节日。}\\ \zh{Ēngòngduō shì Kāmàilóng de chuántǒng jiérì.}\\ \zh{人们过恩贡多，一起唱歌、跳舞、看赛艇。}\\ \zh{Rénmen guò Ēngòngduō, yìqǐ chànggē, tiàowǔ, kàn sàitǐng.}\\ \zh{大家穿漂亮的传统服装。}\\ \zh{Dàjiā chuān piàoliang de chuántǒng fúzhuāng.}\\ \zh{音乐和舞蹈很有意思。}\\ \zh{Yīnyuè hé wǔdǎo hěn yǒu yìsi.}\\ \zh{文化让我们互相了解。}\\ \zh{Wénhuà ràng wǒmen hùxiāng liǎojiě.}}

\textit{Vocabulaire complémentaire utilisé : \zh{是} (shì, être), \zh{的} (de, possession/attribut), \zh{大家} (dàjiā, tout le monde), \zh{新衣服} (xīn yīfu, nouveaux vêtements), \zh{有意思} (yǒu yìsi, intéressant), \zh{赛艇} (sàitǐng, course de pirogues/aviron).}
\end{corrige}

\courssub{Expression orale : Échange interculturel (Tandem)}

\begin{experience}[Jeu de rôle : « Invitation à la fête »]
\textbf{Contexte :} Un élève chinois (A) et un élève camerounais (B) discutent sur WeChat / WhatsApp.
\textbf{Consignes :}
\begin{enumerate}
\item A demande à B : « Quelle est ta fête préférée ? » (\zh{你最喜欢哪个节日？} Nǐ zuì xǐhuan nǎ ge jiérì?)
\item B répond : « C'est le Ngondo / la Fête du Maïs / Noël... » (\zh{是...} Shì...)
\item A demande : « Qu'est-ce qu'on fait ? » (\zh{都做什么？} Dōu zuò shénme?)
\item B décrit 3 activités (manger, danser, porter des habits, chanter...).
\item A compare : « En Chine, pour le Printemps, on fait... » (\zh{中国过春节...} Zhōngguó guò Chūnjié...)
\item Conclusion commune : « La culture nous rapproche. » (\zh{文化让我们更近。} Wénhuà ràng wǒmen gèng jìn.)
\end{enumerate}
\textbf{Évaluation par les pairs :} Grille sur 5 pts (Prononciation des tons / Vocabulaire de la leçon / Structures \zh{过}, \zh{一起}, \zh{都}, \zh{让} / Fluidité).
\end{experience}

\courssub{Élément socioculturel : Le Ngondo et le Printemps chinois}

\begin{savaistu}[Deux fêtes, une même âme : le Ngondo (Cameroun) et Chūnjié (Chine)]
\begin{itemize}
\item \textbf{Le Ngondo (Cameroun, peuple Sawa)} : Grande fête annuelle (novembre/décembre) sur les rives du Wouri à Douala. Cérémonie du \emph{Jebalè} (plongeon du vase sacré), courses de pirogues (\zh{赛艇} sàitǐng), chants \emph{Essèwé}, danses, tenues traditionnelles (\zh{传统服装} chuántǒng fúzhuāng). Moment de communion avec les ancêtres (\zh{祖先} zǔxiān) et de cohésion sociale.
\item \textbf{La Fête du Printemps \zh{春节} (Chine)} : Premier jour du calendrier lunaire (janvier/février). Réveillon familial (\zh{年夜饭} niányèfàn), raviolis (\zh{饺子} jiǎozi), enveloppes rouges (\zh{红包} hóngbāo), feux d'artifice, danse du dragon/lion (\zh{舞龙舞狮}), vœux (\zh{拜年} bài nián). Moment de retour aux sources (\zh{回家} huíjiā) et de renouveau.
\item \textbf{Points communs structurels} : 1. \textbf{Rituel de passage} (fin d'un cycle, début d'un autre). 2. \textbf{Rassemblement familial/clanique} (\zh{全家} / \zh{家族}). 3. \textbf{Performance artistique collective} (musique, danse, spectacle). 4. \textbf{Transmission intergénérationnelle} (les aînés transmettent aux jeunes).
\item \textbf{Différence notable} : Le Ngondo est fortement lié à l'eau (fleuve Wouri) et au culte des ancêtres aquatiques (Jengu) ; Chūnjié est lié au cycle agraire lunaire, au feu (pétards) et aux divinités domestiques (Dieu du Foyer).
\end{itemize}
\end{savaistu}

\begin{aretenir}[Bilan de la leçon 19 -- Vocabulaire \& Compréhension]
\begin{itemize}
\item \textbf{Lexique maîtrisé} : 30+ mots (noms de fêtes, verbes d'action, vêtements, arts, structure \zh{让}, adverbes \zh{都/一起/互相}).
\item \textbf{Grammaire ancrée} : \zh{过}+fête, \zh{两} vs \zh{二}, place de \zh{都/一起}, causatif \zh{让}, \zh{互相}+V.
\item \textbf{Compétence lecture} : Comprendre un texte comparatif court, identifier la structure Sujet-Verbe-Complément et les énumérations (\zh{、}).
\item \textbf{Compétence écriture} : Écrire 8 caractères clés (ordre traits, radicaux), produire un paragraphe descriptif simple.
\item \textbf{Ouverture culturelle} : Connaissance factuelle du Ngondo et de Chūnjié, capacité à établir un parallèle respectueux Cameroun-Chine.
\end{itemize}
\end{aretenir}

\coursec{Glossaire du thème}

Dans la section précédente, nous avons lu et traduit le texte support sur la vie culturelle au Cameroun et en Chine. Vous avez remarqué que certains mots revenaient souvent : \zh{文化} « culture », \zh{节日} « fête », \zh{表演} « représentation »... C'est exactement le signe qu'il s'agit du \cle{vocabulaire central} du thème. Dans cette section, nous allons étudier ces mots un par un : leur prononciation exacte (pinyin et tons), leur nature grammaticale, leur emploi dans une phrase, et les pièges qu'ils tendent aux francophones. Un bon glossaire n'est pas une liste à réciter : c'est une boîte à outils que vous devez savoir ouvrir au bon moment.

\courssub{Le tableau du glossaire}

Voici les dix-sept mots essentiels de la leçon. Lisez chaque ligne à voix haute en suivant le pinyin, puis cachez la colonne des traductions et testez-vous.

\begin{center}
\begin{tabularx}{\linewidth}{|X|X|X|X|}
\hline
\rowcolor{popblueL} Caractères & Pinyin & Nature & Traduction \\
\hline
文化 & wénhuà & n. & culture \\
\hline
节日 & jiérì & n. & fête, jour de fête \\
\hline
春节 & Chūnjié & n. propre & fête du Printemps (Nouvel An chinois) \\
\hline
传统 & chuántǒng & adj. / n. & traditionnel / tradition \\
\hline
服装 & fúzhuāng & n. & vêtement, costume \\
\hline
音乐 & yīnyuè & n. & musique \\
\hline
舞蹈 & wǔdǎo & n. & danse (l'art) \\
\hline
艺术 & yìshù & n. & art \\
\hline
唱歌 & chànggē & v. & chanter \\
\hline
跳舞 & tiàowǔ & v. & danser \\
\hline
表演 & biǎoyǎn & v. / n. & jouer, présenter / représentation, spectacle \\
\hline
庆祝 & qìngzhù & v. & célébrer, fêter \\
\hline
丰富 & fēngfù & adj. & riche, abondant \\
\hline
漂亮 & piàoliang & adj. & beau, joli \\
\hline
互相 & hùxiāng & adv. & mutuellement, l'un l'autre \\
\hline
了解 & liǎojiě & v. & comprendre, connaître \\
\hline
人民 & rénmín & n. & peuple \\
\hline
\end{tabularx}
\end{center}

\begin{popfigure}
\begin{tikzpicture}
\node[font=\bfseries\small, fill=popblue, text=white, minimum width=2.6cm, minimum height=0.7cm] (t1) at (0,0) {Caractères};
\node[font=\bfseries\small, fill=popblue, text=white, minimum width=2.4cm, minimum height=0.7cm] (t2) at (2.7,0) {Pinyin};
\node[font=\bfseries\small, fill=popblue, text=white, minimum width=2.2cm, minimum height=0.7cm] (t3) at (5.1,0) {Nature};
\node[font=\bfseries\small, fill=popblue, text=white, minimum width=3.4cm, minimum height=0.7cm] (t4) at (8.1,0) {Traduction};
\node[font=\small, fill=popblueL, minimum width=2.6cm, minimum height=0.65cm] at (0,-0.75) {文化};
\node[font=\small\itshape, fill=popblueL, minimum width=2.4cm, minimum height=0.65cm] at (2.7,-0.75) {wénhuà};
\node[font=\small, fill=popblueL, minimum width=2.2cm, minimum height=0.65cm] at (5.1,-0.75) {n.};
\node[font=\small, fill=popblueL, minimum width=3.4cm, minimum height=0.65cm] at (8.1,-0.75) {culture};
\node[font=\small, fill=popcream, minimum width=2.6cm, minimum height=0.65cm] at (0,-1.45) {节日};
\node[font=\small\itshape, fill=popcream, minimum width=2.4cm, minimum height=0.65cm] at (2.7,-1.45) {jiérì};
\node[font=\small, fill=popcream, minimum width=2.2cm, minimum height=0.65cm] at (5.1,-1.45) {n.};
\node[font=\small, fill=popcream, minimum width=3.4cm, minimum height=0.65cm] at (8.1,-1.45) {fête};
\node[font=\small, fill=popblueL, minimum width=2.6cm, minimum height=0.65cm] at (0,-2.15) {传统};
\node[font=\small\itshape, fill=popblueL, minimum width=2.4cm, minimum height=0.65cm] at (2.7,-2.15) {chuántǒng};
\node[font=\small, fill=popblueL, minimum width=2.2cm, minimum height=0.65cm] at (5.1,-2.15) {adj./n.};
\node[font=\small, fill=popblueL, minimum width=3.4cm, minimum height=0.65cm] at (8.1,-2.15) {traditionnel / tradition};
\node[font=\small, fill=popcream, minimum width=2.6cm, minimum height=0.65cm] at (0,-2.85) {音乐};
\node[font=\small\itshape, fill=popcream, minimum width=2.4cm, minimum height=0.65cm] at (2.7,-2.85) {yīnyuè};
\node[font=\small, fill=popcream, minimum width=2.2cm, minimum height=0.65cm] at (5.1,-2.85) {n.};
\node[font=\small, fill=popcream, minimum width=3.4cm, minimum height=0.65cm] at (8.1,-2.85) {musique};
\node[font=\small, fill=popblueL, minimum width=2.6cm, minimum height=0.65cm] at (0,-3.55) {跳舞};
\node[font=\small\itshape, fill=popblueL, minimum width=2.4cm, minimum height=0.65cm] at (2.7,-3.55) {tiàowǔ};
\node[font=\small, fill=popblueL, minimum width=2.2cm, minimum height=0.65cm] at (5.1,-3.55) {v.};
\node[font=\small, fill=popblueL, minimum width=3.4cm, minimum height=0.65cm] at (8.1,-3.55) {danser};
\node[font=\small, fill=popcream, minimum width=2.6cm, minimum height=0.65cm] at (0,-4.25) {庆祝};
\node[font=\small\itshape, fill=popcream, minimum width=2.4cm, minimum height=0.65cm] at (2.7,-4.25) {qìngzhù};
\node[font=\small, fill=popcream, minimum width=2.2cm, minimum height=0.65cm] at (5.1,-4.25) {v.};
\node[font=\small, fill=popcream, minimum width=3.4cm, minimum height=0.65cm] at (8.1,-4.25) {célébrer};
\node[font=\small, fill=popblueL, minimum width=2.6cm, minimum height=0.65cm] at (0,-4.95) {丰富};
\node[font=\small\itshape, fill=popblueL, minimum width=2.4cm, minimum height=0.65cm] at (2.7,-4.95) {fēngfù};
\node[font=\small, fill=popblueL, minimum width=2.2cm, minimum height=0.65cm] at (5.1,-4.95) {adj.};
\node[font=\small, fill=popblueL, minimum width=3.4cm, minimum height=0.65cm] at (8.1,-4.95) {riche};
\node[font=\small, fill=popcream, minimum width=2.6cm, minimum height=0.65cm] at (0,-5.65) {互相};
\node[font=\small\itshape, fill=popcream, minimum width=2.4cm, minimum height=0.65cm] at (2.7,-5.65) {hùxiāng};
\node[font=\small, fill=popcream, minimum width=2.2cm, minimum height=0.65cm] at (5.1,-5.65) {adv.};
\node[font=\small, fill=popcream, minimum width=3.4cm, minimum height=0.65cm] at (8.1,-5.65) {mutuellement};
\end{tikzpicture}
\legende{Extrait visuel du glossaire : huit mots représentatifs des quatre familles (noms, verbes, adjectifs, adverbe).}
\end{popfigure}

\courssub{Étude mot par mot : observations et remarques}

\courssub{La famille des fêtes : 文化, 节日, 春节, 传统, 庆祝}

\begin{definition}[传统 chuántǒng]
\zh{传统} désigne ce qui est \cle{transmis de génération en génération} : les coutumes, les fêtes, les arts, les façons de vivre héritées des anciens. Comme nom, il signifie « tradition » ; comme adjectif, « traditionnel ». Exemple : \zh{传统服装} \emph{chuántǒng fúzhuāng} « costume traditionnel » ; \zh{中国的传统} \emph{Zhōngguó de chuántǒng} « les traditions de la Chine ».
\end{definition}

Observez la construction des mots : \zh{节日} \emph{jiérì} est composé de \zh{节} « fête » et \zh{日} « jour » : littéralement « jour de fête ». \zh{春节} \emph{Chūnjié} combine \zh{春} « printemps » et \zh{节} « fête » : la « fête du Printemps », c'est-à-dire le Nouvel An chinois, la fête la plus importante de Chine. Comme c'est un nom propre, on écrit \emph{Chūnjié} avec une majuscule en pinyin.

Le verbe \zh{庆祝} \emph{qìngzhù} « célébrer » prend directement un complément d'objet, comme en français : \zh{庆祝节日} « célébrer la fête », \zh{庆祝新年} « fêter le Nouvel An ».

\courssub{La famille des arts : 音乐, 舞蹈, 艺术, 服装, 表演}

Attention à la prononciation de \zh{音乐} \emph{yīnyuè} : le caractère \zh{乐} se lit \emph{yuè} quand il signifie « musique », mais \emph{lè} quand il signifie « joyeux » (comme dans \zh{快乐} \emph{kuàilè} « heureux »). C'est un caractère à plusieurs lectures, fréquent en chinois.

\zh{表演} \emph{biǎoyǎn} fonctionne à la fois comme verbe (« jouer, présenter un spectacle ») et comme nom (« représentation, spectacle ») : \zh{他们表演舞蹈。} \emph{Tāmen biǎoyǎn wǔdǎo.} « Ils présentent une danse. » / \zh{表演很精彩。} \emph{Biǎoyǎn hěn jīngcǎi.} « Le spectacle est magnifique. »

\courssub{Les verbes à complément incorporé : 唱歌 et 跳舞}

\begin{attention}[唱歌 et 跳舞 : verbe + objet figé]
\zh{唱歌} \emph{chànggē} et \zh{跳舞} \emph{tiàowǔ} sont des \cle{verbes à complément incorporé} : le mot contient déjà un verbe et son objet figé. \zh{唱} « chanter » + \zh{歌} « chanson » = « chanter-une-chanson » ; \zh{跳} « sauter » + \zh{舞} « danse » = « danser-une-danse ». Conséquence pratique : pour préciser \emph{ce que} l'on chante ou danse, on \cle{sépare les deux éléments} et on insère le complément au milieu : \zh{唱一首中文歌} \emph{chàng yì shǒu Zhōngwén gē} « chanter une chanson chinoise » (avec le classificateur \zh{首} pour les chansons). De même : \zh{跳一个舞} « danser une danse ». Erreur typique du francophone : coller un second objet après le mot entier, comme *\zh{唱歌一首歌} (« chanter-une-chanson une chanson »), ce qui est impossible : le complément doit passer \emph{entre} le verbe et l'objet incorporé.
\end{attention}

\courssub{Adjectifs, adverbe et mots de lien : 丰富, 漂亮, 互相, 了解, 人民}

Les adjectifs \zh{丰富} \emph{fēngfù} « riche, abondant » et \zh{漂亮} \emph{piàoliang} « beau, joli » s'emploient comme prédicats avec \zh{很} \emph{hěn}, sans verbe « être » : \zh{文化很丰富。} « La culture est riche. » Notez que \zh{丰富} qualifie l'abondance (culture, contenu, vie), tandis que \zh{漂亮} qualifie l'apparence (vêtements, personnes, objets).

L'adverbe \zh{互相} \emph{hùxiāng} « mutuellement » se place \cle{avant le verbe} : \zh{互相了解} \emph{hùxiāng liǎojiě} « se comprendre mutuellement, apprendre à se connaître ». En français, « se ... mutuellement » est exprimé par le pronom réfléchi ; en chinois, c'est l'adverbe \zh{互相} qui porte cette idée.

\begin{exemplebox}[Exemples traduits]
\zh{我们庆祝节日。} \emph{Wǒmen qìngzhù jiérì.} « Nous célébrons la fête. » \\
\zh{喀麦隆的文化很丰富。} \emph{Kāmàilóng de wénhuà hěn fēngfù.} « La culture camerounaise est très riche. » \\
\zh{春节的时候，中国人穿传统服装。} \emph{Chūnjié de shíhou, Zhōngguórén chuān chuántǒng fúzhuāng.} « Au moment de la fête du Printemps, les Chinois portent des costumes traditionnels. » \\
\zh{她喜欢唱歌，我喜欢跳舞。} \emph{Tā xǐhuan chànggē, wǒ xǐhuan tiàowǔ.} « Elle aime chanter, moi j'aime danser. » \\
\zh{两国人民互相了解。} \emph{Liǎng guó rénmín hùxiāng liǎojiě.} « Les peuples des deux pays apprennent à se connaître mutuellement. » \\
\zh{这个表演很漂亮。} \emph{Zhège biǎoyǎn hěn piàoliang.} « Ce spectacle est très beau. »
\end{exemplebox}

\courssub{Comparer le chinois et le français}

\begin{center}
\begin{tabularx}{\linewidth}{|X|X|X|}
\hline
\rowcolor{popblueL} Notion & En chinois & En français \\
\hline
Adjectif prédicat & 文化很丰富。(sans verbe « être ») & « La culture \emph{est} riche. » (verbe obligatoire) \\
\hline
« Chanter une chanson » & 唱一首中文歌 (complément inséré au milieu) & un seul verbe + objet à la suite \\
\hline
« Se comprendre » & 互相了解 (adverbe avant le verbe) & pronom réfléchi « se » \\
\hline
Nom propre de fête & 春节 (majuscule : Chūnjié) & « la fête du Printemps » \\
\hline
\end{tabularx}
\end{center}

\courssub{Méthode de mémorisation}

\begin{methode}[Mémoriser le glossaire durablement]
\begin{enumerate}
\item \textbf{Regroupez par familles sémantiques} : les fêtes (\zh{节日}, \zh{春节}, \zh{庆祝}), les arts (\zh{音乐}, \zh{舞蹈}, \zh{艺术}, \zh{表演}), les actions (\zh{唱歌}, \zh{跳舞}), les qualités (\zh{丰富}, \zh{漂亮}), les relations (\zh{互相}, \zh{了解}, \zh{人民}).
\item \textbf{Associez chaque mot à une image mentale} : pour \zh{春节}, imaginez les lanternes rouges ; pour \zh{跳舞}, les danseurs du ballet national du Cameroun ; pour \zh{丰富}, une table de ndolé bien garnie.
\item \textbf{Lisez à voix haute} en marquant les tons avec la main (la main monte pour le deuxième ton, descend pour le quatrième).
\item \textbf{Testez-vous dans les deux sens} : caractères vers traduction, puis traduction vers caractères et pinyin.
\item \textbf{Réutilisez chaque mot dans une phrase personnelle} : un mot employé est un mot mémorisé.
\end{enumerate}
\end{methode}

\begin{exoresolu}[Quiz pinyin vers traduction]
Énoncé. Traduisez en français et écrivez les caractères : 1. \emph{jiérì} ; 2. \emph{yīnyuè} ; 3. \emph{qìngzhù} ; 4. \emph{hùxiāng} ; 5. \emph{fúzhuāng}.
\tcblower
Solution détaillée. 1. \emph{jiérì} = \zh{节日} « fête » (nom) ; 2. \emph{yīnyuè} = \zh{音乐} « musique » (attention : \zh{乐} se lit \emph{yuè} ici, pas \emph{lè}) ; 3. \emph{qìngzhù} = \zh{庆祝} « célébrer » (verbe, suivi directement du complément : \zh{庆祝节日}) ; 4. \emph{hùxiāng} = \zh{互相} « mutuellement » (adverbe placé avant le verbe : \zh{互相了解}) ; 5. \emph{fúzhuāng} = \zh{服装} « vêtement, costume » (nom ; on dit \zh{传统服装} « costume traditionnel »).
\end{exoresolu}

\begin{experience}[Cache-cache des caractères et quiz éclair]
En binômes. Étape 1 : un élève lit le tableau du glossaire pendant une minute, puis le ferme ; son camarade écrit les traductions françaises de dix mots et l'élève doit retrouver caractères et pinyin. Étape 2 : quiz pinyin vers traduction — le camarade dicte cinq pinyin au hasard (\emph{wǔdǎo}, \emph{piàoliang}, \emph{liǎojiě}...), l'élève donne la traduction et une phrase d'exemple. Étape 3 : inversez les rôles. Le binôme qui obtient dix bonnes réponses complètes (caractères + pinyin + traduction + phrase) gagne la partie.
\end{experience}

\begin{aretenir}[L'essentiel du glossaire]
Dix-sept mots, quatre familles : fêtes, arts, actions, relations. Trois points de vigilance : \zh{乐} se lit \emph{yuè} dans \zh{音乐} ; \zh{唱歌} et \zh{跳舞} sont des verbes à complément incorporé (\zh{唱一首中文歌}) ; \zh{互相} se place avant le verbe. Avec ces mots, vous pouvez déjà décrire toute la vie culturelle du Cameroun et de la Chine.
\end{aretenir}

\coursec{Clés (radicaux) et ordre des traits}

Après avoir appris le vocabulaire de la vie culturelle dans le glossaire, nous allons maintenant entrer \emph{à l'intérieur} des caractères. En chinois, écrire n'est pas un geste au hasard : chaque caractère obéit à une logique de construction (la \cle{clé} ou \cle{radical}, qui donne une indication de sens) et à un \cle{ordre des traits} précis. Comprendre cette logique, c'est apprendre deux fois plus vite : au lieu de mémoriser un dessin compliqué, on mémorise des pièces simples que l'on assemble.

\courssub{Rappel : le rôle des clés dans le sens des caractères}

\begin{definition}[Qu'est-ce qu'une clé (radical) ?]
Une \cle{clé} (en chinois \zh{部首}, \emph{bùshǒu}) est la partie d'un caractère qui indique sa \textbf{famille de sens}. C'est aussi l'élément sous lequel le caractère est classé dans les dictionnaires. L'autre partie du caractère donne souvent une indication de \textbf{prononciation} : on l'appelle le \cle{composant phonétique}.
\end{definition}

Observez les caractères de notre thème :

\begin{exemplebox}[La clé donne le sens, l'autre partie donne souvent le son]
\zh{祝} (zhù, souhaiter) contient la clé \zh{礻} (autel, rites) : souhaiter, c'est un geste rituel.\\
\zh{节} (jié, fête) contient la clé \zh{艹} (herbe) : à l'origine, le caractère désignait les nœuds du bambou, une plante.\\
\zh{歌} (gē, chanson) contient la clé \zh{欠} (bâiller, ouvrir la bouche) : chanter, c'est ouvrir la bouche ! Le composant \zh{哥} (gē) donne la prononciation.
\end{exemplebox}

\begin{propriete}[La clé 礻 (shì), l'autel]
La clé \zh{礻} (shì) représente un \textbf{autel} et apparaît dans les caractères liés aux \textbf{rites, aux célébrations et au sacré} :
\begin{itemize}
\item \zh{祝} (zhù) : souhaiter, célébrer ;
\item \zh{礼} (lǐ) : le rite, la politesse, le cadeau (\zh{礼物} lǐwù, cadeau) ;
\item \zh{神} (shén) : l'esprit, le dieu.
\end{itemize}
Quand vous voyez \zh{礻} à gauche d'un caractère, pensez : « cérémonie, fête, sacré ».
\end{propriete}

\begin{attention}[Ne pas confondre 礻 et 衤]
Piège classique des francophones : la clé de l'autel \zh{礻} s'écrit avec \textbf{un seul point} en haut, tandis que la clé du vêtement \zh{衤} (yī) s'écrit avec \textbf{deux points}. Un trait de différence change tout le sens :
\begin{itemize}
\item \zh{祝} (zhù, célébrer) $\neq$ \zh{初} (chū, début) — clé vêtement, car autrefois « couper le tissu » marquait le début du vêtement ;
\item \zh{礼} (lǐ, rite) $\neq$ \zh{补} (bǔ, réparer, clé vêtement).
\end{itemize}
Astuce : l'autel, c'est \textbf{un} point, comme une seule flamme de bougie ; le vêtement a \textbf{deux} points, comme deux boutons.
\end{attention}

\courssub{Méthode générale : l'ordre des traits}

\begin{methode}[Les quatre règles d'or de l'écriture]
\begin{enumerate}
\item \textbf{De haut en bas} : on commence par le haut du caractère (\zh{文} : le point du haut d'abord).
\item \textbf{De gauche à droite} : la partie gauche s'écrit avant la partie droite (\zh{祝} : \zh{礻} d'abord, \zh{兄} ensuite).
\item \textbf{L'horizontal avant le vertical} : quand deux traits se croisent, le trait horizontal \zh{横} (héng) passe avant le trait vertical \zh{竖} (shù).
\item \textbf{La clé souvent en premier} : dans les caractères gauche-droite, on écrit la clé (à gauche) avant le reste.
\end{enumerate}
Respecter l'ordre des traits n'est pas du formalisme : cela rend l'écriture plus rapide, plus belle, et permet d'utiliser les dictionnaires et les claviers à reconnaissance de tracé.
\end{methode}

\courssub{Étude détaillée des caractères du thème}

\textbf{1. 节 (jié), fête — 5 traits.}\par
Clé : \zh{艹} (herbe), en haut. Partie inférieure : \zh{卩}. Ordre des traits :
\begin{enumerate}
\item \zh{横} (héng) : le trait horizontal de la clé herbe ;
\item \zh{竖} (shù) : le premier vertical de l'herbe (à gauche) ;
\item \zh{竖} (shù) : le deuxième vertical de l'herbe (à droite) ;
\item \zh{横折钩} (héng zhé gōu) : horizontal-plié-crochet de \zh{卩} ;
\item \zh{竖} (shù) : le vertical final qui descend.
\end{enumerate}

\textbf{2. 文 (wén), culture, écriture — 4 traits.}\par
Ordre : \zh{点} (diǎn, le point en haut) $\rightarrow$ \zh{横} (héng, l'horizontal) $\rightarrow$ \zh{撇} (piě, le jeté vers la gauche) $\rightarrow$ \zh{捺} (nà, l'appuyé vers la droite). Ce petit caractère est la base de toute notre leçon : \zh{文化} (wénhuà, culture), \zh{中文} (Zhōngwén, le chinois), \zh{文明} (wénmíng, civilisation).

\textbf{3. 祝 (zhù), souhaiter, célébrer — 9 traits.}\par
Clé : \zh{礻} (autel), à gauche, 4 traits : point, horizontal-jeté, vertical, point. Partie droite : \zh{兄} (xiōng, frère aîné), 5 traits : le carré \zh{口} (vertical, horizontal-plié, horizontal de fermeture) puis le jeté et le vertical-courbé. Structure : \zh{祝} = \zh{礻} + \zh{兄}.

\textbf{4. 歌 (gē), chanson — 14 traits.}\par
Clé : \zh{欠} (qiàn, bâiller), à droite, 4 traits. Composant phonétique : \zh{哥} (gē, grand frère), à gauche, 10 traits, formé de deux \zh{可} superposés. Remarquez : \zh{哥} et \zh{歌} se prononcent exactement pareil (gē) — la phonétique est parfaite ici. On écrit d'abord tout le \zh{哥} (de haut en bas), puis \zh{欠}.

\textbf{5. 舞 (wǔ), danse — 14 traits.}\par
Clé : \zh{夕} (xī, soir). C'est le caractère le plus complexe de la leçon ; décomposons-le étape par étape :
\begin{enumerate}
\item Le haut : le jeté \zh{撇} puis l'horizontal \zh{横} ;
\item Les quatre verticaux parallèles \zh{竖} \zh{竖} \zh{竖} \zh{竖} (de gauche à droite !) ;
\item Le long horizontal \zh{横} qui traverse tout ;
\item Le bas gauche : \zh{夕} (jeté, horizontal-jeté, point) ;
\item Le bas droit : jeté, horizontal, vertical final.
\end{enumerate}
Astuce de mémorisation : imaginez quatre danseuses en ligne (les quatre verticaux) sur une scène (le long horizontal), le soir (\zh{夕}).

\begin{popfigure}
\begin{tikzpicture}[font=\small]
\node[draw=popblue, fill=popblueL, rounded corners, minimum width=2.2cm, minimum height=1cm] (zhu) at (0,2.2) {\zh{祝} zhù};
\node[draw=poporange, fill=poporangeL, rounded corners, minimum width=1.6cm, minimum height=0.8cm] (shi) at (-1.8,0.6) {\zh{礻} autel};
\node[draw=popgreen, fill=popgreenL, rounded corners, minimum width=1.6cm, minimum height=0.8cm] (xiong) at (1.8,0.6) {\zh{兄} xiōng};
\draw[-{Stealth}, thick, popblue] (shi) -- (zhu);
\draw[-{Stealth}, thick, popblue] (xiong) -- (zhu);
\node[draw=poppurple, fill=poppurpleL, rounded corners, minimum width=2.2cm, minimum height=1cm] (jie) at (6.5,2.2) {\zh{节} jié};
\node[draw=popgreen, fill=popgreenL, rounded corners, minimum width=1.6cm, minimum height=0.8cm] (cao) at (4.7,0.6) {\zh{艹} herbe};
\node[draw=poporange, fill=poporangeL, rounded corners, minimum width=1.6cm, minimum height=0.8cm] (jie2) at (8.3,0.6) {\zh{卩}};
\draw[-{Stealth}, thick, poppurple] (cao) -- (jie);
\draw[-{Stealth}, thick, poppurple] (jie2) -- (jie);
\node[align=center, popmuted] at (3.2,-0.6) {sens (clé) + partie phonétique ou graphique};
\end{tikzpicture}
\legende{Décomposition des caractères \zh{祝} = \zh{礻} + \zh{兄} et \zh{节} = \zh{艹} + \zh{卩} : la clé porte le sens.}
\end{popfigure}

\begin{center}
\begin{tabularx}{\linewidth}{|X|X|X|X|}\hline
\rowcolor{popblueL} Caractère & Pinyin et sens & Clé & Nombre de traits \\ \hline
\zh{节} & jié, fête & \zh{艹} (herbe) & 5 \\ \hline
\zh{文} & wén, culture & \zh{文} (lui-même) & 4 \\ \hline
\zh{祝} & zhù, célébrer & \zh{礻} (autel) & 9 \\ \hline
\zh{歌} & gē, chanson & \zh{欠} (bâiller) & 14 \\ \hline
\zh{舞} & wǔ, danse & \zh{夕} (soir) & 14 \\ \hline
\end{tabularx}
\end{center}

\begin{textebox}[Les clés dans notre texte culturel]
\zh{文化节到了，我们唱歌、跳舞，祝大家快乐！}\\
\emph{Wénhuà jié dào le, wǒmen chàng gē, tiào wǔ, zhù dàjiā kuàilè !}\\
« La fête de la culture arrive, nous chantons, nous dansons, nous souhaitons à tous beaucoup de joie ! »\\
Retrouvez dans cette phrase les cinq caractères étudiés : \zh{文}, \zh{节}, \zh{歌}, \zh{舞}, \zh{祝}. Pour chacun, identifiez sa clé avant de le recopier.
\end{textebox}

\begin{savaistu}[L'origine du caractère 文]
À l'origine, il y a plus de trois mille ans, \zh{文} (wén) représentait un \textbf{homme vu de face, le buste tatoué} d'un motif décoratif. Le sens premier était donc « motif, dessin ». De « dessin » on est passé à « écriture », puis à tout ce que l'écriture transmet : \zh{文化} (wénhuà, la culture), \zh{中文} (Zhōngwén, le chinois écrit), \zh{文明} (wénmíng, la civilisation). Au Cameroun comme en Chine, les motifs corporels et les dessins traditionnels sont aussi une forme d'écriture de la culture !
\end{savaistu}

\begin{exoresolu}[Identifier les clés]
Énoncé : indiquez la clé de chaque caractère et devinez son domaine de sens : \zh{祝}, \zh{神}, \zh{礼}, \zh{歌}.
\tcblower
Solution détaillée :
\begin{itemize}
\item \zh{祝} : clé \zh{礻} (autel) $\rightarrow$ domaine des rites et célébrations ; zhù, souhaiter.
\item \zh{神} : clé \zh{礻} (autel) $\rightarrow$ domaine du sacré ; shén, esprit, dieu.
\item \zh{礼} : clé \zh{礻} (autel) $\rightarrow$ domaine des rites ; lǐ, rite, politesse.
\item \zh{歌} : clé \zh{欠} (ouvrir la bouche) $\rightarrow$ action de la bouche ; gē, chanson. Le composant \zh{哥} (gē) donne le son.
\end{itemize}
\end{exoresolu}

\begin{exercice}[Écriture guidée]
Sur votre cahier, écrivez \textbf{5 lignes} de chacun des cinq caractères étudiés (\zh{节}, \zh{文}, \zh{祝}, \zh{歌}, \zh{舞}), en respectant l'ordre des traits vu en classe. Pour chaque caractère, écrivez d'abord la clé seule, puis le caractère complet, puis un mot de vocabulaire qui l'emploie (exemple : \zh{节日} jiérì, la fête).
\end{exercice}

\begin{corrige}[Écriture guidée]
Vérifiez votre travail avec cette grille : \zh{节} : 5 traits, herbe en haut d'abord ; \zh{文} : 4 traits, point-horizontal-jeté-appuyé ; \zh{祝} : 9 traits, \zh{礻} à gauche (un seul point !) puis \zh{兄} ; \zh{歌} : 14 traits, \zh{哥} complet avant \zh{欠} ; \zh{舞} : 14 traits, les quatre verticaux de gauche à droite, \zh{夕} en bas à gauche. Mots possibles : \zh{节日} (fête), \zh{文化} (culture), \zh{庆祝} (célébrer), \zh{唱歌} (chanter), \zh{跳舞} (danser).
\end{corrige}

\begin{aretenir}[L'essentiel de la section]
\begin{itemize}
\item La \cle{clé} indique la famille de sens ; l'autre composant indique souvent la prononciation.
\item \zh{礻} (autel, un point) = rites et fêtes : \zh{祝}, \zh{礼}, \zh{神}. Ne pas confondre avec \zh{衤} (vêtement, deux points).
\item Ordre des traits : haut avant bas, gauche avant droite, horizontal avant vertical, clé en premier.
\item \zh{节} 5 traits, \zh{文} 4 traits, \zh{祝} 9 traits, \zh{歌} et \zh{舞} 14 traits.
\end{itemize}
\end{aretenir}

\coursec{Compréhension du texte et collocations}

Après avoir étudié les clés (radicaux) et l'ordre des traits des caractères du thème — comme la clé \zh{人} « personne » dans \zh{众} ou la clé \zh{衣} « vêtement » cachée dans \zh{装} — nous sommes maintenant prêts à vérifier notre compréhension globale du texte. Savoir écrire un caractère, c'est bien ; comprendre le message qu'il porte, c'est l'objectif final de toute lecture. Dans cette section, nous allons apprendre une méthode rigoureuse pour répondre aux questions de compréhension, puis nous enrichirons notre expression avec les collocations clés du thème de la vie culturelle au Cameroun et en Chine.

\courssub{Méthode : répondre à une question de compréhension}

\begin{methode}[Répondre à une question de compréhension en trois étapes]
Pour répondre correctement à une question de compréhension en chinois, suis toujours ces trois étapes :
\begin{enumerate}
\item \textbf{Repérer le mot interrogatif.} Identifie l'outil de question : \zh{什么} (shénme, quoi/quel), \zh{谁} (shéi, qui), \zh{哪儿} (nǎr, où), \zh{怎么} (zěnme, comment), \zh{为什么} (wèishénme, pourquoi). Ce mot te dit \emph{quel type d'information} chercher : une chose, une personne, un lieu, une manière, une cause.
\item \textbf{Retrouver la phrase du texte.} Cherche dans le texte la phrase qui contient les mots-clés de la question (par exemple \zh{节日} fête, \zh{传统服装} vêtements traditionnels). Souligne-la mentalement.
\item \textbf{Répondre par une phrase complète en chinois.} Ne réponds jamais par un seul mot. Reprends la structure de la question en remplaçant le mot interrogatif par l'information trouvée : \zh{中国人过什么节日？} → \zh{中国人过春节。}
\end{enumerate}
\end{methode}

\begin{attention}[La question chinoise se répond « en miroir »]
En chinois, l'ordre des mots de la réponse est le même que celui de la question : on remplace simplement le mot interrogatif par la réponse. Contrairement au français, il n'y a \textbf{pas d'inversion} sujet-verbe et pas de mot interrogatif à déplacer en tête de phrase : \zh{什么} reste à la place du mot qu'il interroge. Erreur fréquente des francophones : écrire \zh{什么中国人过节日？} sous l'influence du français « Quelle fête... ? ». La bonne forme est \zh{中国人过什么节日？} (Zhōngguórén guò shénme jiérì ?).
\end{attention}

\courssub{Questions de compréhension sur le texte}

Appliquons la méthode aux quatre questions du texte support sur la vie culturelle au Cameroun et en Chine.

\begin{exemplebox}[Questions et réponses modèles]
\textbf{Question 1.} \zh{中国人过什么节日？} (Zhōngguórén guò shénme jiérì ? — Quelle fête les Chinois célèbrent-ils ?)\\
Mot interrogatif : \zh{什么} → on cherche une \emph{chose} (une fête).\\
Réponse : \zh{中国人过春节。} (Zhōngguórén guò Chūnjié. — Les Chinois célèbrent la fête du Printemps.)

\textbf{Question 2.} \zh{喀麦隆人在传统节日做什么？} (Kāmàilóngrén zài chuántǒng jiérì zuò shénme ? — Que font les Camerounais pendant les fêtes traditionnelles ?)\\
Mot interrogatif : \zh{什么} après le verbe \zh{做} → on cherche des \emph{activités}.\\
Réponse : \zh{他们穿漂亮的传统服装，表演节目。} (Tāmen chuān piàoliang de chuántǒng fúzhuāng, biǎoyǎn jiémù. — Ils portent de beaux vêtements traditionnels et présentent des numéros.)

\textbf{Question 3.} \zh{两个国家的人民都喜欢什么？} (Liǎng ge guójiā de rénmín dōu xǐhuan shénme ? — Qu'aiment les peuples des deux pays ?)\\
Réponse : \zh{他们都很喜欢音乐和舞蹈。} (Tāmen dōu hěn xǐhuan yīnyuè hé wǔdǎo. — Ils aiment tous beaucoup la musique et la danse.)

\textbf{Question 4.} \zh{文化让我们做什么？} (Wénhuà ràng wǒmen zuò shénme ? — Que nous permet la culture ?)\\
Réponse : \zh{文化让我们互相了解。} (Wénhuà ràng wǒmen hùxiāng liǎojiě. — La culture nous permet de nous comprendre mutuellement.)
\end{exemplebox}

\begin{textebox}[Le point commun des deux peuples]
两个国家的人民都很喜欢音乐和舞蹈。节日的时候，大家一起唱歌、跳舞，非常快乐。文化让我们互相了解，也让我们的世界更丰富。\\
Liǎng ge guójiā de rénmín dōu hěn xǐhuan yīnyuè hé wǔdǎo. Jiérì de shíhou, dàjiā yìqǐ chànggē、tiàowǔ, fēicháng kuàilè. Wénhuà ràng wǒmen hùxiāng liǎojiě, yě ràng wǒmen de shìjiè gèng fēngfù.\\
« Les peuples des deux pays aiment beaucoup la musique et la danse. Au moment des fêtes, tout le monde chante et danse ensemble, dans une grande joie. La culture nous permet de nous comprendre mutuellement et rend notre monde plus riche. »
\end{textebox}

\begin{attention}[\zh{了解} et \zh{知道} : deux verbes à ne pas confondre]
\zh{了解} (liǎojiě) signifie « comprendre, connaître en profondeur » une personne, une culture, une situation : \zh{我想了解中国文化。} (Wǒ xiǎng liǎojiě Zhōngguó wénhuà. — Je veux connaître la culture chinoise.)\\
\zh{知道} (zhīdào) signifie « savoir » un fait, une information précise : \zh{我知道他的名字。} (Wǒ zhīdào tā de míngzi. — Je connais son nom.)\\
Piège : on dit \zh{互相了解} (se comprendre mutuellement), jamais \zh{互相知道}.
\end{attention}

\courssub{Les collocations clés du thème}

\begin{definition}[Qu'est-ce qu'une collocation ?]
Une \cle{collocation} est une association habituelle et stable de deux ou plusieurs mots. En chinois comme en français, certains verbes « appellent » certains noms : on dit \zh{过节日} (célébrer une fête) et non \zh{做节日}. Apprendre les collocations par blocs, c'est parler un chinois naturel au lieu d'un chinois traduit mot à mot du français.
\end{definition}

\begin{center}
\begin{tabularx}{\linewidth}{|X|X|X|X|}
\hline
\rowcolor{popblueL} Collocation & Pinyin & Structure & Traduction \\
\hline
过节日 & guò jiérì & verbe + nom & célébrer une fête \\
\hline
过春节 & guò Chūnjié & verbe + nom propre & célébrer la fête du Printemps \\
\hline
表演节目 & biǎoyǎn jiémù & verbe + nom & présenter un numéro (spectacle) \\
\hline
穿传统服装 & chuān chuántǒng fúzhuāng & verbe + nom & porter des vêtements traditionnels \\
\hline
互相了解 & hùxiāng liǎojiě & adverbe + verbe & se comprendre mutuellement \\
\hline
丰富的文化 & fēngfù de wénhuà & adjectif + 的 + nom & une culture riche \\
\hline
\end{tabularx}
\end{center}

\begin{propriete}[Le verbe \zh{过} pour les fêtes et le temps qui passe]
En chinois, on « passe » une fête comme on passe du temps : \zh{过} (guò) + fête. Structure : \textbf{Sujet + 过 + fête}.\\
\zh{我们过新年。} (Wǒmen guò xīnnián. — Nous célébrons le Nouvel An.)\\
\zh{喀麦隆人过传统节日。} (Kāmàilóngrén guò chuántǒng jiérì. — Les Camerounais célèbrent les fêtes traditionnelles.)\\
De même, \zh{过生日} (guò shēngrì) signifie « fêter son anniversaire ». Ne traduis jamais « célébrer » par un verbe inventé : retiens le bloc \zh{过 + fête}.
\end{propriete}

\begin{propriete}[L'adverbe \zh{互相} « mutuellement »]
\zh{互相} (hùxiāng) se place \textbf{avant le verbe} : \textbf{Sujet + 互相 + verbe}.\\
\zh{我们互相学习。} (Wǒmen hùxiāng xuéxí. — Nous apprenons l'un de l'autre.)\\
\zh{两国人民互相了解。} (Liǎng guó rénmín hùxiāng liǎojiě. — Les peuples des deux pays se comprennent mutuellement.)\\
Attention à la place : en français, « mutuellement » peut se mettre à la fin ; en chinois, \zh{互相} précède toujours le verbe.
\end{propriete}

\begin{exemplebox}[Les collocations en contexte]
\zh{春节的时候，中国人穿传统服装，看舞龙舞狮。} (Chūnjié de shíhou, Zhōngguórén chuān chuántǒng fúzhuāng, kàn wǔlóng wǔshī.) « À la fête du Printemps, les Chinois portent des vêtements traditionnels et regardent les danses du dragon et du lion. »\\
\zh{学校的学生表演节目，庆祝文化节。} (Xuéxiào de xuésheng biǎoyǎn jiémù, qìngzhù wénhuàjié.) « Les élèves de l'école présentent des numéros pour célébrer la fête de la culture. »\\
\zh{中国有丰富的文化，喀麦隆也有。} (Zhōngguó yǒu fēngfù de wénhuà, Kāmàilóng yě yǒu.) « La Chine possède une culture riche, et le Cameroun aussi. »
\end{exemplebox}

\begin{popfigure}
\begin{tikzpicture}[
  node distance=0.6cm and 1.4cm,
  box/.style={draw=popblue, fill=popblueL, rounded corners, thick, align=center, minimum height=0.9cm, text width=3.1cm, font=\small},
  arr/.style={-{Stealth[length=3mm]}, thick, popdark}
]
\node[box, align=center] (texte){\textbf{Texte}\\ lire et souligner\\ les mots-clés};
\node[box, right=of texte, fill=poporangeL, draw=poporange, align=center] (mots){\textbf{Mots-clés}\\ 节日、服装、\\ 音乐、文化};
\node[box, right=of mots, fill=popgreenL, draw=popgreen, align=center] (idees){\textbf{Idées principales}\\ fêtes, activités,\\ points communs};
\node[box, right=of idees, fill=poppurpleL, draw=poppurple, align=center] (rep){\textbf{Réponse complète}\\ phrase entière\\ en chinois};
\draw[arr] (texte) -- (mots);
\draw[arr] (mots) -- (idees);
\draw[arr] (idees) -- (rep);
\end{tikzpicture}
\legende{Organigramme de compréhension : du texte à la réponse complète.}
\end{popfigure}

\courssub{Exercices de consolidation}

\begin{exoresolu}[Appariement : collocation et traduction]
Énoncé. Relie chaque collocation chinoise à sa traduction française :\\
1. 过节日 \hspace{0.5cm} 2. 表演节目 \hspace{0.5cm} 3. 穿传统服装 \hspace{0.5cm} 4. 互相了解 \hspace{0.5cm} 5. 丰富的文化\\
a. porter des vêtements traditionnels \hspace{0.3cm} b. une culture riche \hspace{0.3cm} c. célébrer une fête \hspace{0.3cm} d. se comprendre mutuellement \hspace{0.3cm} e. présenter un numéro
\tcblower
Solution détaillée. On raisonne mot par mot :\\
1. \zh{过} (célébrer, passer) + \zh{节日} (fête) → \textbf{1-c} : célébrer une fête.\\
2. \zh{表演} (présenter, jouer) + \zh{节目} (numéro, programme) → \textbf{2-e} : présenter un numéro.\\
3. \zh{穿} (porter, mettre) + \zh{传统服装} (vêtements traditionnels) → \textbf{3-a} : porter des vêtements traditionnels.\\
4. \zh{互相} (mutuellement) + \zh{了解} (comprendre en profondeur) → \textbf{4-d} : se comprendre mutuellement.\\
5. \zh{丰富} (riche, abondant) + \zh{的} + \zh{文化} (culture) → \textbf{5-b} : une culture riche.
\end{exoresolu}

\begin{exoresolu}[Vrai ou Faux ?]
Énoncé. Lis la phrase et dis si elle est vraie (V) ou fausse (F) d'après le texte \emph{Le point commun des deux peuples} :\\
\zh{中国人过节日的时候一起唱歌、跳舞。} (Zhōngguórén guò jiérì de shíhou yìqǐ chànggē、tiàowǔ.)
\tcblower
Solution. \textbf{Vrai (V).} Le texte indique explicitement : \zh{节日的时候，大家一起唱歌、跳舞} « Au moment des fêtes, tout le monde chante et danse ensemble ». La phrase de l'exercice reprend cette information avec le sujet \zh{中国人} et la structure \zh{……的时候}.
\end{exoresolu}

\begin{exercice}[À toi de répondre]
Réponds par des phrases complètes en chinois :
\begin{enumerate}
\item \zh{中国人过什么节日？}
\item \zh{喀麦隆人在传统节日做什么？}
\item \zh{两个国家的人民都喜欢什么？}
\item \zh{文化让我们做什么？}
\item Traduis en chinois : « Les élèves portent de beaux vêtements traditionnels et présentent des numéros. »
\end{enumerate}
\end{exercice}

\begin{corrige}[Exercice « À toi de répondre »]
\begin{enumerate}
\item \zh{中国人过春节。} (Zhōngguórén guò Chūnjié.)
\item \zh{他们穿漂亮的传统服装，表演节目。} (Tāmen chuān piàoliang de chuántǒng fúzhuāng, biǎoyǎn jiémù.)
\item \zh{他们都很喜欢音乐和舞蹈。} (Tāmen dōu hěn xǐhuan yīnyuè hé wǔdǎo.)
\item \zh{文化让我们互相了解。} (Wénhuà ràng wǒmen hùxiāng liǎojiě.)
\item \zh{学生们穿漂亮的传统服装，表演节目。} (Xuéshengmen chuān piàoliang de chuántǒng fúzhuāng, biǎoyǎn jiémù.)
\end{enumerate}
\end{corrige}

\begin{aretenir}[L'essentiel de la section]
\begin{itemize}
\item Pour répondre à une question : repérer le mot interrogatif (\zh{什么}, \zh{谁}, \zh{哪儿}), retrouver la phrase du texte, répondre par une phrase complète en miroir.
\item Le mot interrogatif reste à la place du mot interrogé : pas d'inversion comme en français.
\item Collocations à connaître par cœur : \zh{过节日}, \zh{表演节目}, \zh{穿传统服装}, \zh{互相了解}, \zh{丰富的文化}.
\item \zh{了解} = comprendre en profondeur ; \zh{知道} = savoir un fait.
\end{itemize}
\end{aretenir}

\coursec{Production guidée et ouverture interculturelle}

Dans la section précédente, nous avons vérifié notre compréhension du texte et appris à combiner les mots du glossaire en collocations naturelles : \zh{过节日} (\emph{guò jiérì}, célébrer une fête), \zh{穿传统服装} (\emph{chuān chuántǒng fúzhuāng}, porter des vêtements traditionnels), \zh{唱歌跳舞} (\emph{chàng gē tiào wǔ}, chanter et danser). Il est temps maintenant de \cle{produire} : parler et écrire sur \emph{notre} vie culturelle, en comparant le Cameroun et la Chine. C'est le cœur de la compétence « ouverture au monde » du programme de Première.

\courssub{Comparer deux fêtes nationales : le 20 mai et le 1er octobre}

Commençons par observer deux phrases modèles, l'une sur le Cameroun, l'autre sur la Chine.

\begin{textebox}[Deux fêtes nationales]
在喀麦隆，我们过五月二十日的国庆节。人们唱歌、跳舞、穿传统服装。我很喜欢我们的文化。\\
\emph{Zài Kāmàilóng, wǒmen guò wǔ yuè èrshí rì de guóqìngjié. Rénmen chàng gē, tiào wǔ, chuān chuántǒng fúzhuāng. Wǒ hěn xǐhuan wǒmen de wénhuà.}\\
« Au Cameroun, nous célébrons la fête nationale du 20 mai. Les gens chantent, dansent et portent des vêtements traditionnels. J'aime beaucoup notre culture. »\\[4pt]
在中国，人们过十月一日的国庆节。很多人看游行，也和家人一起吃饭。\\
\emph{Zài Zhōngguó, rénmen guò shí yuè yī rì de guóqìngjié. Hěn duō rén kàn yóuxíng, yě hé jiārén yìqǐ chīfàn.}\\
« En Chine, on célèbre la fête nationale du 1er octobre. Beaucoup de gens regardent le défilé et mangent aussi en famille. »
\end{textebox}

Observons : les deux textes suivent exactement la même architecture. C'est cette architecture que vous allez réutiliser.

\begin{center}
\begin{tabularx}{\linewidth}{|X|X|X|}
\hline
\rowcolor{popblueL} \textbf{Étape} & \textbf{Structure chinoise} & \textbf{Exemple} \\
\hline
1. Lieu & 在 + pays & 在喀麦隆，… (Au Cameroun, …) \\
\hline
2. Fête et date & 过 + date + 的 + fête & 过五月二十日的国庆节 \\
\hline
3. Activités & 人们 + verbes & 人们唱歌、跳舞。 \\
\hline
4. Avis personnel & 我很 / 我也 + 喜欢 & 我很喜欢我们的文化。 \\
\hline
\end{tabularx}
\end{center}

\begin{propriete}[Le verbe 过 : « célébrer, passer (une fête) »]
Le verbe \cle{过} (\emph{guò}) signifie littéralement « passer » ; avec un nom de fête, il signifie « célébrer, fêter ». Structure : \textbf{Sujet + 过 + (date + 的 +) fête}.\\
Ex. : \zh{我们过国庆节。} \emph{Wǒmen guò guóqìngjié.} « Nous célébrons la fête nationale. »\\
Ex. : \zh{中国人过春节。} \emph{Zhōngguórén guò Chūnjié.} « Les Chinois célèbrent la fête du Printemps. »\\
Remarque : le complément de lieu (\zh{在喀麦隆}, \zh{在中国}) se place \textbf{avant le sujet} ou juste après lui, jamais en fin de phrase comme en français.
\end{propriete}

\begin{attention}[La date chinoise : du grand au petit]
En chinois, la date se dit toujours \cle{du grand au petit} : année → mois → jour. On dit \zh{五月二十日} (\emph{wǔ yuè èrshí rì}, littéralement « mois 5, jour 20 »), jamais l'inverse. Le français dit « le 20 mai » (jour puis mois) : ne traduisez donc jamais mot à mot. De même : 十月一日 = le 1er octobre ; 二月十一日 = le 11 février (fête de la Jeunesse au Cameroun). Rappel : devant 日, le chiffre du jour se lit tel quel (二十日 = le 20, 一日 = le 1er).
\end{attention}

\begin{savaistu}[Les enveloppes rouges du Nouvel An chinois]
Pendant la fête du Printemps (\zh{春节} Chūnjié), les Chinois offrent aux enfants et aux proches des \cle{enveloppes rouges}, \zh{红包} (\emph{hóngbāo}), contenant de l'argent porte-bonheur. Le rouge (\zh{红} hóng) symbolise la joie, la chance et éloigne le malheur : on le retrouve partout pendant les fêtes — lanternes, vêtements, décorations. Au Cameroun aussi, les fêtes sont l'occasion de donner et de partager : offrir un repas, un pagne neuf, recevoir la famille. Les gestes diffèrent, l'esprit est proche.
\end{savaistu}

\begin{popfigure}
\begin{tikzpicture}[scale=1.0]
% Ligne Cameroun
\draw[line width=2pt, popgreen] (0,2.2) -- (12,2.2);
\node[anchor=east, font=\bfseries, text=popgreen] at (-0.2,2.2) {Cameroun};
\foreach \x/\d/\f in {1.2/{1\textsuperscript{er} janv.}/{Jour de l'An}, 5.5/{11 février}/{Fête de la Jeunesse}, 9.8/{20 mai}/{Fête nationale}} {
  \fill[popgreen] (\x,2.2) circle (0.12);
  \node[above, align=center, font=\small] at (\x,2.35) {\textbf{\d}\\\f};
}
% Ligne Chine
\draw[line width=2pt, poporange] (0,0) -- (12,0);
\node[anchor=east, font=\bfseries, text=poporange] at (-0.2,0) {Chine};
\foreach \x/\d/\f in {1.2/{春节}/{Fête du Printemps}, 4.6/{端午节}/{Bateaux-dragons}, 7.9/{中秋节}/{Mi-automne}, 10.8/{国庆节}/{1\textsuperscript{er} octobre}} {
  \fill[poporange] (\x,0) circle (0.12);
  \node[below, align=center, font=\small] at (\x,-0.15) {\textbf{\d}\\\f};
}
\end{tikzpicture}
\legende{Frise comparée des grandes fêtes : Cameroun (en haut) et Chine (en bas).}
\end{popfigure}

\courssub{Les trois connecteurs de la production : 和, 也, 都}

Pour écrire un paragraphe fluide, trois petits mots suffisent. Observons d'abord.

\begin{exemplebox}[Observer avant de généraliser]
\zh{我也喜欢中国的春节。} \emph{Wǒ yě xǐhuan Zhōngguó de Chūnjié.} « Moi aussi, j'aime la fête du Printemps chinoise. »\\
\zh{我和我的朋友都喜欢跳舞。} \emph{Wǒ hé wǒ de péngyou dōu xǐhuan tiàowǔ.} « Mon ami et moi aimons tous les deux danser. »\\
\zh{我们唱歌，也跳舞。} \emph{Wǒmen chàng gē, yě tiàowǔ.} « Nous chantons et nous dansons aussi. »
\end{exemplebox}

\begin{propriete}[Forme et emploi des trois connecteurs]
\begin{itemize}
\item \cle{和} (hé) « et » : coordonne des \textbf{noms} ou des \textbf{pronoms}. Structure : A + 和 + B. Ex. : 我和我弟弟 (mon petit frère et moi). \emph{Exception majeure} : 和 ne relie \textbf{jamais deux phrases ni deux verbes} ; pour les actions, on les juxtapose (avec la virgule chinoise 、 ou ，) ou on utilise 也.
\item \cle{也} (yě) « aussi » : adverbe placé \textbf{après le sujet, avant le verbe}. Structure : Sujet + 也 + verbe. Ex. : 我也喜欢。 (Moi aussi j'aime.)
\item \cle{都} (dōu) « tous, toutes » : adverbe placé \textbf{après le sujet pluriel, avant le verbe}. Structure : Sujet (pluriel) + 都 + verbe. Ex. : 我们都过国庆节。 (Nous célébrons tous la fête nationale.)
\end{itemize}
Ordre quand les deux adverbes se rencontrent : 也 se place avant 都 → \zh{我们也都很喜欢。} « Nous aussi, nous aimons tous. »
\end{propriete}

\begin{attention}[Pièges fréquents des francophones]
\begin{enumerate}
\item Ne dites pas « \zh{我唱歌和跳舞} » pour « je chante et je danse » : 和 ne relie pas deux verbes. Dites \zh{我唱歌、跳舞} ou \zh{我唱歌，也跳舞}.
\item 也 et 都 se placent \textbf{avant} le verbe, jamais en fin de phrase comme le français « aussi » : « \zh{我喜欢春节也} » est fautif ; dites \zh{我也喜欢春节}.
\item Avec 和, le sujet devient pluriel : 我和他都… (lui et moi, nous… tous les deux).
\item La virgule de coordination chinoise 、 (dite « virgule de pause ») s'emploie entre des mots ou groupes de mots coordonnés : \zh{唱歌、跳舞、穿传统服装}. Ne la confondez pas avec la virgule ordinaire ，.
\end{enumerate}
\end{attention}

\begin{exoresolu}[Transformer avec 和, 也, 都]
Énoncé. Combinez en une seule phrase correcte : (a) 我喜欢音乐。/ 我喜欢舞蹈。(avec 也) ; (b) 我哥哥喜欢春节。/ 我喜欢春节。(avec 和…都) ; (c) 喀麦隆人过国庆节。/ 中国人过国庆节。(avec 都).
\tcblower
Solution détaillée.\\
(a) Deux verbes se succèdent : on ajoute 也 avant le second verbe → \zh{我喜欢音乐，也喜欢舞蹈。} \emph{Wǒ xǐhuan yīnyuè, yě xǐhuan wǔdǎo.} « J'aime la musique et j'aime aussi la danse. »\\
(b) On fusionne les sujets avec 和, puis 都 avant le verbe → \zh{我和我哥哥都喜欢春节。} \emph{Wǒ hé wǒ gēge dōu xǐhuan Chūnjié.} « Mon grand frère et moi aimons tous les deux la fête du Printemps. »\\
(c) Les deux sujets forment un groupe : on peut écrire \zh{喀麦隆人和中国人都过国庆节。} \emph{Kāmàilóngrén hé Zhōngguórén dōu guò guóqìngjié.} « Les Camerounais et les Chinois célèbrent tous leur fête nationale. »
\end{exoresolu}

\courssub{Rédaction guidée : « Ma fête culturelle préférée »}

\begin{methode}[Rédiger en cinq étapes]
\begin{enumerate}
\item \textbf{Choisir} sa fête : 国庆节 (fête nationale), 青年节 (fête de la Jeunesse), 春节, 中秋节…
\item \textbf{Situer} : 在 + lieu, puis la date avec 过 + mois + jour + 的 + fête.
\item \textbf{Décrire} 2 ou 3 activités avec le glossaire : 唱歌、跳舞、穿传统服装、吃好吃的、看游行、和家人一起…
\item \textbf{Lier} avec 和, 也, 都 (au moins une fois chacun).
\item \textbf{Conclure} par un avis : 我很喜欢… / 我觉得…很有意思。
\end{enumerate}
Consigne : 5 phrases minimum, au moins 8 mots du glossaire de la leçon.
\end{methode}

\begin{textebox}[Modèle de rédaction — 我最喜欢的节日]
我最喜欢的节日是五月二十日的国庆节。在雅温得，我们都穿传统服装。人们唱歌、跳舞，也看游行。我和我的家人一起吃好吃的。我很喜欢我们的文化，也很喜欢中国的节日。\\
\emph{Wǒ zuì xǐhuan de jiérì shì wǔ yuè èrshí rì de guóqìngjié. Zài Yǎwēndé, wǒmen dōu chuān chuántǒng fúzhuāng. Rénmen chàng gē, tiào wǔ, yě kàn yóuxíng. Wǒ hé wǒ de jiārén yìqǐ chī hǎochī de. Wǒ hěn xǐhuan wǒmen de wénhuà, yě hěn xǐhuan Zhōngguó de jiérì.}\\
« Ma fête préférée, c'est la fête nationale du 20 mai. À Yaoundé, nous portons tous des vêtements traditionnels. Les gens chantent, dansent et regardent aussi le défilé. Ma famille et moi mangeons ensemble de bons plats. J'aime beaucoup notre culture et j'aime aussi beaucoup les fêtes chinoises. »
\end{textebox}

\begin{exercice}[À vous d'écrire]
Rédigez votre paragraphe « Ma fête culturelle préférée » (5 phrases minimum, 8 mots du glossaire, les connecteurs 和, 也, 都). Puis, en binôme : relisez le texte de votre voisin et vérifiez (1) la date du grand au petit, (2) la place de 也 et 都 avant le verbe, (3) l'usage de 和 uniquement entre noms. La classe corrige ensuite collectivement un exemple copié au tableau.
\end{exercice}

\begin{corrige}[Grille de relecture]
Points à valider chez le pair : date correcte (\zh{五月二十日} et non « 二十日五月 ») ; sujet + 也/都 + verbe ; 和 entre noms seulement ; ponctuation chinoise ，。 ; au moins 8 mots du glossaire soulignés. Erreur type à corriger au tableau : « \zh{我也很喜欢} » (correct) contre « \zh{我很喜欢也} » (fautif).
\end{corrige}

\courssub{Complément utile à l'examen : deux fêtes chinoises fréquentes}

\emph{(Complément signalé : ces deux fêtes reviennent souvent dans les textes de compréhension écrite des évaluations.)}

\begin{center}
\begin{tabularx}{\linewidth}{|X|X|X|X|}
\hline
\rowcolor{popblueL} Caractères & Pinyin & Nature & Traduction \\
\hline
中秋节 & Zhōngqiūjié & n. & fête de la mi-automne \\
\hline
端午节 & Duānwǔjié & n. & fête des bateaux-dragons \\
\hline
月饼 & yuèbing & n. & gâteau de lune \\
\hline
龙舟 & lóngzhōu & n. & bateau-dragon \\
\hline
红包 & hóngbāo & n. & enveloppe rouge (porte-bonheur) \\
\hline
团圆 & tuányuán & v./adj. & se réunir (en famille) \\
\hline
\end{tabularx}
\end{center}

\begin{exemplebox}[Deux phrases à retenir]
\zh{中秋节的时候，我们吃月饼。} \emph{Zhōngqiūjié de shíhou, wǒmen chī yuèbing.} « À la fête de la mi-automne, on mange des gâteaux de lune. »\\
\zh{端午节有龙舟比赛，很有意思。} \emph{Duānwǔjié yǒu lóngzhōu bǐsài, hěn yǒu yìsi.} « À la fête des bateaux-dragons, il y a des courses de bateaux-dragons, c'est très intéressant. »
\end{exemplebox}

\begin{experience}[Débat interculturel en binômes]
En binôme, chacun présente « sa » fête (l'un une fête camerounaise, l'autre une fête chinoise) en 4 phrases avec 和, 也, 都. Puis trouvez ensemble une phrase de conclusion commune du type : \zh{我们的节日不一样，但是我们都喜欢和家人一起过。} \emph{Wǒmen de jiérì bù yíyàng, dànshì wǒmen dōu xǐhuan hé jiārén yìqǐ guò.} « Nos fêtes sont différentes, mais nous aimons tous les célébrer en famille. »
\end{experience}

\begin{aretenir}[L'essentiel de la section]
\begin{itemize}
\item La date chinoise va du grand au petit : \zh{五月二十日} = le 20 mai ; \zh{十月一日} = le 1er octobre.
\item 过 + fête = « célébrer » ; le complément de lieu (在 + lieu) se place avant le verbe.
\item 和 relie des noms ; 也 et 都 se placent avant le verbe.
\item Plan de production : lieu → fête et date → activités → avis personnel.
\item Pour l'examen : connaître 春节, 中秋节, 端午节, 国庆节 et le vocabulaire associé.
\end{itemize}
\end{aretenir}

\newpage

\coursec{Activité d'intégration}

\courssub{Objectifs de l'activité}

À la fin de cette activité d'intégration, tu seras capable de :

\begin{itemize}
\item mobiliser le \cle{vocabulaire de la vie culturelle} (fêtes, musique, danse, vêtements traditionnels) dans une production réelle ;
\item réutiliser les structures de \cle{comparaison} et les mots-outils \zh{和} (hé, « et »), \zh{也} (yě, « aussi »), \zh{都} (dōu, « tous/tout ») et \zh{因为} (yīnwèi, « parce que ») ;
\item lire et comprendre un court message authentique en chinois simplifié ;
\item produire une \cle{affiche bilingue} chinois-français et la présenter à l'oral en cinq phrases minimum ;
\item développer ton ouverture interculturelle en comparant une fête camerounaise et une fête chinoise de manière factuelle et respectueuse.
\end{itemize}

\courssub{Situation}

Ton lycée de Yaoundé participe à un \cle{festival interculturel} organisé avec une école partenaire de Pékin. Des correspondants chinois suivront le festival en visioconférence. Ton professeur de chinois a reçu ce message de la professeure chinoise, Madame Wang :

\begin{textebox}[Message de Madame Wang (Wáng lǎoshī)]
亲爱的同学们，你们好！\\
Qīn'ài de tóngxuémen, nǐmen hǎo !\\
« Chers élèves, bonjour ! »\\
下个月我们学校有一个文化节。我们很想知道喀麦隆的年轻人怎么庆祝节日。\\
Xià ge yuè wǒmen xuéxiào yǒu yí ge wénhuàjié. Wǒmen hěn xiǎng zhīdào Kāmàilóng de niánqīngrén zěnme qìngzhù jiérì.\\
« Le mois prochain, notre école organise un festival culturel. Nous aimerions beaucoup savoir comment les jeunes du Cameroun célèbrent leurs fêtes. »\\
请你们做一张海报，介绍一个喀麦隆的节日：名字、时间、活动（唱歌、跳舞、穿传统服装……）。也请你们比较一下中国的节日。\\
Qǐng nǐmen zuò yì zhāng hǎibào, jièshào yí ge Kāmàilóng de jiérì : míngzi, shíjiān, huódòng (chàng gē, tiào wǔ, chuān chuántǒng fúzhuāng……). Yě qǐng nǐmen bǐjiào yíxià Zhōngguó de jiérì.\\
« Faites une affiche qui présente une fête camerounaise : nom, date, activités (chanter, danser, porter des vêtements traditionnels...). Comparez-la aussi avec une fête chinoise, s'il vous plaît. »\\
谢谢！我们等你们的海报！\\
Xièxie ! Wǒmen děng nǐmen de hǎibào !\\
« Merci ! Nous attendons vos affiches ! »
\end{textebox}

Ta classe répond à cette demande : chaque binôme prépare une affiche bilingue et la présente en chinois devant la classe. Les meilleures affiches seront envoyées à Pékin.

\begin{attention}[Lecture du message : trois points de vigilance]
\begin{itemize}
\item \zh{怎么} (zěnme) signifie « comment » : \zh{怎么庆祝节日} = « comment célèbre-t-on les fêtes ». Ne le confonds pas avec \zh{什么} (shénme, « quoi »).
\item \zh{一张海报} : le classificateur des objets plats et en feuille (affiche, papier, photo) est \zh{张} (zhāng), pas \zh{个}.
\item \zh{比较一下} (bǐjiào yíxià) : « comparer un peu » ; le suffixe \zh{一下} adoucit la demande, comme « s'il vous plaît ».
\end{itemize}
\end{attention}

\courssub{Consignes}

\begin{enumerate}
\item \textbf{Compréhension du message.} Lis le message de Madame Wang et réponds par écrit en français : que demande-t-elle exactement ? Quels sont les trois éléments obligatoires de l'affiche ?
\item \textbf{Choix de la fête.} Avec ton binôme, choisis une fête culturelle camerounaise (par exemple la fête de la jeunesse, un festival traditionnel de ta région, la fête de l'indépendance) et une fête chinoise (la Fête du Printemps \zh{春节} Chūnjié, ou la Fête de la Lune \zh{中秋节} Zhōngqiūjié).
\item \textbf{Préparation du lexique.} Surligne dans le glossaire de la leçon au moins \cle{huit mots} que tu vas utiliser. Vérifie les tons en pinyin avant de les écrire sur l'affiche.
\item \textbf{Confection de l'affiche bilingue.} Rédige le titre, le nom de la fête, la date et trois ou quatre activités en chinois avec la traduction française. Soigne l'écriture des caractères : ordre des traits, équilibre, taille régulière.
\item \textbf{Préparation de la présentation orale.} Écris un script d'au moins \cle{cinq phrases} en chinois. Obligatoire : utiliser \zh{和}, \zh{也}, \zh{都} et une comparaison avec \zh{比} ou \zh{和……一样}.
\item \textbf{Présentation.} Chaque binôme présente son affiche en deux minutes. Les autres élèves écoutent et notent les mots du glossaire qu'ils reconnaissent.
\item \textbf{Vote argumenté.} La classe vote pour la fête qu'elle aimerait découvrir. Chaque électeur justifie son choix à l'oral avec \zh{因为} : \zh{我想去中国，因为春节很有意思。} (Wǒ xiǎng qù Zhōngguó, yīnwèi Chūnjié hěn yǒu yìsi. « Je veux aller en Chine parce que la Fête du Printemps est très intéressante. »)
\end{enumerate}

\courssub{Ressources à mobiliser}

\begin{aretenir}[Boîte à outils de l'activité]
\begin{itemize}
\item \textbf{Lexique :} 节日 (jiérì, fête), 庆祝 (qìngzhù, célébrer), 唱歌 (chàng gē, chanter), 跳舞 (tiào wǔ, danser), 穿 (chuān, porter), 传统 (chuántǒng, traditionnel), 服装 (fúzhuāng, vêtement), 文化 (wénhuà, culture), 活动 (huódòng, activité), 海报 (hǎibào, affiche).
\item \textbf{Coordination :} A \zh{和} B (A et B) — uniquement entre noms ou verbes, jamais entre deux phrases.
\item \textbf{Adverbes :} Sujet + \zh{也} + verbe (« aussi ») ; Sujets pluriels + \zh{都} + verbe (« tous »). Attention : \zh{也} et \zh{都} se placent \emph{avant} le verbe.
\item \textbf{Comparaison :} A + \zh{比} + B + adjectif (« A est plus... que B ») ; A \zh{和} B \zh{一样} + adjectif (« A est aussi... que B »).
\item \textbf{Cause :} phrase + \zh{因为} + cause (« parce que »). En chinois soutenu, on peut répondre avec \zh{所以} (suǒyǐ, « donc ») : \zh{因为……所以……}, mais une seule des deux particules suffit à l'oral.
\end{itemize}
\end{aretenir}

\begin{attention}[Pièges fréquents]
Ne dis pas « \zh{我喜欢唱歌，和跳舞。} » pour relier deux phrases : \zh{和} ne relie pas deux propositions ; écris \zh{我喜欢唱歌，也喜欢跳舞。} (Wǒ xǐhuan chàng gē, yě xǐhuan tiào wǔ. « J'aime chanter et j'aime aussi danser. »). Ne place jamais \zh{也} ou \zh{都} après le verbe : on dit \zh{我们都跳舞。} (Wǒmen dōu tiào wǔ. « Nous dansons tous. ») et jamais « \zh{我们跳舞都} ». Enfin, dans la comparaison avec \zh{比}, n'ajoute pas \zh{很} devant l'adjectif : on dit \zh{春节比青年节长。} et non « \zh{春节比青年节很长} ».
\end{attention}

\courssub{Proposition de production et corrigé commenté}

\begin{exoresolu}[Affiche et présentation modèles]
\textbf{Énoncé.} Rédige l'affiche bilingue d'une fête camerounaise et le script de présentation (5 phrases minimum, 8 mots du glossaire, \zh{和}, \zh{也}, \zh{都}).
\tcblower
\textbf{Affiche (texte chinois).}\\
喀麦隆青年节 — Kāmàilóng Qīngniánjié — « La fête de la jeunesse au Cameroun »\\
时间：二月十一日。Shíjiān : èr yuè shíyī rì. « Date : le 11 février. »\\
活动：唱歌、跳舞、穿传统服装、庆祝！Huódòng : chàng gē, tiào wǔ, chuān chuántǒng fúzhuāng, qìngzhù ! « Activités : chanter, danser, porter les vêtements traditionnels, célébrer ! »\\[0.3em]
\textbf{Script de présentation.}\\
大家好！我们介绍喀麦隆的青年节。\\
Dàjiā hǎo ! Wǒmen jièshào Kāmàilóng de Qīngniánjié.\\
« Bonjour à tous ! Nous présentons la fête de la jeunesse du Cameroun. »\\
青年节在二月十一日，年轻人都庆祝这个节日。\\
Qīngniánjié zài èr yuè shíyī rì, niánqīngrén dōu qìngzhù zhè ge jiérì.\\
« La fête de la jeunesse a lieu le 11 février ; les jeunes célèbrent tous cette fête. »\\
我们唱歌和跳舞，也穿漂亮的传统服装。\\
Wǒmen chàng gē hé tiào wǔ, yě chuān piàoliang de chuántǒng fúzhuāng.\\
« Nous chantons et dansons, et nous portons aussi de beaux vêtements traditionnels. »\\
中国的春节和青年节一样有意思，但是春节比青年节长。\\
Zhōngguó de Chūnjié hé Qīngniánjié yíyàng yǒu yìsi, dànshì Chūnjié bǐ Qīngniánjié cháng.\\
« La Fête du Printemps chinoise est aussi intéressante que la fête de la jeunesse, mais elle dure plus longtemps. »\\
我们都很喜欢这两个节日，因为它们很有意思！\\
Wǒmen dōu hěn xǐhuan zhè liǎng ge jiérì, yīnwèi tāmen hěn yǒu yìsi !\\
« Nous aimons beaucoup ces deux fêtes, parce qu'elles sont très intéressantes ! »\\[0.3em]
\textbf{Commentaire des choix de langue.} Le script contient huit mots du glossaire : \zh{节日}, \zh{庆祝}, \zh{唱歌}, \zh{跳舞}, \zh{穿}, \zh{传统}, \zh{服装}, \zh{活动} (sur l'affiche). \zh{和} relie deux verbes (\zh{唱歌和跳舞}) ; \zh{也} et \zh{都} sont bien placés avant le verbe ; la comparaison utilise d'abord \zh{和……一样} puis \zh{比} + adjectif sans \zh{很} ; la justification finale emploie \zh{因为}, structure valorisée à l'examen. Note aussi le classificateur \zh{个} dans \zh{这个节日} et \zh{两个节日} : \zh{两} (liǎng) s'emploie devant un classificateur, jamais \zh{二}.
\end{exoresolu}

\courssub{Bilan de l'activité}

Cette activité t'a permis de passer de la \emph{reconnaissance} du vocabulaire à son \emph{emploi réel} dans une situation de communication authentique : répondre à une demande venue de Chine. Tu as lu un message en chinois sans traduction mot à mot, sélectionné le lexique pertinent, construit des phrases avec les mots-outils \zh{和}, \zh{也}, \zh{都} et justifié une opinion avec \zh{因为}. Tu as aussi comparé deux cultures sans les hiérarchiser : les fêtes camerounaises et chinoises rassemblent toutes les deux la famille, la musique et la joie. Pour t'auto-évaluer, vérifie : ai-je utilisé huit mots du glossaire ? Mes tons sont-ils corrects ? Mes caractères respectent-ils l'ordre des traits ? Si oui, tu es prêt pour l'évaluation du module.

\newpage

\coursec{Méthodes et exercices résolus}

\coursec{Leçon 19 — Vocabulaire et compréhension de texte : vie culturelle au Cameroun et en Chine}

\courssub{Mise en route et découverte du thème}

\begin{textebox}[Situation de communication : une conversation au lycée]
\textbf{Contexte :} Deux élèves, Paul (Cameroun) et Wang Wei (Chine), discutent pendant la récréation au lycée de Yaoundé.\\
\zh{保罗：王伟，周末你去了博物馆吗？}\\
\emph{Bǎoluó: Wáng Wěi, zhōumò nǐ qùle bówùguǎn ma?}\\
\zh{王伟：去了。雅温得的国家博物馆很有意思，我看了很多非洲艺术。你呢？}\\
\emph{Wáng Wěi: Qùle. Yǎwēndé de guójiā bówùguǎn hěn yǒu yìsi, wǒ kànle hěn duō Fēizhōu yìshù. Nǐ ne?}\\
\zh{保罗：我和朋友去看电影了。晚上我们还听了现场音乐。}\\
\emph{Bǎoluó: Wǒ hé péngyou qù kàn diànyǐng le. Wǎnshang wǒmen hái tīngle xiànchǎng yīnyuè.}\\
\zh{王伟：真好！过春节的时候，中国人全家人一起吃饺子、看春晚、放烟花。你想不想去中国过春节？}\\
\emph{Wáng Wěi: Zhēn hǎo! Guò Chūnjié de shíhou, Zhōngguórén quánjiārén yìqǐ chī jiǎozi, kàn Chūnwǎn, fàng yānhuā. Nǐ xiǎng bu xiǎng qù Zhōngguó guò Chūnjié?}\\
\zh{保罗：很想！我对中国文化很感兴趣。}\\
\emph{Bǎoluó: Hěn xiǎng! Wǒ duì Zhōngguó wénhuà hěn gǎn xìngqù.}\\
« Paul : Wang Wei, es-tu allé au musée le week-end ? — Wang Wei : Oui. Le musée national de Yaoundé est très intéressant, j'ai vu beaucoup d'art africain. Et toi ? — Paul : Je suis allé au cinéma avec des amis. Le soir, on a aussi écouté de la musique live. — Wang Wei : Super ! Pendant la Fête du Printemps, les Chinois mangent des raviolis en famille, regardent le Gala du Nouvel An et lancent des feux d'artifice. Tu veux aller en Chine pour la Fête du Printemps ? — Paul : Très envie ! Je m'intéresse beaucoup à la culture chinoise. »
\end{textebox}

\begin{methode}[Entrer dans le thème : observer et classer]
\begin{enumerate}
\item \textbf{Identifie le thème.} Les mots \zh{文化} (culture), \zh{博物馆} (musée), \zh{电影} (cinéma), \zh{音乐} (musique), \zh{春节} (Fête du Printemps), \zh{烟花} (feux d'artifice) reviennent : le thème est la \cle{vie culturelle}.
\item \textbf{Repère les lieux et les actions.} Lieux : \zh{博物馆}、\zh{电影院} (implicite par \zh{看电影}). Actions : \zh{去} (aller), \zh{看} (voir/regarder), \zh{听} (écouter), \zh{吃} (manger), \zh{放} (lancer/tirer).
\item \textbf{Note les marqueurs culturels.} \zh{春节} (Chūnjié), \zh{饺子} (jiǎozi, raviolis), \zh{春晚} (Chūnwǎn, Gala du Nouvel An), \zh{烟花} (yānhuā) sont spécifiques à la Chine ; \zh{雅温得}、\zh{国家博物馆}、\zh{非洲艺术} ancrent le texte au Cameroun.
\end{enumerate}
\end{methode}

\courssub{Texte support : lecture et traduction}

\begin{textebox}[Le week-end de Marie — Texte support principal]
玛丽住在雅温得。她很喜欢喀麦隆的文化生活。周末，她常常和朋友一起去听音乐，有时候也去看电影。她最喜欢跳舞。她的中国朋友李华说，中国人过春节的时候，全家人一起吃饭、看电视，还放烟花。玛丽很想有一天去中国看看。\\
\emph{Mǎlì zhù zài Yǎwēndé. Tā hěn xǐhuan Kāmàilóng de wénhuà shēnghuó. Zhōumò, tā chángcháng hé péngyou yìqǐ qù tīng yīnyuè, yǒu shíhou yě qù kàn diànyǐng. Tā zuì xǐhuan tiàowǔ. Tā de Zhōngguó péngyou Lǐ Huá shuō, Zhōngguórén guò Chūnjié de shíhou, quánjiārén yìqǐ chīfàn, kàn diànshì, hái fàng yānhuā. Mǎlì hěn xiǎng yǒu yì tiān qù Zhōngguó kànkan.}\\
« Marie habite à Yaoundé. Elle aime beaucoup la vie culturelle camerounaise. Le week-end, elle va souvent écouter de la musique avec ses amis, et parfois elle va aussi au cinéma. Ce qu'elle préfère, c'est danser. Son amie chinoise Li Hua dit qu'au moment de la Fête du Printemps, les Chinois mangent en famille, regardent la télévision et lancent aussi des feux d'artifice. Marie aimerait beaucoup aller un jour visiter la Chine. »
\end{textebox}

\courssub{Glossaire du thème}

\begin{center}
\begin{tabularx}{\linewidth}{|X|X|X|X|}
\hline
\rowcolor{popblueL} \textbf{Caractères} & \textbf{Pinyin} & \textbf{Nature} & \textbf{Traduction} \\
\hline
文化 & wénhuà & n. & culture \\
生活 & shēnghuó & n. & vie, vie quotidienne \\
文化生活 & wénhuà shēnghuó & n. & vie culturelle \\
音乐 & yīnyuè & n. & musique \\
电影 & diànyǐng & n. & film, cinéma \\
电影院 & diànyǐngyuàn & n. & salle de cinéma \\
博物馆 & bówùguǎn & n. & musée \\
国家博物馆 & guójiā bówùguǎn & n. & musée national \\
节日 & jiérì & n. & fête, jour férié \\
春节 & Chūnjié & pn. & Fête du Printemps (Nouvel An chinois) \\
国庆节 & Guóqìngjié & pn. & Fête nationale (1er octobre) \\
舞蹈 & wǔdǎo & n. & danse (art) \\
跳舞 & tiàowǔ & v. & danser \\
唱歌 & chànggē & v. & chanter \\
听 & tīng & v. & écouter \\
看 & kàn & v. & voir, regarder, lire \\
过 & guò & v. & passer (fête), fêter \\
参观 & cānguān & v. & visiter (lieu) \\
喜欢 & xǐhuan & v. & aimer, apprécier \\
常常 & chángcháng & adv. & souvent \\
有时候 & yǒu shíhou & adv. & parfois \\
一起 & yìqǐ & adv. & ensemble \\
最 & zuì & adv. & le plus \\
还 & hái & adv. & en plus, aussi \\
想 & xiǎng & v. & vouloir, avoir envie de \\
因为 & yīnwèi & conj. & parce que \\
所以 & suǒyǐ & conj. & donc \\
饺子 & jiǎozi & n. & raviolis (chinois) \\
烟花 & yānhuā & n. & feux d'artifice \\
春晚 & Chūnwǎn & n. & Gala du Nouvel An (TV) \\
雅温得 & Yǎwēndé & pn. & Yaoundé \\
杜阿拉 & Dù'ālā & pn. & Douala \\
北京 & Běijīng & pn. & Pékin \\
喀麦隆 & Kāmàilóng & pn. & Cameroun \\
中国 & Zhōngguó & pn. & Chine \\
\hline
\end{tabularx}
\end{center}

\begin{popfigure}
\begin{tikzpicture}[mindmap, grow cyclic, every node/.style=concept, concept color=popblue, text=white,
    level 1/.append style={level distance=3.5cm, sibling angle=90, font=\small},
    level 2/.append style={level distance=3cm, sibling angle=45, font=\tiny}]
\node {Vie culturelle}
    child[concept color=popgreen] { node {Lieux (地点)}
        child { node {博物馆} }
        child { node {电影院} }
        child { node {公园 \emph{gōngyuán}} } }
    child[concept color=poporange] { node {Arts (艺术)}
        child { node {音乐} }
        child { node {舞蹈} }
        child { node {电影} } }
    child[concept color=poppurple] { node {Fêtes (节日)}
        child { node {春节} }
        child { node {国庆节} }
        child { node {节日} } }
    child[concept color=poppink] { node {Actions (活动)}
        child { node {听音乐} }
        child { node {看电影} }
        child { node {跳舞} }
        child { node {过节} }
        child { node {参观} } };
\end{tikzpicture}
\legende{Carte mentale : familles de mots du thème « vie culturelle »}
\end{popfigure}

\courssub{Savoir-faire 1 : Lire et comprendre un court texte chinois}

\begin{methode}[Lire et comprendre un texte sur la vie culturelle]
\begin{enumerate}
\item \textbf{Première lecture globale.} Lis le texte une fois sans dictionnaire. Repère le \cle{thème} (ici : la vie culturelle), les \cle{personnages} (Marie, Li Hua), le \cle{lieu} (Yaoundé, Chine) et le \cle{moment} (week-end, Fête du Printemps). Les noms propres (\zh{雅温得}、\zh{北京}) et les chiffres te guident.
\item \textbf{Repérage des mots connus.} Souligne les mots du vocabulaire du thème : \zh{文化}、\zh{音乐}、\zh{电影}、\zh{博物馆}、\zh{节日}、\zh{春节}、\zh{跳舞}、\zh{烟花}. Un texte de Première contient toujours 70 à 80 \% de mots déjà étudiés.
\item \textbf{Découpage des phrases.} Une phrase chinoise se découpe : Sujet + (complément de temps) + (complément de lieu avec \zh{在}) + verbe + complément. Repère d'abord le \cle{verbe}, puis remonte vers le sujet.
\item \textbf{Deuxième lecture, phrase par phrase.} Traduis mentalement chaque phrase. Pour un mot inconnu, devine grâce au \cle{contexte} et aux \cle{clés} (radicaux) : \zh{馆} contient la clé \zh{饣} (nourriture) à l'origine mais compose aujourd'hui des noms de lieux (\zh{图书馆}、\zh{博物馆}、\zh{饭馆}) ; \zh{影} a la clé \zh{彡} (traits décoratifs/lumière) → image, ombre.
\item \textbf{Vérification.} Relis le texte entier : ta compréhension doit être cohérente du début à la fin. Si une phrase contredit le reste, reprends-la.
\end{enumerate}
\end{methode}

\begin{exoresolu}[Compréhension globale d'un texte]
\textbf{Énoncé.} Lis le texte « Le week-end de Marie » ci-dessus et réponds en français : (a) Où habite Marie ? (b) Quelles activités culturelles aime-t-elle ? (c) Que font les Chinois pendant la Fête du Printemps ?
\tcblower
\textbf{Solution détaillée.}
\begin{enumerate}
\item \textbf{Question (a).} La première phrase donne la structure \zh{住在} + lieu : \zh{玛丽住在雅温得} = « Marie habite à Yaoundé ». Réponse : elle habite à \textbf{Yaoundé}.
\item \textbf{Question (b).} On repère les verbes d'activité : \zh{听音乐} (écouter de la musique), \zh{看电影} (voir des films), \zh{跳舞} (danser, introduit par \zh{最喜欢} « ce qu'elle préfère »). Réponse : elle aime écouter de la musique, aller au cinéma, et elle préfère danser.
\item \textbf{Question (c).} Le marqueur temporel \zh{过春节的时候} (« au moment où l'on fête le Nouvel An chinois ») introduit trois actions coordonnées : \zh{吃饭} (manger), \zh{看电视} (regarder la télé), \zh{还放烟花} (« et en plus » lancer des feux d'artifice). Réponse : ils mangent en famille, regardent la télévision (souvent le Gala \zh{春晚}) et lancent des feux d'artifice.
\end{enumerate}
\textbf{Conclusion.} Toutes les réponses se trouvent dans l'ordre du texte ; il suffit de repérer les verbes et les compléments de lieu et de temps.
\end{exoresolu}

\courssub{Savoir-faire 2 : Reconnaître, prononcer et employer le vocabulaire du thème}

\begin{methode}[Mémoriser et employer le vocabulaire de la vie culturelle]
\begin{enumerate}
\item \textbf{Apprends par familles de mots.} Regroupe : lieux (\zh{博物馆} musée, \zh{电影院} cinéma, \zh{公园} parc), arts (\zh{音乐} musique, \zh{舞蹈} danse, \zh{电影} film), fêtes (\zh{节日} fête, \zh{春节} Fête du Printemps, \zh{国庆节} fête nationale).
\item \textbf{Retiens les verbes avec leurs compléments habituels} (collocations) : \zh{听音乐} écouter de la musique, \zh{看电影} voir un film, \zh{跳舞} danser, \zh{唱歌} chanter, \zh{过节} / \zh{过春节} fêter une fête / le Nouvel An, \zh{参观博物馆} visiter un musée, \zh{放烟花} lancer des feux d'artifice, \zh{吃饺子} manger des raviolis.
\item \textbf{Prononce en respectant les tons.} Attention aux tons proches : \zh{wénhuà} (2\textsuperscript{e} + 4\textsuperscript{e} ton), \zh{yīnyuè} (1\textsuperscript{er} + 4\textsuperscript{e}), \zh{tiàowǔ} (4\textsuperscript{e} + 3\textsuperscript{e}), \zh{jiǎozi} (3\textsuperscript{e} + ton neutre), \zh{yānhuā} (1\textsuperscript{er} + 1\textsuperscript{er}). Répète chaque mot trois fois à voix haute.
\item \textbf{Construis une phrase complète} pour chaque mot : Sujet + \zh{喜欢/常常} + verbe + complément. Exemple : \zh{我常常听音乐。}
\item \textbf{Teste-toi} en couvrant la colonne des caractères, puis celle du pinyin, dans les deux sens (français → chinois et chinois → français).
\end{enumerate}
\end{methode}

\begin{exoresolu}[Vocabulaire : compléter et employer]
\textbf{Énoncé.} Complète les phrases avec le mot qui convient : \zh{博物馆}、\zh{电影}、\zh{春节}、\zh{音乐}、\zh{舞}。Puis traduis chaque phrase en français.

1. 星期天，我们去看了一个很好的\underline{\hspace{1.5cm}} \quad 2. 中国人过\underline{\hspace{1.5cm}}的时候吃饺子。 \quad 3. 我喜欢听喀麦隆的\underline{\hspace{1.5cm}}。 \quad 4. 雅温得的国家\underline{\hspace{1.5cm}}很有意思。 \quad 5. 晚会的时候，大家又唱歌又跳\underline{\hspace{1.5cm}}。
\tcblower
\textbf{Solution détaillée.}
\begin{enumerate}
\item \zh{电影} : \zh{看} + film ; la collocation est \zh{看电影}. Phrase : \zh{星期天，我们去看了一个很好的电影。} \emph{Xīngqītiān, wǒmen qù kànle yí ge hěn hǎo de diànyǐng.} « Dimanche, nous sommes allés voir un très bon film. » (Classificateur \zh{个} devant \zh{电影} ; \zh{了} marque l'accompli.)
\item \zh{春节} : la collocation est \zh{过春节} « fêter le Nouvel An chinois » ; manger des raviolis (\zh{饺子}) est la tradition du Nord de la Chine. \zh{中国人过春节的时候吃饺子。} \emph{Zhōngguórén guò Chūnjié de shíhou chī jiǎozi.} « Les Chinois mangent des raviolis pendant la Fête du Printemps. »
\item \zh{音乐} : collocation \zh{听音乐}. \zh{我喜欢听喀麦隆的音乐。} \emph{Wǒ xǐhuan tīng Kāmàilóng de yīnyuè.} « J'aime écouter la musique camerounaise. »
\item \zh{博物馆} : \zh{国家博物馆} « musée national ». \zh{雅温得的国家博物馆很有意思。} \emph{Yǎwēndé de guójiā bówùguǎn hěn yǒu yìsi.} « Le musée national de Yaoundé est très intéressant. »
\item \zh{舞} : structure \zh{又……又……} « à la fois … et … » ; \zh{跳舞} se coupe ici en \zh{跳} + \zh{舞} (nom verbal). \zh{晚会的时候，大家又唱歌又跳舞。} \emph{Wǎnhuì de shíhou, dàjiā yòu chànggē yòu tiàowǔ.} « Pendant la soirée, tout le monde chante et danse. »
\end{enumerate}
\textbf{Conclusion.} Chaque réponse repose sur une collocation fixe : c'est le verbe qui appelle son complément habituel.
\end{exoresolu}

\courssub{Savoir-faire 3 : Répondre à des questions de compréhension}

\begin{methode}[Répondre à une question de compréhension en chinois]
\begin{enumerate}
\item \textbf{Analyse la question.} Repère le mot interrogatif : \zh{谁} (qui), \zh{什么} (quoi), \zh{哪儿/哪里} (où), \zh{什么时候} (quand), \zh{为什么} (pourquoi), \zh{怎么样} (comment), \zh{吗} (question fermée).
\item \textbf{Localise la phrase du texte} qui contient la réponse : le mot interrogatif te dit quel élément chercher (personne, lieu, temps, raison).
\item \textbf{Rédige une phrase complète}, pas un mot isolé : reprends la structure de la question en remplaçant le mot interrogatif par la réponse. Question : \zh{玛丽住在哪儿？} → Réponse : \zh{玛丽住在雅温得。}
\item \textbf{Pour \zh{为什么},} réponds avec \zh{因为……（所以）……} : \zh{因为她很喜欢中国文化。}
\item \textbf{Pour une question en \zh{吗},} commence par \zh{对/是的} (oui) ou \zh{不} + phrase négative, puis justifie.
\item \textbf{Vérifie} : ordre des mots (Sujet - Temps - Lieu - Verbe - Complément), tons du pinyin si demandé, ponctuation chinoise 。
\end{enumerate}
\end{methode}

\begin{exoresolu}[Questions de compréhension sur le texte « Le week-end de Marie »]
\textbf{Énoncé.} Réponds en chinois, par des phrases complètes, d'après le texte du savoir-faire 1.

1. 玛丽周末常常做什么？ \quad 2. 她最喜欢什么活动？ \quad 3. 中国人过春节的时候做什么？ \quad 4. 玛丽想去中国吗？为什么？
\tcblower
\textbf{Solution détaillée.}
\begin{enumerate}
\item La question porte sur \zh{什么} (quoi) + fréquence \zh{常常}. Le texte dit : \zh{周末，她常常和朋友一起去听音乐}. Réponse rédigée : \zh{她周末常常和朋友一起去听音乐，有时候也去看电影。} \emph{Tā zhōumò chángcháng hé péngyou yìqǐ qù tīng yīnyuè, yǒu shíhou yě qù kàn diànyǐng.} « Le week-end, elle va souvent écouter de la musique avec ses amis, et parfois elle va aussi au cinéma. »
\item Le marqueur \zh{最喜欢} désigne la préférence : \zh{她最喜欢跳舞。} \emph{Tā zuì xǐhuan tiàowǔ.} « Ce qu'elle préfère, c'est danser. »
\item Question plurielle : on reprend les trois actions avec \zh{还} : \zh{他们全家人一起吃饭、看电视，还放烟花。} \emph{Tāmen quánjiārén yìqǐ chīfàn, kàn diànshì, hái fàng yānhuā.} « Ils mangent en famille, regardent la télévision et lancent aussi des feux d'artifice. »
\item Question fermée en \zh{吗} + \zh{为什么} : le texte dit \zh{玛丽很想有一天去中国看看}. Réponse : \zh{对，她很想去中国，因为她对中国的文化生活很感兴趣。} \emph{Duì, tā hěn xiǎng qù Zhōngguó, yīnwèi tā duì Zhōngguó de wénhuà shēnghuó hěn gǎn xìngqù.} « Oui, elle a très envie d'aller en Chine, parce que la vie culturelle chinoise l'intéresse beaucoup. » (Structure à retenir : \zh{对……感兴趣} « s'intéresser à ».)
\end{enumerate}
\textbf{Conclusion.} Chaque réponse reprend la structure de la question et remplace le mot interrogatif par l'information du texte : c'est la méthode la plus sûre pour ne pas perdre de points.
\end{exoresolu}

\courssub{Savoir-faire 4 : Reconnaître les clés et l'ordre des traits}

\begin{methode}[Analyser un caractère : clé et ordre des traits]
\begin{enumerate}
\item \textbf{Identifie la clé (radical)} : elle donne un indice de sens. Dans le thème culturel : \zh{讠} (parole) dans \zh{话}, \zh{艹} (herbe) dans \zh{节}, \zh{宀} (toit) dans \zh{家}, \zh{日} (soleil) dans \zh{春}, \zh{彡} (traits/lumière) dans \zh{影}, \zh{罒/四} (structure) dans \zh{舞}.
\item \textbf{Applique les règles d'ordre des traits} : de haut en bas (\zh{三}), de gauche à droite (\zh{川}), l'horizontal avant le vertical (\zh{十}), l'extérieur avant l'intérieur (\zh{国}), on ferme la boîte en dernier (\zh{口} → \zh{日} → \zh{田}), le trait central avant les côtés symétriques (\zh{小}、\zh{水}).
\item \textbf{Compte les traits} en écrivant lentement : \zh{文} 4 traits, \zh{化} 4 traits, \zh{节} 5 traits, \zh{春} 9 traits, \zh{影} 15 traits.
\item \textbf{Distingue les caractères proches} (pièges visuels) : \zh{文} (wén, culture) ≠ \zh{又} (yòu, encore) ; \zh{化} (huà, transformer) ≠ \zh{北} (běi, nord) ; \zh{馆} (guǎn, lieu public) ≠ \zh{管} (guǎn, tube / gérer) ; \zh{影} (yǐng, ombre) ≠ \zh{景} (jǐng, paysage).
\item \textbf{Écris chaque caractère 5 fois} en nommant la clé à voix haute : « \zh{话}, clé de la parole \zh{讠}, donc lié au langage. »
\end{enumerate}
\end{methode}

\begin{exoresolu}[Clés et ordre des traits]
\textbf{Énoncé.} (a) Donne la clé de chaque caractère et explique le lien avec le sens : \zh{话}、\zh{节}、\zh{家}、\zh{影}。(b) Décris l'ordre des traits de \zh{文} et \zh{春}。(c) Combien de traits a \zh{文} ?
\tcblower
\textbf{Solution détaillée.}
\begin{enumerate}
\item[(a)] \zh{话} (huà, parole/phrase) : clé \textbf{讠} (parole, simplifié de \zh{言}) à gauche → tout caractère avec \zh{讠} concerne le langage (\zh{说} shuō, \zh{语} yǔ, \zh{读} dú, \zh{讨论} tǎolùn).\\
\zh{节} (jié, fête / nœud / section) : clé \textbf{艹} (herbe, simplifié de \zh{艸}) en haut → sens ancien lié aux plantes (nœud de bambou \zh{节}, puis par extension « fête », « section », « économiser »).\\
\zh{家} (jiā, famille / maison / foyer) : clé \textbf{宀} (toit) → un cochon \zh{豕} (shǐ) sous un toit = la maison, le foyer, la famille.\\
\zh{影} (yǐng, ombre / image / reflet) : clé \textbf{彡} (shān, traits décoratifs, cheveux, lumière) à droite → \zh{电影} « images électriques » = cinéma ; \zh{影子} ombre.
\item[(b)] \zh{文} (wén, 4 traits) : 1) le point \zh{丶} en haut, 2) l'horizontal \zh{一}, 3) le trait descendant gauche \zh{丿} (piě), 4) le trait descendant droit \zh{㇏} (nà) — les deux derniers se croisent, gauche avant droite.\\
\zh{春} (chūn, 9 traits) : de haut en bas : les trois horizontaux \zh{三} (trois traits), puis le \zh{丿} (4\textsuperscript{e}) et le \zh{㇏} (5\textsuperscript{e}) qui les traversent, puis le composant \zh{日} (soleil, 4 traits : vertical, pli, deux horizontaux, fermeture) en bas — règle : extérieur avant intérieur, fermeture en dernier.
\item[(c)] \zh{文} a \textbf{4 traits}.
\end{enumerate}
\textbf{Conclusion.} Nommer la clé aide à retenir le sens ; respecter l'ordre des traits garantit une écriture lisible, rapide et standardisée, exigée à l'examen.
\end{exoresolu}

\courssub{Savoir-faire 5 : Employer les collocations et produire des phrases}

\begin{methode}[Construire des phrases correctes avec les collocations du thème]
\begin{enumerate}
\item \textbf{Verbe + complément inséparable} : ne traduis jamais mot à mot. « Écouter de la musique » = \zh{听音乐} ; « voir un film » = \zh{看电影} ; « fêter une fête » = \zh{过节} / \zh{过春节} ; « visiter un musée » = \zh{参观博物馆}.
\item \textbf{Place des compléments} : le complément de temps va \emph{avant} le verbe (après le sujet) : \zh{我周末去看电影。} — jamais à la fin comme en français. Le lieu avec \zh{在} précède aussi le verbe : \zh{我在雅温得学习。}
\item \textbf{Fréquence} : \zh{常常} / \zh{经常} / \zh{有时候} se placent avant le verbe : \zh{她常常听音乐。}
\item \textbf{Coordonner deux activités} : \zh{又……又……} (simultané, même sujet) : \zh{我们又唱歌又跳舞。} ; \zh{……还……} (addition, séquentiel ou emphasis) : \zh{我们吃饭，还看电视。}
\item \textbf{Exprimer l'envie et la préférence} : \zh{想} + verbe : \zh{我想去中国看看。} ; \zh{最喜欢} + verbe : \zh{我最喜欢跳舞。} ; \zh{对……感兴趣} : \zh{我对中国文化很感兴趣。}
\item \textbf{Rédaction guidée (4-6 phrases)} : 1. Présentation (nom, ville : \zh{住在}). 2. Activités fréquentes (temps + \zh{常常} + verbe). 3. Préférence (\zh{最喜欢}). 4. Fête du pays (\zh{过} + fête). 5. Projet (\zh{想} + aller + \zh{因为}). Relis en vérifiant ordre des mots, tons, \zh{的} (possession), classificateurs.
\end{enumerate}
\end{methode}

\begin{exoresolu}[Production guidée : ma vie culturelle]
\textbf{Énoncé.} Écris 5 phrases en chinois (avec pinyin et traduction) pour présenter ta vie culturelle : où tu habites, deux activités que tu fais souvent, ton activité préférée, une fête de ton pays, et un pays que tu veux visiter.
\tcblower
\textbf{Solution détaillée (modèle de rédaction).}

\zh{我住在杜阿拉。} \emph{Wǒ zhù zài Dù'ālā.} « J'habite à Douala. » (Structure : sujet + \zh{住在} + ville.)

\zh{周末，我常常和朋友一起听音乐，有时候也去看足球比赛。} \emph{Zhōumò, wǒ chángcháng hé péngyou yìqǐ tīng yīnyuè, yǒu shíhou yě qù kàn zúqiú bǐsài.} « Le week-end, j'écoute souvent de la musique avec mes amis, et parfois je vais aussi voir des matchs de football. » (Temps avant le verbe ; \zh{也} avant le verbe ; \zh{足球比赛} match de foot.)

\zh{我最喜欢跳舞。} \emph{Wǒ zuì xǐhuan tiàowǔ.} « Ce que je préfère, c'est danser. »

\zh{我们国家有很多节日，大家都很喜欢过节。} \emph{Wǒmen guójiā yǒu hěn duō jiérì, dàjiā dōu hěn xǐhuan guò jié.} « Notre pays a beaucoup de fêtes, et tout le monde aime les fêter. » (\zh{都} se place avant le verbe ; \zh{过节} collocation.)

\zh{我想有一天去中国看看，因为我对中国文化很感兴趣。} \emph{Wǒ xiǎng yǒu yì tiān qù Zhōngguó kànkan, yīnwèi wǒ duì Zhōngguó wénhuà hěn gǎn xìngqù.} « J'aimerais aller un jour visiter la Chine, parce que je m'intéresse beaucoup à la culture chinoise. » (\zh{想} + verbe ; \zh{因为} pour la justification ; \zh{看看} verbe redoublé « aller voir un peu ».)

\textbf{Conclusion.} Cinq phrases suffisent si chacune respecte : sujet → temps → lieu → verbe → complément, avec des collocations correctes et les mots-outils (\zh{的}、\zh{了}、\zh{过}、\zh{在}).
\end{exoresolu}

\courssub{Ouverture interculturelle : Cameroun – Chine}

\begin{savaistu}[Vie culturelle : points de rencontre et différences]
\begin{itemize}
\item \textbf{Fêtes nationales.} Cameroun : \zh{国庆节} (20 mai, fête de l'unité), \zh{青年节} (11 février). Chine : \zh{国庆节} (1er octobre, fondation de la RPC), \zh{春节} (Nouvel An lunaire, date variable janv./févr.), \zh{中秋节} (Fête de la Mi-Automne).
\item \textbf{Repas de fête.} Cameroun : \zh{恩多莱} (ndolé), \zh{香蕉炖鱼} (poisson braisé/banane plantain), \zh{基辛} (koki). Chine Nord : \zh{饺子} (raviolis) ; Chine Sud : \zh{年糕} (gâteau de riz gluant), \zh{汤圆} (boulettes de riz gluant).
\item \textbf{Arts populaires.} Cameroun : \zh{马卡萨} (makossa), \zh{本皮} (bikutsi), danses traditionnelles (Bassa, Bamiléké), \zh{足球} (football, passion nationale). Chine : \zh{京剧} (opéra de Pékin), \zh{太极拳} (tai-chi), \zh{书法} (calligraphie), \zh{乒乓球} (tennis de table, sport national).
\item \textbf{Lieux culturels.} Yaoundé : Musée national, Musée d'art camerounais, Centre culturel français. Pékin : Cité interdite (\zh{故宫}), Temple du Ciel (\zh{天坛}), Musée national de Chine (\zh{中国国家博物馆}), Opéra de Pékin (\zh{国家大剧院}).
\item \textbf{Valeurs partagées.} Respect des aînés (\zh{尊老爱幼}), importance de la famille (\zh{家庭}), hospitalité (\zh{好客}), attachement aux traditions (\zh{传统}) tout en s'ouvrant à la modernité (\zh{现代化}).
\end{itemize}
\end{savaistu}

\begin{experience}[Activité orale : « Mon week-end culturel idéal »]
\textbf{Consigne.} En binôme (A et B). A : Tu es un élève camerounais. B : Tu es un élève chinois en échange. Échangez 3 minutes sur vos activités de week-end préférées.
\begin{enumerate}
\item A demande : \zh{你周末喜欢做什么？} B répond avec \zh{常常}、\zh{有时候}、\zh{最喜欢}.
\item B demande : \zh{喀麦隆有什么节日？怎么过？} A répond avec \zh{过节}、\zh{吃}、\zh{跳舞}、\zh{足球}.
\item Ensemble : Trouvez 2 points communs et 2 différences. Présentez à la classe en chinois.
\end{enumerate}
\textbf{Support lexicale :} \zh{一样} (pareil), \zh{不一样} (différent), \zh{都} (tous les deux), \zh{但是} (mais).
\end{experience}

\begin{attention}[Les 5 erreurs qui coûtent des points au Bac]
\begin{enumerate}
\item \textbf{Tons négligés.} Confondre \zh{wén} (2\textsuperscript{e} ton, dans \zh{文化}) avec wěn (3\textsuperscript{e} ton, stable) ou wèn (4\textsuperscript{e} ton, demander), ou écrire le pinyin sans tons : à l'examen, un pinyin sans ton est un pinyin faux. Révise : \zh{yīnyuè} (musique) n'est pas \zh{yīnyuē} ; \zh{jiǎozi} (raviolis) ton 3 + neutre, pas \zh{jiāozi}.
\item \textbf{Caractères confondus (pièges visuels).} \zh{文} (wén, culture) ≠ \zh{又} (yòu, encore) ; \zh{化} (huà) ≠ \zh{北} (běi) ; \zh{馆} (guǎn, lieu : musée, restaurant, ambassade) ≠ \zh{管} (guǎn, tube / gérer) ; \zh{影} (yǐng, image/ombre) ≠ \zh{景} (jǐng, paysage). Écris chaque caractère en respectant l'ordre des traits : un trait mal placé peut créer un autre caractère.
\item \textbf{Ordre des mots calqué du français.} Ne dis jamais \zh{我看电影周末} : le complément de temps se place \emph{avant} le verbe : \zh{我周末看电影。} De même, le lieu avec \zh{在} précède le verbe : \zh{我在雅温得学习。} (Sujet + Temps + Lieu + Verbe).
\item \textbf{Collocations inventées (verbe + objet).} On ne peut pas dire \zh{看音乐} (regarder la musique) : la musique s'\emph{écoute} (\zh{听音乐}), le film se \emph{voit} (\zh{看电影}), la fête se \emph{fête} (\zh{过节}), le musée se \emph{visite} (\zh{参观博物馆}). Apprends toujours le verbe \textbf{avec} son complément habituel.
\item \textbf{Oubli des mots-outils et classificateurs.} Un film = \zh{一个电影} (classificateur \zh{个}) ; une chanson = \zh{一首歌} (classificateur \zh{首}) ; la possession exige \zh{的} : \zh{喀麦隆的文化} (la culture \emph{du} Cameroun) ; l'aspect accompli \zh{了} après le verbe : \zh{我看了电影} ; l'expérience \zh{过} : \zh{我去过中国} ; la coordination de deux qualités/actions simultanées exige \zh{又……又……}, pas un simple « et » (\zh{和} ne lie que des noms).
\end{enumerate}
\end{attention}

\newpage

\coursec{Exercices}

\courssub{Je vérifie mes connaissances}

\begin{exercice}[QCM : le vocabulaire de la vie culturelle]
Pour chaque question, entoure la bonne réponse.

\begin{enumerate}
\item \zh{节日} (jiérì) signifie :
\begin{enumerate}
\item la musique \item la fête \item le vêtement \item la danse
\end{enumerate}
\item Le mot qui signifie « célébrer, fêter » est :
\begin{enumerate}
\item \zh{传统} \item \zh{文化} \item \zh{庆祝} \item \zh{互相}
\end{enumerate}
\item \zh{表演节目} (biǎoyǎn jiémù) signifie :
\begin{enumerate}
\item chanter une chanson \item regarder la télévision \item présenter un spectacle \item porter un habit
\end{enumerate}
\item Quel mot désigne les « vêtements traditionnels » ?
\begin{enumerate}
\item \zh{现代服装} \item \zh{传统服装} \item \zh{新衣服} \item \zh{工作服}
\end{enumerate}
\end{enumerate}
\end{exercice}

\begin{exercice}[Vrai ou faux, justifié]
Lis le texte support de la leçon, puis dis si chaque phrase est vraie (\zh{对} duì) ou fausse (\zh{不对} bú duì). Justifie ta réponse en citant une phrase du texte en chinois.

\begin{enumerate}
\item \zh{喀麦隆人不喜欢音乐和舞蹈。} (Kāmàilóng rén bù xǐhuan yīnyuè hé wǔdǎo.)
\item \zh{春节是中国人最重要的传统节日。} (Chūnjié shì Zhōngguó rén zuì zhòngyào de chuántǒng jiérì.)
\item \zh{过春节的时候，中国人不穿新衣服。} (Guò Chūnjié de shíhou, Zhōngguó rén bù chuān xīn yīfu.)
\item \zh{喀麦隆和中国都有丰富的文化。} (Kāmàilóng hé Zhōngguó dōu yǒu fēngfù de wénhuà.)
\end{enumerate}
\end{exercice}

\begin{exercice}[Phrases à trous avec la banque de mots]
Complète les phrases avec les mots suivants : \zh{文化} (wénhuà), \zh{节日} (jiérì), \zh{传统} (chuántǒng), \zh{庆祝} (qìngzhù), \zh{表演} (biǎoyǎn), \zh{互相} (hùxiāng).

\begin{enumerate}
\item \zh{春节是中国最重要的\_\_\_\_\_\_。}
\item \zh{过节日的时候，我们\_\_\_\_\_\_说« 新年快乐 »。}
\item \zh{学生们在学校里\_\_\_\_\_\_节目。}
\item \zh{喀麦隆的\_\_\_\_\_\_很丰富：有音乐、舞蹈和美食。}
\item \zh{我们全家一起吃饭，\_\_\_\_\_\_这个快乐的日子。}
\item \zh{跳舞的时候，人们常常穿\_\_\_\_\_\_服装。}
\end{enumerate}
\end{exercice}

\begin{exercice}[Associe le pinyin aux caractères]
Relie chaque caractère ou mot à son pinyin.

\begin{center}
\begin{tabularx}{\linewidth}{|X|X|}
\hline
\rowcolor{popblueL} Caractères & Pinyin \\
\hline
\zh{舞蹈} & chuántǒng \\
\hline
\zh{传统} & qìngzhù \\
\hline
\zh{庆祝} & yīnyuè \\
\hline
\zh{音乐} & wǔdǎo \\
\hline
\zh{祝福} & zhùfú \\
\hline
\end{tabularx}
\end{center}
\end{exercice}

\courssub{Je m'entraîne}

\begin{exercice}[Traduis en français]
Traduis les phrases suivantes en français correct.

\begin{enumerate}
\item \zh{喀麦隆人非常喜欢音乐和舞蹈。}\\ (Kāmàilóng rén fēicháng xǐhuan yīnyuè hé wǔdǎo.)
\item \zh{过春节的时候，中国人穿新衣服，互相祝福。}\\ (Guò Chūnjié de shíhou, Zhōngguó rén chuān xīn yīfu, hùxiāng zhùfú.)
\item \zh{我们学校的学生表演了传统节目。}\\ (Wǒmen xuéxiào de xuésheng biǎoyǎn le chuántǒng jiémù.)
\item \zh{中国和喀麦隆都有丰富的文化。}\\ (Zhōngguó hé Kāmàilóng dōu yǒu fēngfù de wénhuà.)
\end{enumerate}
\end{exercice}

\begin{exercice}[Appariement de collocations]
Associe chaque verbe de la colonne A au complément de la colonne B, puis traduis chaque collocation en français.

\begin{center}
\begin{tabularx}{\linewidth}{|X|X|X|}
\hline
\rowcolor{popblueL} A : verbe & B : complément & Traduction \\
\hline
\zh{过} & \zh{传统服装} & \ldots \\
\hline
\zh{穿} & \zh{节日} & \ldots \\
\hline
\zh{表演} & \zh{歌} & \ldots \\
\hline
\zh{唱} & \zh{节目} & \ldots \\
\hline
\end{tabularx}
\end{center}
\end{exercice}

\begin{textebox}[Le Nouvel An chinois — texte à trous]
春节是中国人\_\_\_\_\_\_重要\_\_\_\_\_\_节日。\\
(Chūnjié shì Zhōngguó rén \_\_\_\_\_\_ zhòngyào \_\_\_\_\_\_ jiérì.)\\
过春节的时候，人们穿新衣服，\_\_\_\_\_\_祝福。\\
(Guò Chūnjié de shíhou, rénmen chuān xīn yīfu, \_\_\_\_\_\_ zhùfú.)\\
我们全家一起吃年夜饭，还看\_\_\_\_\_\_很多节目。\\
(Wǒmen quánjiā yìqǐ chī niányèfàn, hái kàn \_\_\_\_\_\_ hěn duō jiémù.)\\
唱歌\_\_\_\_\_\_跳舞都是春节的\_\_\_\_\_\_活动。\\
(Chàng gē \_\_\_\_\_\_ tiào wǔ dōu shì Chūnjié de \_\_\_\_\_\_ huódòng.)
\end{textebox}

\begin{exercice}[Texte à trous : caractères et mots-outils]
Complète le texte « Le Nouvel An chinois » ci-dessus avec les mots de la banque : \zh{的} (de), \zh{了} (le), \zh{和} (hé), \zh{最} (zuì), \zh{互相} (hùxiāng), \zh{传统} (chuántǒng). Chaque mot n'est utilisé qu'une seule fois.
\end{exercice}

\begin{exercice}[Remets les mots dans l'ordre]
Remets les mots dans le bon ordre pour faire des phrases correctes, puis traduis-les en français.

\begin{enumerate}
\item \zh{喜欢 / 喀麦隆人 / 音乐 / 非常 / 和 / 舞蹈}
\item \zh{春节 / 是 / 节日 / 中国 / 的 / 传统 / 最 / 重要}
\item \zh{穿 / 人们 / 时候 / 过节日的 / 传统服装 / 漂亮的}
\item \zh{表演 / 学生们 / 了 / 节目 / 精彩的}
\end{enumerate}
\end{exercice}

\courssub{Je me prépare au Bac}

\begin{textebox}[中国的春节]
春节是中国人最重要的传统节日。\\
(Chūnjié shì Zhōngguó rén zuì zhòngyào de chuántǒng jiérì.)\\
过春节以前，人们打扫房子，买很多好吃的东西。\\
(Guò Chūnjié yǐqián, rénmen dǎsǎo fángzi, mǎi hěn duō hǎochī de dōngxi.)\\
除夕晚上，全家一起吃年夜饭，互相祝福。\\
(Chúxī wǎnshang, quánjiā yìqǐ chī niányèfàn, hùxiāng zhùfú.)\\
孩子们最喜欢春节，因为他们可以穿新衣服，还可以拿到红包。\\
(Háizimen zuì xǐhuan Chūnjié, yīnwèi tāmen kěyǐ chuān xīn yīfu, hái kěyǐ ná dào hóngbāo.)\\
节日的时候，很多地方都有舞龙和舞狮表演，非常热闹。\\
(Jiérì de shíhou, hěn duō dìfang dōu yǒu wǔlóng hé wǔshī biǎoyǎn, fēicháng rènao.)\\
春节让人们感受到中国文化的丰富和美丽。\\
(Chūnjié ràng rénmen gǎnshòu dào Zhōngguó wénhuà de fēngfù hé měilì.)
\end{textebox}

\begin{exercice}[Compréhension écrite type Bac : la fête du Printemps]
Lis le texte « 中国的春节 » ci-dessus, puis réponds en chinois, par des phrases complètes, aux cinq questions.

Questions :
\begin{enumerate}
\item \zh{春节是什么节日？}
\item \zh{过春节以前，人们做什么？}
\item \zh{除夕晚上，全家一起做什么？}
\item \zh{孩子们为什么最喜欢春节？}
\item \zh{节日的时候，人们可以看什么表演？这些表演怎么样？}
\end{enumerate}

Question de langue : relève dans le texte une phrase avec \zh{因为} et traduis-la en français ; puis fais une phrase personnelle avec \zh{因为}.
\end{exercice}

\begin{exercice}[Thème et version]
\textbf{Version.} Traduis en chinois (caractères simplifiés) le texte suivant :

« Les Camerounais aiment beaucoup la musique et la danse. Pendant les fêtes, les gens portent de beaux vêtements traditionnels. Les familles se retrouvent, mangent ensemble et se souhaitent mutuellement le bonheur. »

\textbf{Thème.} Traduis en français le texte « À traduire » ci-dessous.
\end{exercice}

\begin{textebox}[À traduire]
中国的文化很丰富。\\
(Zhōngguó de wénhuà hěn fēngfù.)\\
过节日的时候，人们唱歌、跳舞、表演节目。\\
(Guò jiérì de shíhou, rénmen chàng gē, tiào wǔ, biǎoyǎn jiémù.)\\
大家都非常高兴。\\
(Dàjiā dōu fēicháng gāoxìng.)
\end{textebox}

\begin{exercice}[Expression écrite type Bac]
Sujet : \zh{介绍一个喀麦隆的传统节日} — Présente une fête traditionnelle du Cameroun (par exemple une fête de ton village ou de ta région).

Consignes :
\begin{enumerate}
\item Rédige un court texte en chinois simplifié de \textbf{60 à 80 caractères}.
\item Tu dois utiliser au moins cinq mots du vocabulaire de la leçon : \zh{文化}, \zh{节日}, \zh{传统}, \zh{庆祝}, \zh{表演}, \zh{互相}, \zh{服装}, \zh{音乐}, \zh{舞蹈}.
\item Emploie au moins une structure avec \zh{的时候} et une structure avec \zh{互相}.
\item Soigne l'ordre des mots et les tons du pinyin si tu recopies ton texte en pinyin.
\end{enumerate}

Aide au démarrage : \zh{在我的家乡，有一个很重要的节日……} (Zài wǒ de jiāxiāng, yǒu yí ge hěn zhòngyào de jiérì...)
\end{exercice}



\newpage

\coursec{Corrigés des exercices}

\begin{corrige}[Exercice 1 — QCM : le vocabulaire de la vie culturelle]
\begin{enumerate}
\item Réponse \textbf{b} : \zh{节日} (jiérì) = la fête, le jour de fête. \textbf{Justification} : \zh{节} signifie « fête, nœud » et \zh{日} « jour » ; littéralement « jour de fête ».
\item Réponse \textbf{c} : \zh{庆祝} (qìngzhù) = célébrer, fêter. \textbf{Distracteurs} : \zh{传统} (chuántǒng) = tradition ; \zh{文化} (wénhuà) = culture ; \zh{互相} (hùxiāng) = mutuellement.
\item Réponse \textbf{c} : \zh{表演节目} (biǎoyǎn jiémù) = présenter un spectacle, un numéro. \textbf{Justification} : \zh{表演} est le verbe « représenter, jouer » et \zh{节目} « programme, numéro ».
\item Réponse \textbf{b} : \zh{传统服装} (chuántǒng fúzhuāng) = les vêtements traditionnels. \textbf{Distracteurs} : \zh{现代服装} = vêtements modernes ; \zh{新衣服} = vêtements neufs ; \zh{工作服} = tenue de travail.
\end{enumerate}
\textbf{Point de vigilance} : ne pas confondre \zh{节日} (la fête) et \zh{节目} (le programme, le spectacle) ; les deux contiennent \zh{节} mais le second caractère diffère (\zh{日} « jour » contre \zh{目} « œil »).
\end{corrige}

\begin{corrige}[Exercice 2 — Vrai ou faux, justifié]
\begin{enumerate}
\item \textbf{不对} (faux). Le texte dit : \zh{喀麦隆人非常喜欢音乐和舞蹈。} (Kāmàilóng rén fēicháng xǐhuan yīnyuè hé wǔdǎo.) « Les Camerounais aiment beaucoup la musique et la danse. » La phrase de l'énoncé contient la négation \zh{不}, contraire au texte.
\item \textbf{对} (vrai). Le texte dit : \zh{春节是中国人最重要的传统节日。} (Chūnjié shì Zhōngguó rén zuì zhòngyào de chuántǒng jiérì.)
\item \textbf{不对} (faux). Le texte dit : \zh{过春节的时候，中国人穿新衣服，互相祝福。} « Pour la fête du Printemps, les Chinois portent des vêtements neufs et échangent leurs vœux. »
\item \textbf{对} (vrai). Le texte dit : \zh{中国和喀麦隆都有丰富的文化。} (Zhōngguó hé Kāmàilóng dōu yǒu fēngfù de wénhuà.)
\end{enumerate}
\textbf{Point de vigilance} : la justification doit citer une phrase du texte en chinois, pas seulement une traduction française ; c'est la méthode attendue en compréhension écrite.
\end{corrige}

\begin{corrige}[Exercice 3 — Phrases à trous avec la banque de mots]
\begin{enumerate}
\item \zh{春节是中国最重要的节日。} (Chūnjié shì Zhōngguó zuì zhòngyào de jiérì.) « La fête du Printemps est la fête la plus importante de Chine. »
\item \zh{过节日的时候，我们互相说“新年快乐”。} (Guò jiérì de shíhou, wǒmen hùxiāng shuō “Xīnnián kuàilè”.) « Pendant les fêtes, nous nous souhaitons mutuellement "Bonne année". » \textbf{Justification} : \zh{互相} se place toujours avant le verbe.
\item \zh{学生们在学校里表演节目。} (Xuéshengmen zài xuéxiào lǐ biǎoyǎn jiémù.) « Les élèves présentent des spectacles à l'école. »
\item \zh{喀麦隆的文化很丰富：有音乐、舞蹈和美食。} (Kāmàilóng de wénhuà hěn fēngfù: yǒu yīnyuè, wǔdǎo hé měishí.) « La culture du Cameroun est très riche : il y a la musique, la danse et la gastronomie. »
\item \zh{我们全家一起吃饭，庆祝这个快乐的日子。} (Wǒmen quánjiā yíqǐ chīfàn, qìngzhù zhège kuàilè de rìzi.) « Toute notre famille mange ensemble et célèbre ce jour heureux. »
\item \zh{跳舞的时候，人们常常穿传统服装。} (Tiàowǔ de shíhou, rénmen chángcháng chuān chuántǒng fúzhuāng.) « Quand on danse, les gens portent souvent des vêtements traditionnels. »
\end{enumerate}
\textbf{Point de vigilance} : chaque mot de la banque n'est utilisé qu'une fois ; vérifier que la phrase complétée reste grammaticale (sujet + verbe + complément).
\end{corrige}

\begin{corrige}[Exercice 4 — Associe le pinyin aux caractères]
\begin{itemize}
\item \zh{舞蹈} — wǔdǎo (la danse)
\item \zh{传统} — chuántǒng (la tradition)
\item \zh{庆祝} — qìngzhù (célébrer)
\item \zh{音乐} — yīnyuè (la musique)
\item \zh{祝福} — zhùfú (souhaiter, bénir ; les vœux)
\end{itemize}
\textbf{Point de vigilance} : attention aux tons de \zh{传统} (chuántǒng, 2ᵉ puis 3ᵉ ton) et de \zh{庆祝} (qìngzhù, deux 4ᵉ tons) ; ne pas confondre \zh{音乐} yīnyuè avec \zh{月} yuè « lune » seul.
\end{corrige}

\begin{corrige}[Exercice 5 — Traduis en français]
\begin{enumerate}
\item \zh{喀麦隆人非常喜欢音乐和舞蹈。} : « Les Camerounais aiment beaucoup la musique et la danse. » \textbf{Justification} : \zh{非常} = « très, beaucoup », adverbe placé avant le verbe \zh{喜欢}.
\item \zh{过春节的时候，中国人穿新衣服，互相祝福。} : « Au moment de la fête du Printemps, les Chinois portent des vêtements neufs et se souhaitent mutuellement leurs vœux. » \textbf{Justification} : \zh{……的时候} introduit le complément de temps « au moment de… ».
\item \zh{我们学校的学生表演了传统节目。} : « Les élèves de notre école ont présenté des spectacles traditionnels. » \textbf{Justification} : le \zh{了} après le verbe \zh{表演} marque l'action accomplie, d'où le passé composé en français.
\item \zh{中国和喀麦隆都有丰富的文化。} : « La Chine et le Cameroun ont tous les deux une culture riche. » \textbf{Justification} : \zh{都} « tous les deux » se place après les deux sujets reliés par \zh{和}.
\end{enumerate}
\end{corrige}

\begin{corrige}[Exercice 6 — Appariement de collocations]
\begin{itemize}
\item \zh{过节日} (guò jiérì) : « passer, fêter une fête ».
\item \zh{穿传统服装} (chuān chuántǒng fúzhuāng) : « porter des vêtements traditionnels ».
\item \zh{表演节目} (biǎoyǎn jiémù) : « présenter un spectacle ».
\item \zh{唱歌} (chàng gē) : « chanter (une chanson) ».
\end{itemize}
\textbf{À retenir} : en chinois, chaque verbe s'associe à des compléments précis ; on dit \zh{过节} « fêter », jamais \zh{做节}. De même, on dit \zh{穿} pour les vêtements et \zh{戴} (dài) pour les accessoires (chapeau, lunettes).
\end{corrige}

\begin{corrige}[Exercice 7 — Texte à trous : caractères et mots-outils]
\begin{enumerate}
\item \zh{春节是中国人最重要的节日。} — \zh{最} + adjectif forme le superlatif (« le plus ») ; \zh{的} relie le complément \zh{重要} au nom \zh{节日}.
\item \zh{人们穿新衣服，互相祝福。} — \zh{互相} « mutuellement » se place avant le verbe \zh{祝福}.
\item \zh{还看了很多节目。} — \zh{了} après le verbe \zh{看} marque l'action accomplie.
\item \zh{唱歌和跳舞都是春节的传统活动。} — \zh{和} « et » relie deux sujets ; \zh{传统} qualifie \zh{活动}.
\end{enumerate}
Texte complet traduit : « La fête du Printemps est la fête la plus importante des Chinois. Au moment de la fête, les gens portent des vêtements neufs et échangent leurs vœux. Toute notre famille mange ensemble le repas du réveillon et regarde aussi beaucoup de spectacles. Chanter et danser sont des activités traditionnelles de la fête du Printemps. »

\textbf{Point de vigilance} : distinguer les trois emplois des mots-outils : \zh{的} (lien complément-nom), \zh{了} (accompli, après le verbe), \zh{和} (coordination de noms ou de sujets, jamais de phrases entières).
\end{corrige}

\begin{corrige}[Exercice 8 — Remets les mots dans l'ordre]
\begin{enumerate}
\item \zh{喀麦隆人非常喜欢音乐和舞蹈。} (Kāmàilóng rén fēicháng xǐhuan yīnyuè hé wǔdǎo.) « Les Camerounais aiment beaucoup la musique et la danse. »
\item \zh{春节是中国最重要的传统节日。} (Chūnjié shì Zhōngguó zuì zhòngyào de chuántǒng jiérì.) « La fête du Printemps est la fête traditionnelle la plus importante de Chine. »
\item \zh{过节日的时候，人们穿漂亮的传统服装。} (Guò jiérì de shíhou, rénmen chuān piàoliang de chuántǒng fúzhuāng.) « Pendant les fêtes, les gens portent de beaux vêtements traditionnels. » \textbf{Justification} : le complément de temps \zh{过节日的时候} se place en tête de phrase.
\item \zh{学生们表演了精彩的节目。} (Xuéshengmen biǎoyǎn le jīngcǎi de jiémù.) « Les élèves ont présenté un spectacle magnifique. »
\end{enumerate}
\textbf{Point de vigilance} : l'ordre chinois est Sujet + (complément de temps) + verbe + complément ; l'adjectif qualificatif (\zh{漂亮的}, \zh{精彩的}) se place avant le nom, suivi de \zh{的}.
\end{corrige}

\begin{corrige}[Exercice 9 — Compréhension écrite type Bac : la fête du Printemps]
\textbf{Barème indicatif (10 points)} : 1,5 point par question (exactitude du sens + correction de la phrase en chinois), soit 7,5 points ; question de langue : 2,5 points (1 point pour la phrase relevée, 0,5 pour la traduction, 1 pour la phrase personnelle correcte).

\begin{enumerate}
\item \zh{春节是中国人最重要的传统节日。} (Chūnjié shì Zhōngguó rén zuì zhòngyào de chuántǒng jiérì.) « C'est la fête traditionnelle la plus importante des Chinois. »
\item \zh{过春节以前，人们打扫房子，买很多好吃的东西。} (Guò Chūnjié yǐqián, rénmen dǎsǎo fángzi, mǎi hěn duō hǎochī de dōngxi.) « Avant la fête, les gens nettoient la maison et achètent beaucoup de bonnes choses à manger. »
\item \zh{除夕晚上，全家一起吃年夜饭，互相祝福。} (Chúxī wǎnshang, quánjiā yíqǐ chī niányèfàn, hùxiāng zhùfú.) « Le soir du réveillon, toute la famille mange ensemble le repas du Nouvel An et échange ses vœux. »
\item \zh{因为他们可以穿新衣服，还可以拿到红包。} (Yīnwèi tāmen kěyǐ chuān xīn yīfu, hái kěyǐ ná dào hóngbāo.) « Parce qu'ils peuvent porter des vêtements neufs et recevoir des enveloppes rouges. »
\item \zh{人们可以看舞龙和舞狮表演，这些表演非常热闹。} (Rénmen kěyǐ kàn wǔlóng hé wǔshī biǎoyǎn, zhèxiē biǎoyǎn fēicháng rènao.) « On peut voir des danses du dragon et du lion ; ces spectacles sont très animés. »
\end{enumerate}

\textbf{Question de langue} : phrase relevée : \zh{孩子们最喜欢春节，因为他们可以穿新衣服，还可以拿到红包。} (Háizimen zuì xǐhuan Chūnjié, yīnwèi tāmen kěyǐ chuān xīn yīfu, hái kěyǐ ná dào hóngbāo.) « Les enfants aiment le plus la fête du Printemps parce qu'ils peuvent porter des vêtements neufs et recevoir des enveloppes rouges. » Phrase personnelle (exemple) : \zh{我喜欢喀麦隆，因为喀麦隆的文化很丰富。} (Wǒ xǐhuan Kāmàilóng, yīnwèi Kāmàilóng de wénhuà hěn fēngfù.) « J'aime le Cameroun parce que sa culture est très riche. »

\textbf{Points de vigilance} : répondre par des phrases complètes (sujet + verbe), recopier les caractères sans erreur, et placer \zh{因为} en tête de la subordonnée de cause.
\end{corrige}

\begin{corrige}[Exercice 10 — Thème et version]
\textbf{Barème indicatif (10 points)} : version 5 points (1 point par segment correct, caractères et ordre des mots) ; thème 5 points (fidélité du sens et qualité du français).

\textbf{Version (chinois → français)} \\
\textit{Texte source :} \zh{喀麦隆人非常喜欢音乐和舞蹈。过节日的时候，人们穿漂亮的传统服装。家人团聚，一起吃饭，互相祝福。} \\
\textit{Traduction corrigée :} « Les Camerounais aiment beaucoup la musique et la danse. Pendant les fêtes, les gens portent de beaux vêtements traditionnels. La famille se retrouve, mange ensemble et s'échange des vœux. »

\textbf{Points de méthode} : « pendant les fêtes » se traduit par le complément de temps \zh{过节日的时候} placé avant le sujet ou juste après ; « mutuellement » = \zh{互相}, placé avant le verbe ; « se souhaiter le bonheur » se rend naturellement par \zh{互相祝福} ; « se retrouver » = \zh{团聚} (tuánjù).

\textbf{Thème (français → chinois)} \\
\textit{Texte source :} « La culture chinoise est très riche. Pendant les fêtes, les gens chantent, dansent et présentent des spectacles. Tout le monde est très joyeux. » \\
\textit{Traduction corrigée :} \zh{中国文化非常丰富。过节日的时候，人们唱歌、跳舞、表演节目。大家都非常高兴。} (Zhōngguó wénhuà fēicháng fēngfù. Guò jiérì de shíhou, rénmen chàng gē, tiào wǔ, biǎoyǎn jiémù. Dàjiā dōu fēicháng gāoxìng.)

\textbf{Points de vigilance} : ne pas traduire \zh{大家} par « grande famille » mais par « tout le monde » ; la virgule énumérative chinoise \zh{、} correspond à des virgules ou à « et » en français ; \zh{过节日的时候} structure le temps en tête de phrase.
\end{corrige}

\begin{corrige}[Exercice 11 — Expression écrite type Bac — modèle de production]
\textbf{Barème indicatif (10 points)} : respect de la longueur (60 à 80 caractères) : 2 points ; au moins cinq mots du thème : 2 points ; structure en \zh{的时候} : 1,5 point ; structure en \zh{互相} : 1,5 point ; correction des caractères et de l'ordre des mots : 3 points.

\textbf{Modèle de production} :\\
\zh{在我的家乡，有一个很重要的传统节日。过这个节日的时候，人们穿漂亮的传统服装，唱歌、跳舞、表演节目。家人一起吃饭，互相祝福。大家都非常高兴。我很喜欢我们喀麦隆的文化。}\\
(Zài wǒ de jiāxiāng, yǒu yí ge hěn zhòngyào de chuántǒng jiérì. Guò zhège jiérì de shíhou, rénmen chuān piàoliang de chuántǒng fúzhuāng, chàng gē, tiào wǔ, biǎoyǎn jiémù. Jiārén yíqǐ chīfàn, hùxiāng zhùfú. Dàjiā dōu fēicháng gāoxìng. Wǒ hěn xǐhuan wǒmen Kāmàilóng de wénhuà.)

Traduction : « Dans mon village, il y a une fête traditionnelle très importante. Au moment de cette fête, les gens portent de beaux vêtements traditionnels, chantent, dansent et présentent des spectacles. Les familles mangent ensemble et échangent leurs vœux. Tout le monde est très joyeux. J'aime beaucoup notre culture camerounaise. »

\textbf{Critères de réussite} : texte de 60 à 80 caractères ; au moins cinq mots du thème ; une structure en \zh{的时候} ; une structure en \zh{互相} ; ordre des mots correct (sujet + verbe + complément).

\textbf{Points de vigilance} : ne pas écrire \zh{互相} après le verbe ; ne pas oublier \zh{的} entre l'adjectif et le nom (\zh{漂亮的传统服装}) ; compter les caractères (la ponctuation ne compte pas) ; éviter de recopier mot à mot le texte support sans adaptation personnelle.
\end{corrige}

\newpage

\coursec{Fiche bilan}

\begin{popfigure}
\begin{tikzpicture}[mindmap, grow cyclic, every node/.style={concept, font=\small}, concept color=popblue,
  level 1/.append style={level distance=3.1cm, sibling angle=51.4, font=\footnotesize},
  level 2/.append style={level distance=1.9cm, font=\scriptsize}]
\node[concept color=popblue, text=white, align=center]{Leçon 19\\ Vie culturelle\\ Cameroun -- Chine}
  child[concept color=poporange]{ node[align=center]{Lexique\\ culture}
    child{ node[align=center]{文化 wénhuà\\ 艺术 yìshù} }
    child{ node[align=center]{节日 jiérì\\ 音乐 yīnyuè} }
  }
  child[concept color=popgreen]{ node[align=center]{Lire un\\ texte}
    child{ node[align=center]{repérer\\ thème + mots-clés} }
  }
  child[concept color=poppurple]{ node[align=center]{Grammaire}
    child{ node[align=center]{会 huì :\\ savoir-faire} }
    child{ node[align=center]{喜欢 xǐhuan\\ + verbe} }
  }
  child[concept color=poppink]{ node[align=center]{Clés et\\ traits}
    child{ node[align=center]{艹 herbe\\ 讠parole} }
  }
  child[concept color=popgold]{ node[align=center]{Comprendre\\ répondre}
    child{ node[align=center]{questions\\ 谁 什么 为什么} }
  }
  child[concept color=popgreen!70!popblue]{ node[align=center]{Cultures\\ comparées}
    child{ node[align=center]{ndolé, makossa\\ opéra, calligraphie} }
  };
\end{tikzpicture}
\legende{Carte mentale de la leçon 19 : vie culturelle au Cameroun et en Chine}
\end{popfigure}

\begin{bilan}
\textbf{1. Lexique clé à connaître par cœur}
\begin{itemize}\setlength{\itemsep}{1pt}
  \item \zh{文化} wénhuà : la culture ; \zh{艺术} yìshù : l'art ; \zh{节日} jiérì : la fête ; \zh{音乐} yīnyuè : la musique.
  \item \zh{电影} diànyǐng : le film ; \zh{博物馆} bówùguǎn : le musée ; \zh{舞蹈} wǔdǎo : la danse ; \zh{书法} shūfǎ : la calligraphie.
  \item \zh{传统} chuántǒng : la tradition ; \zh{表演} biǎoyǎn : le spectacle, représenter ; \zh{参观} cānguān : visiter ; \zh{有意思} yǒu yìsi : intéressant.
  \item \zh{春节} Chūnjié : la fête du Printemps (Nouvel An chinois) ; \zh{喀麦隆} Kāmàilóng : le Cameroun.
\end{itemize}

\textbf{2. Règles de grammaire à maîtriser}
\begin{center}
\begin{tabularx}{\linewidth}{|X|X|X|}
\hline
\rowcolor{popblueL} Structure & Exemple & Traduction \\
\hline
Sujet + 会 + verbe (savoir-faire acquis) & \zh{他会跳舞。} Tā huì tiàowǔ. & Il sait danser. \\
\hline
Sujet + 喜欢 + verbe / nom & \zh{我喜欢听音乐。} Wǒ xǐhuan tīng yīnyuè. & J'aime écouter de la musique. \\
\hline
Sujet + 跟 / 和 + quelqu'un + verbe & \zh{我跟朋友参观博物馆。} Wǒ gēn péngyou cānguān bówùguǎn. & Je visite le musée avec des amis. \\
\hline
Phrase d'existence : lieu + 有 + nom & \zh{雅温得有很多文化活动。} Yǎwēndé yǒu hěn duō wénhuà huódòng. & Il y a beaucoup d'activités culturelles à Yaoundé. \\
\hline
\end{tabularx}
\end{center}

\textbf{3. Clés (radicaux) et ordre des traits}
\begin{itemize}\setlength{\itemsep}{1pt}
  \item \zh{艺} (yì, art) : clé de l'herbe \zh{艹} en haut ; on écrit d'abord la clé, puis le bas (haut $\rightarrow$ bas).
  \item \zh{话} (huà, parole) : clé de la parole \zh{讠} à gauche, puis \zh{舌} à droite (gauche $\rightarrow$ droite).
  \item \zh{节} (jié, fête) : clé \zh{艹} en haut, puis la partie inférieure ; traits horizontaux avant verticaux.
  \item \zh{博} (bó, vaste) : clé \zh{十} à gauche ; de gauche à droite, de haut en bas.
\end{itemize}

\textbf{4. Méthodes express}
\begin{itemize}\setlength{\itemsep}{1pt}
  \item \textbf{Lire un texte} : 1) lire le titre et la première phrase ; 2) souligner les mots du thème (\zh{文化}, \zh{节日}, \zh{音乐}) ; 3) repérer les mots interrogatifs des questions ; 4) chercher la réponse phrase par phrase.
  \item \textbf{Répondre} : reprendre la structure de la question dans la réponse, sans oublier le sujet ni la ponctuation chinoise 。
  \item \textbf{Mémoriser} : apprendre les mots par familles (fêtes / arts / lieux) et avec leur clé.
\end{itemize}
\end{bilan}

\begin{aretenir}[Checklist avant le Bac]
\begin{itemize}\setlength{\itemsep}{1pt}
  \item[$\square$] Je connais le lexique de la vie culturelle : \zh{文化}, \zh{艺术}, \zh{节日}, \zh{音乐}, \zh{电影}, \zh{博物馆}, \zh{舞蹈}, \zh{书法}.
  \item[$\square$] Je prononce correctement les tons du vocabulaire du thème (pinyin sans faute).
  \item[$\square$] Je sais employer \zh{会} + verbe pour exprimer un savoir-faire.
  \item[$\square$] Je sais employer \zh{喜欢} + verbe ou nom pour exprimer un goût culturel.
  \item[$\square$] Je construis une phrase d'existence avec \zh{有} : lieu + \zh{有} + nom.
  \item[$\square$] Je reconnais les clés \zh{艹} (herbe) et \zh{讠} (parole) et je respecte l'ordre des traits.
  \item[$\square$] Je sais lire un court texte et repérer rapidement les mots-clés du thème.
  \item[$\square$] Je réponds aux questions de compréhension en phrase complète, avec la ponctuation chinoise.
  \item[$\square$] Je distingue \zh{参观} (visiter un lieu) et \zh{看} (regarder un spectacle, un film).
  \item[$\square$] Je peux comparer en quelques phrases la vie culturelle du Cameroun (makossa, ndolé) et de la Chine (calligraphie, fête du Printemps).
  \item[$\square$] Je n'oublie pas le classificateur : \zh{一场电影} (un film/séance), \zh{一个节日} (une fête).
  \item[$\square$] Je relis ma production : ordre Sujet -- Verbe -- Complément, sans mot français laissé dans la phrase.
\end{itemize}
\end{aretenir}

\findecours

\end{document}
